\appendix

\section{Proof Details}\label{sec:proof-details}

\subsection{Static Certificate and Projective Closure}\label{app:step-001}

\begin{lemma}[Anchor nonvanishing] \label{lem:step-001-anchor}

Under Assumptions~\ref{assump:parameter-regime}, \ref{assump:balcan-common-chain}, and
\ref{assump:anchored-derivative-closure}, for every \(\theta\in\Theta\),

\begin{equation}
\label{eq:appendix-display-001}
F_{j_*}(\theta)=1,
\qquad
\|F(\theta)\|_2\geq1,
\qquad
F(\theta)\neq0.
\end{equation}

Consequently \(\gamma_F=F/\|F\|_2\) is well-defined and differentiable along \(\Theta\).


\end{lemma}

\begin{proof}
By the anchored closure assumption, \(Q_{j_*}\equiv1\). Therefore, using the
setting definition of an output feature,

\begin{equation}
\label{eq:appendix-display-002}
F_{j_*}(\theta)
=Q_{j_*}(\theta,\eta_1(\theta),\ldots,\eta_q(\theta))
=1
\end{equation}

for every \(\theta\in U\), and in particular for every \(\theta\in\Theta\). Thus

\begin{equation}
\label{eq:appendix-display-003}
\|F(\theta)\|_2^2
=\sum_{i=1}^N F_i(\theta)^2
\geq F_{j_*}(\theta)^2
=1.
\end{equation}

The common-chain presentation makes the outputs differentiable on \(U\);
equivalently, Assumption~\ref{assump:anchored-derivative-closure} supplies
their derivatives there. Since the denominator is bounded below by one on a
neighborhood of every point of \(\Theta\), normalization is differentiable, including at the endpoints of
\(\Theta\). This proves the lemma. 
\end{proof}

\begin{lemma}[Literal coefficient-height matrix certificate]
\label{lem:step-001-height}

Under Assumptions~\ref{assump:parameter-regime} and
\ref{assump:anchored-derivative-closure}, for every \(\theta\in\Theta\),

\begin{equation}
\label{eq:appendix-display-004}
\|B(\theta)\|_{\mathrm{op}}
\leq \|B(\theta)\|_{\mathrm F}
\leq \widehat\Lambda_{B,T}.
\end{equation}

Consequently

\begin{equation}
\label{eq:appendix-display-005}
\sup_{\theta\in\Theta}\|B(\theta)\|_{\mathrm{op}}
\leq\widehat\Lambda_{B,T}.
\end{equation}


\end{lemma}

\begin{proof}
Retain the setting-defined \(T_*:=\max\{1,T\}\), and introduce only for this proof

\begin{equation}
\label{eq:appendix-display-006}
c_{rs}:=\sum_{\ell=0}^{m}|b_{rs,\ell}|T_*^\ell,
\qquad 0\leq r,s\leq N.
\end{equation}

For \(\theta\in\Theta\), \(|\theta|\leq T\leq T_*\). Because \(\ell\) is a nonnegative integer and
\(T_*\geq1\), \(|\theta|^\ell\leq T_*^\ell\) also for \(\ell=0\). Hence, entry by entry,

\begin{equation}
\label{eq:appendix-display-007}
\begin{aligned}
|B_{rs}(\theta)|
&=\left|\sum_{\ell=0}^{m}b_{rs,\ell}\theta^\ell\right|\\
&\leq\sum_{\ell=0}^{m}|b_{rs,\ell}|\,|\theta|^\ell\\
&\leq\sum_{\ell=0}^{m}|b_{rs,\ell}|T_*^\ell
=c_{rs}.
\end{aligned}
\end{equation}

Summing the squared entrywise bounds over exactly the \((N+1)^2\) entries gives

\begin{equation}
\label{eq:appendix-display-008}
\|B(\theta)\|_{\mathrm F}^2
=\sum_{r=0}^{N}\sum_{s=0}^{N}|B_{rs}(\theta)|^2
\leq\sum_{r=0}^{N}\sum_{s=0}^{N}c_{rs}^2
=\sum_{r=0}^{N}\sum_{s=0}^{N}
\left(\sum_{\ell=0}^{m}|b_{rs,\ell}|T_*^\ell\right)^2
=\widehat\Lambda_{B,T}^2.
\end{equation}

For completeness, the operator/Frobenius comparison is also direct. For any
\(x=(x_0,\ldots,x_N)\in\mathbb R^{N+1}\), applying the scalar Cauchy--Schwarz inequality separately in each
row yields

\begin{equation}
\label{eq:appendix-display-009}
\begin{aligned}
\|B(\theta)x\|_2^2
&=\sum_{r=0}^{N}\left|\sum_{s=0}^{N}B_{rs}(\theta)x_s\right|^2\\
&\leq\sum_{r=0}^{N}
\left(\sum_{s=0}^{N}|B_{rs}(\theta)|^2\right)
\left(\sum_{s=0}^{N}|x_s|^2\right)\\
&=\|B(\theta)\|_{\mathrm F}^2\|x\|_2^2.
\end{aligned}
\end{equation}

Taking the supremum over unit vectors proves
\(\|B(\theta)\|_{\mathrm{op}}\leq\|B(\theta)\|_{\mathrm F}\). Combining the two pointwise inequalities and
then taking the supremum over \(\Theta\) proves the lemma. No coefficient, power, matrix entry, or dimensional
factor has been absorbed. 
\end{proof}

\begin{lemma}[Homogeneous block extraction]
\label{lem:step-001-homogeneous-block}

Under Assumption~\ref{assump:anchored-derivative-closure}, if \(F_0\equiv0\) and \(B_F(\theta)\) denotes the
\(N\times N\) block \((B_{rs}(\theta))_{1\leq r,s\leq N}\), then

\begin{equation}
\label{eq:appendix-display-010}
F'(\theta)=B_F(\theta)F(\theta)
\qquad\text{for every }\theta\in\Theta.
\end{equation}


\end{lemma}

\begin{proof}
In the homogeneous specialization,

\begin{equation}
\label{eq:appendix-display-011}
\widetilde F(\theta)=
\begin{pmatrix}0\\F(\theta)\end{pmatrix},
\qquad
\widetilde F'(\theta)=
\begin{pmatrix}0\\F'(\theta)\end{pmatrix}.
\end{equation}

The exact closure identity therefore reads

\begin{equation}
\label{eq:appendix-display-012}
\begin{pmatrix}0\\F'(\theta)\end{pmatrix}
=B(\theta)
\begin{pmatrix}0\\F(\theta)\end{pmatrix}.
\end{equation}

Taking rows \(1,\ldots,N\), the column-zero terms are multiplied by \(F_0(\theta)=0\), while the remaining
columns are exactly \(B_F(\theta)F(\theta)\). Thus \(F'=B_FF\). No vanishing of any off-block entry of \(B\)
has been assumed. 
\end{proof}

\begin{proposition}[Normalized derivative and projective certificate]
\label{prop:step-001-projective}

Under Assumptions~\ref{assump:parameter-regime}, \ref{assump:balcan-common-chain}, and
\ref{assump:anchored-derivative-closure}, and Lemmas~\ref{lem:step-001-anchor},
\ref{lem:step-001-height}, and \ref{lem:step-001-homogeneous-block}, if \(F_0\equiv0\), then for every
\(\theta\in\Theta\),

\begin{equation}
\label{eq:appendix-display-013}
\gamma_F'(\theta)
=\bigl(I_N-\gamma_F(\theta)\gamma_F(\theta)^{\mathsf T}\bigr)
B_F(\theta)\gamma_F(\theta),
\end{equation}

and

\begin{equation}
\label{eq:appendix-display-014}
\Gamma_{\mathrm{proj}}(F)
=\sup_{\theta\in\Theta}\|\gamma_F'(\theta)\|_2
\leq\widehat\Lambda_{B,T}.
\end{equation}


\end{proposition}

\begin{proof}
Fix \(\theta\in\Theta\), and write
\(r(\theta):=\|F(\theta)\|_2\) within this proof. Lemma~\ref{lem:step-001-anchor} gives \(r(\theta)\geq1\), so
all following divisions are valid. Differentiating \(r^2=\langle F,F\rangle\) gives

\begin{equation}
\label{eq:appendix-display-015}
2rr'=2\langle F,F'\rangle,
\qquad
r'=\frac{\langle F,F'\rangle}{r}.
\end{equation}

Therefore direct differentiation of \(\gamma_F=F/r\) gives

\begin{equation}
\label{eq:appendix-display-016}
\begin{aligned}
\gamma_F'
&=\frac{F'}{r}-\frac{F r'}{r^2}\\
&=\frac{F'}{r}-\frac{F\langle F,F'\rangle}{r^3}\\
&=\bigl(I_N-\gamma_F\gamma_F^{\mathsf T}\bigr)\frac{F'}{r}.
\end{aligned}
\end{equation}

By Lemma~\ref{lem:step-001-homogeneous-block},
\(F'=B_FF=rB_F\gamma_F\). Substitution proves the exact identity

\begin{equation}
\label{eq:appendix-display-017}
\gamma_F'
=\bigl(I_N-\gamma_F\gamma_F^{\mathsf T}\bigr)B_F\gamma_F.
\end{equation}

It remains to prove the norm bound without hiding a dimensional constant. Since \(\|\gamma_F\|_2=1\), for every
\(v\in\mathbb R^N\),

\begin{equation}
\label{eq:appendix-display-018}
\left\|\bigl(I_N-\gamma_F\gamma_F^{\mathsf T}\bigr)v\right\|_2^2
=\|v\|_2^2-|\langle\gamma_F,v\rangle|^2
\leq\|v\|_2^2.
\end{equation}

Thus the projector has Euclidean operator norm at most one. Also, if \(x\in\mathbb R^N\), embed it as
\(\bar x=(0,x)\in\mathbb R^{N+1}\), and let \(\pi\) discard coordinate zero. Then

\begin{equation}
\label{eq:appendix-display-019}
B_F(\theta)x=\pi B(\theta)\bar x,
\end{equation}

so

\begin{equation}
\label{eq:appendix-display-020}
\|B_F(\theta)x\|_2
\leq\|B(\theta)\bar x\|_2
\leq\|B(\theta)\|_{\mathrm{op}}\|\bar x\|_2
=\|B(\theta)\|_{\mathrm{op}}\|x\|_2.
\end{equation}

Taking the supremum over unit \(x\) proves
\(\|B_F(\theta)\|_{\mathrm{op}}\leq\|B(\theta)\|_{\mathrm{op}}\). Hence the exact identity and
Lemma~\ref{lem:step-001-height} yield, pointwise,

\begin{equation}
\label{eq:appendix-display-021}
\begin{aligned}
\|\gamma_F'(\theta)\|_2
&\leq
\|I_N-\gamma_F(\theta)\gamma_F(\theta)^{\mathsf T}\|_{\mathrm{op}}
\|B_F(\theta)\|_{\mathrm{op}}
\|\gamma_F(\theta)\|_2\\
&\leq\|B(\theta)\|_{\mathrm{op}}\\
&\leq\widehat\Lambda_{B,T}.
\end{aligned}
\end{equation}

The right-hand side is finite deterministic instance data. Taking the supremum over \(\theta\in\Theta\) both
proves finiteness of \(\Gamma_{\mathrm{proj}}(F)\) and gives the claimed bound. 
\end{proof}

\begin{proposition}[Boundary and baseline consistency]
\label{prop:step-001-boundary}

Under Assumptions~\ref{assump:parameter-regime}, \ref{assump:balcan-common-chain}, and
\ref{assump:anchored-derivative-closure}, Lemmas~\ref{lem:step-001-anchor},
\ref{lem:step-001-height}, and \ref{lem:step-001-homogeneous-block}, and
Proposition~\ref{prop:step-001-projective}, all conclusions of those results remain valid without modification
under the static specializations \(q=0\), \(m=0\), constant \(B\), constant \(\widetilde F\), \(N=1\), or
\(\widehat\Lambda_{B,T}=0\), at both endpoints of \(\Theta\), and for a stationary \(\gamma_F\) in the
homogeneous branch. In particular,

\begin{equation}
\label{eq:appendix-display-022}
F_{j_*}=1,
\qquad
F\neq0,
\qquad
\sup_{\theta\in\Theta}\|B(\theta)\|_{\mathrm{op}}
\leq\widehat\Lambda_{B,T}
\end{equation}

still hold. Whenever \(F_0\equiv0\), the same regimes also satisfy

\begin{equation}
\label{eq:appendix-display-023}
F'=B_FF,
\qquad
\gamma_F'=(I_N-\gamma_F\gamma_F^{\mathsf T})B_F\gamma_F,
\qquad
\Gamma_{\mathrm{proj}}(F)\leq\widehat\Lambda_{B,T}.
\end{equation}

In addition:

1. If \(N=d\geq1\), \(m=0\), and the only nonzero entries of the constant augmented monomial shift are
   \(B_{0d}=d\) and \(B_{k+1,k}=k\) for \(1\leq k\leq d-1\), then
\begin{equation}
\label{eq:appendix-indented-001}
   \widehat\Lambda_{B,T}=\left(\sum_{k=1}^d k^2\right)^{1/2}.
\end{equation}
2. If \(N=2\), \(m=0\), \(0<\delta\leq1\), and the only nonzero entry of the constant shear is
   \(B_{2,1}=1/\delta\), then \(\widehat\Lambda_{B,T}=1/\delta\).


\end{proposition}

\begin{proof}
If \(q=0\), the chain is empty and \(M=0\); none of the calculations in
Lemmas~\ref{lem:step-001-anchor}--\ref{lem:step-001-homogeneous-block} or
Proposition~\ref{prop:step-001-projective} uses a chain coordinate. Thus the proof remains unchanged and the
additional dependence on \(q,M,\Delta\), after \(B\) is fixed, is exactly degree zero.

If \(m=0\), then \(B(\theta)=(b_{rs,0})_{r,s=0}^N\) is constant and the defining certificate becomes exactly

\begin{equation}
\label{eq:appendix-display-024}
\widehat\Lambda_{B,T}
=\left(\sum_{r=0}^{N}\sum_{s=0}^{N}|b_{rs,0}|^2\right)^{1/2}
=\|B\|_{\mathrm F},
\end{equation}

because \(T_*^0=1\). This also covers any constant \(B\) presented with a larger formal value of \(m\) and
zero higher coefficients. When \(N=1\), the anchor forces the sole feature to be \(F_1=1\), so
\(\gamma_F=1\) and \(I_1-\gamma_F\gamma_F^{\mathsf T}=0\); the normalized identity and the zero projective
speed therefore hold exactly. All derivatives used in these conclusions exist on the open interval \(U\supseteq\Theta\), so
the same calculations apply at both endpoints of \(\Theta\).

If \(\widetilde F\) is constant, then \(\widetilde F'=0\); in the homogeneous branch the anchored nonzero
vector \(F\) and its normalization \(\gamma_F\) are constant, so \(\Gamma_{\mathrm{proj}}(F)=0\). More
generally, a stationary homogeneous normalized curve, meaning \(\gamma_F'\equiv0\) on \(\Theta\), has
\(\Gamma_{\mathrm{proj}}(F)=0\leq\widehat\Lambda_{B,T}\) directly from the setting definition. This stationary
specialization does not require \(B=0\), but the literal anchor rules out nonzero radial motion along the
current feature direction. Indeed, write \(F=r\gamma_F\), where \(r=\|F\|_2\geq1\). The anchored coordinate
identity is

\begin{equation}
\label{eq:appendix-display-025}
1=F_{j_*}=r\gamma_{F,j_*},
\qquad
\gamma_{F,j_*}=\frac1r>0.
\end{equation}

Differentiating \(F=r\gamma_F\) and its anchored coordinate, and using \(\gamma_F'=0\), gives

\begin{equation}
\label{eq:appendix-display-026}
0=F_{j_*}'=r'\gamma_{F,j_*}+r\gamma_{F,j_*}'
=r'\gamma_{F,j_*},
\end{equation}

so \(r'=0\). Hence \(F'=r'\gamma_F+r\gamma_F'=0\), and
Lemma~\ref{lem:step-001-homogeneous-block} yields

\begin{equation}
\label{eq:appendix-display-027}
0=F'=B_FF=rB_F\gamma_F,
\qquad
B_F\gamma_F=0.
\end{equation}

Thus the projector identity permits radial components algebraically for a generally normalized curve, while
the anchored stationary branch has zero radial speed. The full matrix \(B\) may still be nonzero: annihilating
the current augmented feature direction does not force the entire matrix to vanish.

If \(\widehat\Lambda_{B,T}=0\), the defining sum of squares first gives
\(\sum_{\ell=0}^m|b_{rs,\ell}|T_*^\ell=0\) for every \((r,s)\). Every term in each of these sums is
nonnegative and \(T_*>0\), so every coefficient \(b_{rs,\ell}\) is zero, and hence \(B\equiv0\). The closure
identity gives \(\widetilde F'=0\) on the interval \(U\), so the tuple is
constant; in the homogeneous branch \(\gamma_F'=0\) and \(\Gamma_{\mathrm{proj}}(F)=0\). Thus the zero right-hand
side is valid rather than a degenerate failure of the bound.

For item 1, take the actual augmented monic tuple

\begin{equation}
\label{eq:appendix-display-028}
\widetilde F(\theta)=(\theta^d,1,\theta,\ldots,\theta^{d-1}),
\qquad d\geq1.
\end{equation}

This is the \(q=0\), \(M=0\), \(N=d\), \(\Delta=d\) specialization. Its derivative is given by the constant
shift with \(B_{0d}=d\) and \(B_{k+1,k}=k\) for \(1\leq k\leq d-1\), because
\((\theta^d)'=dF_d\) and \((F_{k+1})'=kF_k\). Thus \(m=0\), and the nonzero coefficient magnitudes are exactly
\(d,1,2,\ldots,d-1\). The literal formula gives

\begin{equation}
\label{eq:appendix-display-029}
\widehat\Lambda_{B,T}^2
=d^2+\sum_{k=1}^{d-1}k^2
=\sum_{k=1}^d k^2.
\end{equation}

The sum over \(1\leq k\leq d-1\) is empty when \(d=1\), so that boundary case gives height one. This proves
only the certificate specialization allocated to this row; it does not import the later affine-monic sweep.

For item 2, the actual tuple \(\widetilde F=(0,1,\theta/\delta)\) has
\(q=0\), \(M=0\), \(\Delta=1\), \(N=2\), and \(m=0\), and it has derivative closure under the constant
one-entry shear \(B_{2,1}=1/\delta\), with all other entries zero. The same literal formula gives

\begin{equation}
\label{eq:appendix-display-030}
\widehat\Lambda_{B,T}^2=|1/\delta|^2=1/\delta^2,
\end{equation}

and \(\delta>0\) gives \(\widehat\Lambda_{B,T}=1/\delta\). These are algebraic reductions of the certificate
itself: they do not assume the later monomial-instance or counter-example conclusions. 
\end{proof}


\begin{proof}[Synthesis of the static certificate]
Lemma~\ref{lem:step-001-anchor} gives
$F_{j_*}=1$, hence $F\neq0$ and $\|F\|_2\geq1$.  From the literal
coefficient list and $|\theta|\leq T_*$, Lemma~\ref{lem:step-001-height}
gives
\begin{equation}
\label{eq:appendix-display-031}
\sup_{\theta\in\Theta}\|B(\theta)\|_{\mathrm{op}}
\leq\widehat\Lambda_{B,T}.
\end{equation}
If $F_0\equiv0$, Lemma~\ref{lem:step-001-homogeneous-block} gives
$F'=B_FF$, and Proposition~\ref{prop:step-001-projective} therefore gives
\begin{equation}
\label{eq:appendix-display-032}
\gamma_F'
=(I_N-\gamma_F\gamma_F^{\mathsf T})B_F\gamma_F,
\qquad
\Gamma_{\mathrm{proj}}(F)\leq\widehat\Lambda_{B,T}.
\end{equation}
Proposition~\ref{prop:step-001-boundary} shows that constant and stationary
curves, endpoints of $\Theta$, zero coefficient height, and the two baseline
specializations preserve these conclusions with their literal constants.
\end{proof}

\subsection{Persistent-Root Nullity}\label{app:step-002}

\begin{lemma}[Persistent-root affine translate and anchor exclusion]
\label{lem:step-002-affine-locus}

Under Assumption~\ref{assump:anchored-derivative-closure} and the Anchor nonvanishing
Lemma~\ref{lem:step-001-anchor}, for every interval \(I\subseteq\Theta\) with \(\lvert I\rvert>0\), the set
\(Z_\infty(I)\) is empty or a proper affine subspace of \(\mathbb R^N\). More precisely, define the common
kernel

\begin{equation}
\label{eq:appendix-display-033}
K_I
:=\bigcap_{\theta\in I}
\ker\!\left(v\longmapsto\langle v,F(\theta)\rangle\right)
=\left\{v\in\mathbb R^N:
\langle v,F(\theta)\rangle=0
\text{ for every }\theta\in I\right\}.
\end{equation}

If \(a^0\in Z_\infty(I)\), then

\begin{equation}
\label{eq:appendix-display-034}
Z_\infty(I)=a^0+K_I,
\qquad
e_{j_*}\notin K_I,
\qquad
a^0+e_{j_*}\notin Z_\infty(I).
\end{equation}


\end{lemma}

\begin{proof}
For each \(\theta\in I\), let

\begin{equation}
\label{eq:appendix-display-035}
H_\theta
:=\left\{a\in\mathbb R^N:
F_0(\theta)+\langle a,F(\theta)\rangle=0\right\}.
\end{equation}

The pointwise meaning of identically zero gives the exact identity

\begin{equation}
\label{eq:appendix-display-036}
Z_\infty(I)=\bigcap_{\theta\in I}H_\theta.
\end{equation}

This intersection may be indexed by uncountably many points; no finite-rank reduction is being assumed. If it
is empty, the first branch of the target is complete. Suppose it is nonempty and choose
\(a^0\in Z_\infty(I)\). For every \(a\in\mathbb R^N\),

\begin{equation}
\label{eq:appendix-display-037}
\begin{aligned}
a\in Z_\infty(I)
&\Longleftrightarrow
F_0(\theta)+\langle a,F(\theta)\rangle=0
\quad\text{for every }\theta\in I\\
&\Longleftrightarrow
\langle a-a^0,F(\theta)\rangle=0
\quad\text{for every }\theta\in I\\
&\Longleftrightarrow a-a^0\in K_I.
\end{aligned}
\end{equation}

Hence \(Z_\infty(I)=a^0+K_I\). Each set in the defining intersection for \(K_I\) is the kernel of a linear
functional, so their arbitrary intersection contains zero and is closed under addition and scalar
multiplication. Thus \(K_I\) is a linear subspace.

The anchor is expressed in the original random-coefficient coordinates and gives

\begin{equation}
\label{eq:appendix-display-038}
\langle e_{j_*},F(\theta)\rangle
=F_{j_*}(\theta)
=1
\qquad\text{for every }\theta\in I.
\end{equation}

Therefore \(e_{j_*}\notin K_I\). In particular \(K_I\neq\mathbb R^N\), so its translate
\(a^0+K_I\) is a proper affine subspace. The required anchor-direction exclusion can also be checked directly:

\begin{equation}
\label{eq:appendix-display-039}
\begin{aligned}
F_0(\theta)+\langle a^0+e_{j_*},F(\theta)\rangle
&=
F_0(\theta)+\langle a^0,F(\theta)\rangle+F_{j_*}(\theta)\\
&=0+1=1
\qquad\text{for every }\theta\in I.
\end{aligned}
\end{equation}

Thus \(a^0+e_{j_*}\notin Z_\infty(I)\).

The argument does not require the nonanchor features to be distinct, nonconstant, or linearly independent. If
\(F_0\equiv0\) on \(I\), then \(0\in Z_\infty(I)\) and the exact identity becomes
\(Z_\infty(I)=K_I\), still proper because it excludes \(e_{j_*}\). If the features are constant or dependent,
\(K_I\) may have any dimension from \(0\) through \(N-1\), and its translate may be unbounded, but properness
is unchanged. When \(N=1\), \(j_*=1\) and

\begin{equation}
\label{eq:appendix-display-040}
K_I=\{v\in\mathbb R:vF_1(\theta)=0\text{ for every }\theta\in I\}=\{0\},
\end{equation}

so a nonempty locus is exactly one point; otherwise it is empty.

Finally, the proof uses the literal point set \(I\). Whether an interval includes neither, one, or both of its
relative endpoints changes only which \(\theta\)'s occur in this intersection. No passage to
\(\overline I\), analytic identity theorem, or endpoint continuation is used. This proves the lemma.

\end{proof}

\begin{lemma}[Borel and Lebesgue nullity of the persistent-root locus]
\label{lem:step-002-lebesgue-null}

Under the Anchor nonvanishing Lemma~\ref{lem:step-001-anchor} and
Lemma~\ref{lem:step-002-affine-locus}, for every interval \(I\subseteq\Theta\) with
\(\lvert I\rvert>0\), the set \(Z_\infty(I)\) is Borel and

\begin{equation}
\label{eq:appendix-display-041}
\lambda_N(Z_\infty(I))=0.
\end{equation}

This conclusion includes empty loci, unbounded nonempty affine loci, and the \(N=1\) empty-or-singleton case.


\end{lemma}

\begin{proof}
For every fixed \(\theta\in I\), the map

\begin{equation}
\label{eq:appendix-display-042}
a\longmapsto F_0(\theta)+\langle a,F(\theta)\rangle
\end{equation}

is a continuous affine map on \(\mathbb R^N\), so \(H_\theta\) is closed. Therefore the exact target set

\begin{equation}
\label{eq:appendix-display-043}
Z_\infty(I)=\bigcap_{\theta\in I}H_\theta
\end{equation}

is closed, hence Borel. Arbitrary intersections of closed sets are closed, so this remains valid when \(I\) is
uncountable and also when the intersection is empty.

Because \(\lvert I\rvert>0\), the interval is nonempty; choose any \(\bar\theta\in I\). The literal
anchor gives \(F_{j_*}(\bar\theta)=1\), and hence the fixed-time set has the graph representation

\begin{equation}
\label{eq:appendix-display-044}
H_{\bar\theta}
=\left\{a\in\mathbb R^N:
a_{j_*}
=-F_0(\bar\theta)
-\sum_{i\neq j_*}a_iF_i(\bar\theta)\right\}.
\end{equation}

Since a persistent identity in particular holds at \(\bar\theta\),

\begin{equation}
\label{eq:appendix-display-045}
Z_\infty(I)\subseteq H_{\bar\theta}.
\end{equation}

If \(N=1\), the sum is empty and
\(H_{\bar\theta}=\{-F_0(\bar\theta)\}\), a singleton of one-dimensional Lebesgue measure zero. Monotonicity
then gives \(\lambda_1(Z_\infty(I))=0\). This agrees with
Lemma~\ref{lem:step-002-affine-locus}, which says that the nonempty locus itself is a singleton.

Suppose \(N\geq2\). Let
\(\beta=(a_i)_{i\neq j_*}\in\mathbb R^{N-1}\) denote the original nonanchor coordinates. For every fixed
\(\beta\), the section of \(H_{\bar\theta}\) in the original coordinate \(a_{j_*}\) is exactly the singleton

\begin{equation}
\label{eq:appendix-display-046}
\left\{
-F_0(\bar\theta)-\sum_{i\neq j_*}\beta_iF_i(\bar\theta)
\right\}.
\end{equation}

After the measure-preserving permutation that places \(a_{j_*}\) last, Tonelli's theorem applied to the
nonnegative Borel indicator of \(H_{\bar\theta}\) yields

\begin{equation}
\label{eq:appendix-display-047}
\begin{aligned}
\lambda_N(H_{\bar\theta})
&=\int_{\mathbb R^{N-1}}
\lambda_1\!\left(
\{t\in\mathbb R:(\beta,t)\in H_{\bar\theta}\}
\right)\,d\lambda_{N-1}(\beta)\\
&=\int_{\mathbb R^{N-1}}0\,d\lambda_{N-1}(\beta)
=0.
\end{aligned}
\end{equation}

Here \((\beta,t)\) means insertion into the original coordinate slots before the harmless permutation. The
graph may be unbounded, but Tonelli is applied to a nonnegative indicator on the full product space, and the
fiber-measure integrand is identically zero. Therefore

\begin{equation}
\label{eq:appendix-display-048}
0\leq\lambda_N(Z_\infty(I))
\leq\lambda_N(H_{\bar\theta})
=0.
\end{equation}

No regularity in \(\theta\), root multiplicity statement, or analyticity is used in this geometric nullity
calculation. This proves the lemma. 
\end{proof}

\begin{proposition}[Nullity under arbitrary full joint laws]
\label{prop:step-002-law-null}

Under Assumption~\ref{assump:cube-density-laws} and
Lemma~\ref{lem:step-002-lebesgue-null}, for every
\(\mu\in\mathcal D_{N,R,\kappa}\) and every interval \(I\subseteq\Theta\) with
\(\lvert I\rvert>0\),

\begin{equation}
\label{eq:appendix-display-049}
\Pr_{\alpha\sim\mu}
\left[
F_0+\langle\alpha,F\rangle\equiv0\text{ on }I
\right]
=\Pr_{\alpha\sim\mu}[\alpha\in Z_\infty(I)]
=0.
\end{equation}

The probability is computed from one full \(N\)-dimensional joint density. No independence, marginal-density
bound, or conditional-density bound is required.


\end{proposition}

\begin{proof}
Fix an arbitrary
\(\mu\in\mathcal D_{N,R,\kappa}\). Assumption~\ref{assump:cube-density-laws} supplies a nonnegative joint
density \(f_\mu\) on the actual space \(\mathbb R^N\), supported on the closed cube \([-R,R]^N\), and satisfying
\(f_\mu\leq\kappa\) almost everywhere. By Lemma~\ref{lem:step-002-lebesgue-null},
\(Z_\infty(I)\) is Borel and \(\lambda_N(Z_\infty(I))=0\). Consequently,

\begin{equation}
\label{eq:appendix-display-050}
\begin{aligned}
\mu(Z_\infty(I))
&=\int_{\mathbb R^N}
\mathbf 1_{Z_\infty(I)}(a)f_\mu(a)\,d\lambda_N(a)\\
&=\int_{Z_\infty(I)\cap[-R,R]^N}
f_\mu(a)\,d\lambda_N(a)\\
&\leq
\kappa\,\lambda_N\!\left(Z_\infty(I)\cap[-R,R]^N\right)\\
&\leq
\kappa\,\lambda_N(Z_\infty(I))
=0.
\end{aligned}
\end{equation}

This is a single joint-density calculation. It never factors \(f_\mu\), integrates a marginal-density bound,
or conditions on one coordinate, so every correlation allowed by the law class is retained. If the affine
locus is unbounded, only its support-cube intersection is integrated, and that intersection remains Borel and
Lebesgue-null. The closed-cube boundary causes no separate term. Empty loci, nonempty proper loci,
\(F_0\equiv0\), dependent or constant features, and the \(N=1\) singleton branch all enter the same zero
calculation. This proves the proposition. 
\end{proof}


\begin{proof}[Synthesis of persistent-root nullity]
Fix an interval $I\subseteq\Theta$ of positive length and
$\mu\in\mathcal D_{N,R,\kappa}$.  By
Lemma~\ref{lem:step-002-affine-locus}, either $Z_\infty(I)$ is empty or,
for any $a^0\in Z_\infty(I)$,
\begin{equation}
\label{eq:appendix-display-051}
Z_\infty(I)=a^0+K_I,
\qquad e_{j_*}\notin K_I.
\end{equation}
Thus it is a proper affine subspace; in particular, $a^0+e_{j_*}$ is not
in the locus.  Lemma~\ref{lem:step-002-lebesgue-null} proves that this
possibly unbounded locus is Borel and Lebesgue-null, and
Proposition~\ref{prop:step-002-law-null} integrates its indicator against
the single full joint density to obtain
\begin{equation}
\label{eq:appendix-display-052}
\Pr_{\alpha\sim\mu}
\left[F_0+\langle\alpha,F\rangle\equiv0\text{ on }I\right]=0.
\end{equation}
This argument retains arbitrary correlation and does not exclude a nonempty
persistent-root locus by assumption.
\end{proof}

\subsection{Measurable Affine Pivot Sweep}\label{app:step-003}

\begin{theorem}[Equal-dimensional Euclidean area formula]\label{thm:area-formula}
Let $D\subseteq\mathbb R^N$ be Lebesgue measurable and let $G:D\to\mathbb R^N$ be Lipschitz, with a
$C^1$ formula on an open neighborhood of $D$. Then, with the derivative of that ambient formula at almost
every density point of $D$,
\begin{equation}
\label{eq:appendix-display-053}
\int_D|\det DG(x)|\,d\lambda_N(x)
=\int_{\mathbb R^N}N(G,D,y)\,d\lambda_N(y),
\qquad
N(G,D,y):=\#\{x\in D:G(x)=y\}.
\end{equation}
In particular,
\begin{equation}
\label{eq:appendix-display-054}
\lambda_N(G(D))\leq\int_D|\det DG(x)|\,d\lambda_N(x),
\end{equation}
without injectivity, noncriticality, simple-root, or finite-fiber hypotheses \citep{federer1969gmt}.
\end{theorem}

In Proposition~\ref{prop:step-003-finite-area}, $D=D_{j,n}\subseteq\mathbb R\times\mathbb R^{N-1}$ and $G=\Psi_j$; source
and target both have dimension $N$. Lemma~\ref{lem:step-003-measurable-domains} supplies completed
measurability, and Lemma~\ref{lem:step-003-finite-chart} supplies the Lipschitz restriction, ambient $C^1$
formula, and exact determinant. Thus every hypothesis of Theorem~\ref{thm:area-formula} is checked in the
same coefficient coordinates used by the probability law.

\begin{lemma}[Measurable incidence and finite pivot exhaustion]
\label{lem:step-003-measurable-domains}

Under Assumptions~\ref{assump:parameter-regime}, \ref{assump:balcan-common-chain},
\ref{assump:anchored-derivative-closure}, and \ref{assump:cube-density-laws}, and the Anchor
nonvanishing Lemma~\ref{lem:step-001-anchor}, fix an arbitrary
\(\mu\in\mathcal D_{N,R,\kappa}\), an interval \(I\subseteq\Theta\) with
\(\lvert I\rvert>0\) and its literal endpoint convention, and a Lebesgue-measurable legal partition
\begin{equation}
\label{eq:appendix-display-055}
I=\bigsqcup_{j=1}^N E_j,\qquad F_j\neq0\text{ on }E_j.
\end{equation}
Define
\begin{equation}
\label{eq:appendix-display-056}
\mathcal R_I
:=\{a\in[-R,R]^N:\exists\theta\in I,\
F_0(\theta)+\langle a,F(\theta)\rangle=0\},
\end{equation}
\begin{equation}
\label{eq:appendix-display-057}
E_{j,n}:=\{\theta\in E_j:\lvert F_j(\theta)\rvert\geq1/n\},
\end{equation}
and
\begin{equation}
\label{eq:appendix-display-058}
D_{j,n}
:=\{(\theta,\beta):
\theta\in E_{j,n},\ \beta\in[-R,R]^{N-1},\
\lvert T_j(\theta,\beta)\rvert\leq R\}.
\end{equation}
Then \(\mathcal R_I\) is Borel; \(E_{j,n}\) and \(D_{j,n}\) are completed-Lebesgue measurable;
\(D_{j,n}\subseteq D_{j,n+1}\); and
\begin{equation}
\label{eq:appendix-display-059}
E_{j,n}\uparrow E_j.
\end{equation}
Moreover \(T_j\), \(\partial_\theta T_j\), and
\begin{equation}
\label{eq:appendix-display-060}
\mathbf 1\{\lvert T_j\rvert\leq R\}\lvert\partial_\theta T_j\rvert
\end{equation}
are measurable on the relevant pivot domains. For \(N=1\), the beta domain is
\([-R,R]^0=\{()\}\), its \(0\)-dimensional Lebesgue mass is one, and all nonpivot sums are empty.


\end{lemma}

\begin{proof}
The primitive common-chain presentation makes each \(F_i\) continuously
differentiable on \(U\). Hence
\begin{equation}
\label{eq:appendix-display-061}
h(\theta,a):=F_0(\theta)+\langle a,F(\theta)\rangle
\end{equation}
is continuous on \(U\times\mathbb R^N\).

Because \(I\) is a bounded interval, it admits compact sets \(K_r\subseteq I\) with
\(K_r\uparrow I\), for every choice of open, closed, or half-open endpoints. Concretely, trim every excluded
endpoint by \(1/r\) and retain every included endpoint once the trimmed interval is nonempty. For each \(r\),
\begin{equation}
\label{eq:appendix-display-062}
C_r:=\{(\theta,a)\in K_r\times[-R,R]^N:h(\theta,a)=0\}
\end{equation}
is closed in a compact product, hence compact. Its coefficient projection is compact. The exact identity
\begin{equation}
\label{eq:appendix-display-063}
\mathcal R_I=\bigcup_{r=1}^{\infty}\operatorname{proj}_a(C_r)
\end{equation}
therefore makes \(\mathcal R_I\) an \(F_\sigma\) subset of \(\mathbb R^N\). This proves Borel measurability
without deleting an included parameter endpoint or a coefficient-cube face.

For fixed \(j,n\), continuity of \(F_j\) makes
\(\{\theta\in\Theta:\lvert F_j(\theta)\rvert\geq1/n\}\) closed. Its intersection with the
Lebesgue-measurable cell \(E_j\) is \(E_{j,n}\), hence is Lebesgue measurable. The product
\(E_{j,n}\times[-R,R]^{N-1}\) is measurable in the Euclidean product completion.

On the open nonzero-pivot set
\begin{equation}
\label{eq:appendix-display-064}
\Omega_j:=\{(\theta,\beta)\in U\times\mathbb R^{N-1}:F_j(\theta)\neq0\},
\end{equation}
write
\begin{equation}
\label{eq:appendix-display-065}
G_j(\theta,\beta)
:=F_0(\theta)+\sum_{i\neq j}\beta_iF_i(\theta).
\end{equation}
Then
\begin{equation}
\label{eq:appendix-display-066}
T_j=-\frac{G_j}{F_j},
\qquad
\partial_\theta T_j
=-\frac{(\partial_\theta G_j)F_j-G_jF_j'}{F_j^2}.
\end{equation}
Both functions are continuous on \(\Omega_j\). Thus \(D_{j,n}\) is the intersection of the measurable
product domain with the measurable weak-threshold set
\(\{\lvert T_j\rvert\leq R\}\); it is completed measurable. This nonnegative integrand is
measurable there as a product of measurable functions.

The threshold \(1/n\) decreases with \(n\), so \(E_{j,n}\subseteq E_{j,n+1}\) and
\(D_{j,n}\subseteq D_{j,n+1}\). If \(\theta\in E_j\), legality gives
\(\lvert F_j(\theta)\rvert>0\). Choosing any integer
\begin{equation}
\label{eq:appendix-display-067}
n\geq\frac{1}{\lvert F_j(\theta)\rvert}
\end{equation}
puts \(\theta\) in \(E_{j,n}\). Hence \(\bigcup_nE_{j,n}=E_j\), with no uniform lower bound on the pivot.
For \(N=1\), the same proof has a one-point beta space and empty sums. 
\end{proof}

\begin{lemma}[Original-coordinate finite charts and exact Jacobian]
\label{lem:step-003-finite-chart}

Under Assumptions~\ref{assump:parameter-regime}, \ref{assump:balcan-common-chain}, and
\ref{assump:anchored-derivative-closure}, and Lemma~\ref{lem:step-003-measurable-domains}, for every
\(j\in\{1,\ldots,N\}\) and \(n\geq1\), define
\begin{equation}
\label{eq:appendix-display-068}
K_{j,n}:=\{\theta\in\Theta:\lvert F_j(\theta)\rvert\geq1/n\}.
\end{equation}
The restrictions of \(T_j\) and \(\Psi_j\) to
\(K_{j,n}\times[-R,R]^{N-1}\), and hence to \(D_{j,n}\), are Lipschitz even when the restriction is
disconnected. The map inserts the solved root in the original \(j\)-th coefficient coordinate:
\begin{equation}
\label{eq:appendix-display-069}
(\Psi_j(\theta,\beta))_j=T_j(\theta,\beta),
\qquad
(\Psi_j(\theta,\beta))_i=\beta_i\quad(i\neq j),
\end{equation}
and it satisfies
\begin{equation}
\label{eq:appendix-display-070}
F_0(\theta)+\langle\Psi_j(\theta,\beta),F(\theta)\rangle=0.
\end{equation}
On the ambient open pivot set, and therefore almost everywhere on every \(D_{j,n}\),
\begin{equation}
\label{eq:appendix-display-071}
\lvert\det D\Psi_j(\theta,\beta)\rvert
=\lvert\partial_\theta T_j(\theta,\beta)\rvert.
\end{equation}
For \(N=1\), this is the determinant of the \(1\times1\) derivative matrix.


\end{lemma}

\begin{proof}
Compactness of \(\Theta\) and primitive \(C^1\) regularity give finite
proof-local quantities
\begin{equation}
\label{eq:appendix-display-072}
M_i:=\sup_{\theta\in\Theta}\lvert F_i(\theta)\rvert,
\qquad
L_i:=\sup_{\theta\in\Theta}\lvert F_i'(\theta)\rvert,
\qquad 0\leq i\leq N.
\end{equation}
If \(K_{j,n}\) is empty, every Lipschitz assertion is vacuous. Otherwise take
\(\theta,\vartheta\in K_{j,n}\). For every \(i\in\{0,\ldots,N\}\setminus\{j\}\), direct quotient algebra,
the two endpoint denominator bounds, and the mean-value bound on the full interval \(\Theta\) give
\begin{equation}
\label{eq:appendix-display-073}
\begin{aligned}
\left\lvert
\frac{F_i(\theta)}{F_j(\theta)}
-\frac{F_i(\vartheta)}{F_j(\vartheta)}
\right\rvert
&\leq
n\lvert F_i(\theta)-F_i(\vartheta)\rvert
+n^2M_i\lvert F_j(\theta)-F_j(\vartheta)\rvert\\
&\leq
(nL_i+n^2M_iL_j)\lvert\theta-\vartheta\rvert.
\end{aligned}
\end{equation}
Also
\begin{equation}
\label{eq:appendix-display-074}
\left\lvert\frac{F_i(\theta)}{F_j(\theta)}\right\rvert\leq nM_i.
\end{equation}
No segment is required to stay in \(K_{j,n}\); only its two endpoints use the pivot margin, while the
mean-value estimate is taken on the ambient interval. The calculation therefore remains valid across
different connected components and if \(F_j\) has opposite signs on those components.

For \(\beta,\gamma\in[-R,R]^{N-1}\), split the difference first in \(\theta\) and then in beta. These
bounds and Cauchy--Schwarz yield
\begin{equation}
\label{eq:appendix-display-075}
\lvert T_j(\theta,\beta)-T_j(\vartheta,\gamma)\rvert
\leq
C_{\theta,j,n}\lvert\theta-\vartheta\rvert
+C_{\beta,j,n}\lVert\beta-\gamma\rVert_2,
\end{equation}
where
\begin{equation}
\label{eq:appendix-display-076}
C_{\theta,j,n}
:=nL_0+n^2M_0L_j
+R\sum_{\substack{1\leq i\leq N\\i\neq j}}
(nL_i+n^2M_iL_j),
\end{equation}
\begin{equation}
\label{eq:appendix-display-077}
C_{\beta,j,n}
:=n\left(\sum_{\substack{1\leq i\leq N\\i\neq j}}M_i^2\right)^{1/2}.
\end{equation}
These are finite explicit proof-level constants, not theorem-facing margins. The resulting bound proves that
\(T_j\) is Lipschitz for the Euclidean product metric. Since \(\Psi_j\) inserts \(T_j\) and copies every beta
coordinate, \(\Psi_j\) is Lipschitz on the same set.

The root identity is exact:
\begin{equation}
\label{eq:appendix-display-078}
\begin{aligned}
F_0(\theta)+\langle\Psi_j(\theta,\beta),F(\theta)\rangle
&=F_0(\theta)+T_j(\theta,\beta)F_j(\theta)
+\sum_{i\neq j}\beta_iF_i(\theta)\\
&=0.
\end{aligned}
\end{equation}
Thus \(\Psi_j\) outputs the original coefficient vector, not a transformed or augmented coordinate.

Order the nonpivot indices as \(i_1,\ldots,i_{N-1}\), use input coordinates
\((\theta,\beta_{i_1},\ldots,\beta_{i_{N-1}})\), and temporarily permute output rows into the order
\((j,i_1,\ldots,i_{N-1})\). The derivative matrix is
\begin{equation}
\label{eq:appendix-display-079}
\begin{pmatrix}
\partial_\theta T_j&
\partial_{\beta_{i_1}}T_j&\cdots&
\partial_{\beta_{i_{N-1}}}T_j\\
0&1&\cdots&0\\
\vdots&\vdots&\ddots&\vdots\\
0&0&\cdots&1
\end{pmatrix}.
\end{equation}
Its determinant is \(\partial_\theta T_j\). Undoing the row permutation changes only the sign, so the
absolute determinant is exactly \(\lvert\partial_\theta T_j\rvert\). The formula is \(C^1\) on
\(\Omega_j\); derivative locality therefore supplies this ambient derivative almost everywhere on the
measurable restriction \(D_{j,n}\). When \(N=1\), there is no row permutation or beta coordinate, and the
matrix is \((\partial_\theta T_1)\). 
\end{proof}

\begin{proposition}[Multiplicity-safe finite-level area and joint-density bound]
\label{prop:step-003-finite-area}

Under Assumption~\ref{assump:cube-density-laws},
Lemmas~\ref{lem:step-003-measurable-domains} and \ref{lem:step-003-finite-chart}, and the restated
equal-dimensional Euclidean area formula, define
\begin{equation}
\label{eq:appendix-display-080}
A_{j,n}:=\Psi_j(D_{j,n}),
\qquad
A_n:=\bigcup_{j=1}^N A_{j,n}.
\end{equation}
For every \(\mu\in\mathcal D_{N,R,\kappa}\), every interval
\(I\subseteq\Theta\) with \(\lvert I\rvert>0\), and every measurable legal pivot partition fixed as in
Lemma~\ref{lem:step-003-measurable-domains}, every \(A_{j,n}\) and \(A_n\) is completed-Lebesgue measurable,
\begin{equation}
\label{eq:appendix-display-081}
A_n\subseteq[-R,R]^N,
\qquad
A_n\subseteq A_{n+1},
\end{equation}
and
\begin{equation}
\label{eq:appendix-display-082}
\mu(A_n)
\leq
\kappa\sum_{j=1}^N
\int_{E_{j,n}}\int_{[-R,R]^{N-1}}
\mathbf 1\{\lvert T_j(\theta,\beta)\rvert\leq R\}
\lvert\partial_\theta T_j(\theta,\beta)\rvert
\,d\beta\,d\theta.
\end{equation}
The conclusion retains zero-Jacobian charts, tangent and multiple roots, finite and infinite fibers, included
interval endpoints, every closed coefficient-cube face, and \(N=1\).


\end{proposition}

\begin{proof}
Fix \(j,n\). The domain \(D_{j,n}\) is completed measurable by
Lemma~\ref{lem:step-003-measurable-domains}, while
Lemma~\ref{lem:step-003-finite-chart} gives a Lipschitz restriction of the ambient \(C^1\) chart and its exact
    determinant. Applying Theorem~\ref{thm:area-formula} gives
\begin{equation}
\label{eq:appendix-display-083}
\begin{aligned}
\lambda_N(A_{j,n})
&\leq
\int_{D_{j,n}}\lvert\det D\Psi_j(\theta,\beta)\rvert\,d\theta\,d\beta\\
&=
\int_{E_{j,n}}\int_{[-R,R]^{N-1}}
\mathbf 1\{\lvert T_j(\theta,\beta)\rvert\leq R\}
\lvert\partial_\theta T_j(\theta,\beta)\rvert
\,d\beta\,d\theta.
\end{aligned}
\end{equation}
Tonelli applies because the integrand is nonnegative and measurable, including if the integral is infinite.

For completeness, completed image measurability is not being assumed. Choose a Borel representative of
\(D_{j,n}\) modulo a null set and intersect that representative with
\begin{equation}
\label{eq:appendix-display-084}
K_{j,n}\times[-R,R]^{N-1}\subseteq\Omega_j.
\end{equation}
This does not change it modulo a null set because \(D_{j,n}\) is already contained in that finite-pivot
product. The map is continuous and Lipschitz there, so the image of the Borel part is analytic and Lebesgue
measurable.
The finite-level Lipschitz map sends the null symmetric difference to a null set by the equal-dimensional
area inequality. Hence \(A_{j,n}\), and then the finite union \(A_n\), is measurable in the completed
\(\lambda_N\)-space and in the completion of \(\mu\).

The exact multiplicity identity is
\begin{equation}
\label{eq:appendix-display-085}
\int_{D_{j,n}}\lvert\partial_\theta T_j\rvert
=
\int_{\mathbb R^N}
N(\Psi_j,D_{j,n},a)\,da.
\end{equation}
Every \(a\in A_{j,n}\) has multiplicity at least one. Several roots or several parameter-beta preimages are
all counted. If the fiber is infinite, the multiplicity is \(+\infty\); the pointwise inequality
\begin{equation}
\label{eq:appendix-display-086}
\mathbf 1_{A_{j,n}}(a)
\leq N(\Psi_j,D_{j,n},a)
\end{equation}
still has the same direction. Thus no injectivity, finite-root count, or simple-root condition enters the
volume estimate.

Differentiate the exact chart identity from Lemma~\ref{lem:step-003-finite-chart} at fixed beta:
\begin{equation}
\label{eq:appendix-display-087}
F_0'(\theta)
+\langle\Psi_j(\theta,\beta),F'(\theta)\rangle
+F_j(\theta)\partial_\theta T_j(\theta,\beta)
=0.
\end{equation}
On \(D_{j,n}\), \(F_j\neq0\). Therefore a tangent root, or any multiple root whose fixed-coefficient first
derivative vanishes, is a critical preimage with \(\partial_\theta T_j=0\). For the measurable critical set
\begin{equation}
\label{eq:appendix-display-088}
C_{j,n}
:=\{(\theta,\beta)\in D_{j,n}:\partial_\theta T_j(\theta,\beta)=0\},
\end{equation}
the same area formula gives
\begin{equation}
\label{eq:appendix-display-089}
\lambda_N(\Psi_j(C_{j,n}))
\leq\int_{C_{j,n}}0=0.
\end{equation}
Critical roots remain in the image; they are controlled by the area formula rather than excluded by a
transversality assumption. An identically zero chart Jacobian is the same zero-integral case.

If an endpoint \(\theta_0\) belongs to \(I\), it belongs literally to its partition cell and, when active, to
the corresponding finite-level domain. The slice
\(\{\theta_0\}\times[-R,R]^{N-1}\) has \(N\)-dimensional measure zero, so its Lipschitz image is null by the
same area formula. Thus endpoint roots are included and controlled; an excluded endpoint never belongs to the
event. The beta cube and the constraint \(\lvert T_j\rvert\leq R\) use weak inequalities, so every cube face
and corner is included.

Every output has nonpivot coordinates in \([-R,R]\) and pivot coordinate in \([-R,R]\), hence
\(A_n\subseteq[-R,R]^N\). Domain nesting gives image nesting. Since \(f_\mu\leq\kappa\) almost everywhere
under the one actual \(N\)-dimensional joint law,
\begin{equation}
\label{eq:appendix-display-090}
\begin{aligned}
\mu(A_n)
&=\int_{A_n}f_\mu(a)\,da\\
&\leq\kappa\lambda_N(A_n)\\
&\leq\kappa\sum_{j=1}^N\lambda_N(A_{j,n}).
\end{aligned}
\end{equation}
Substituting the chart area bounds proves the asserted inequality. Nothing factors \(f_\mu\), conditions on
a coordinate, or replaces it by a marginal density. For \(N=1\), the beta integral is integration over the
one-point \(0\)-dimensional cube with mass
\begin{equation}
\label{eq:appendix-display-091}
\lambda_0([-R,R]^0)=1=(2R)^0.
\end{equation}

\end{proof}

\begin{lemma}[Complete root coverage and persistent-root removal]
\label{lem:step-003-root-coverage}

Under the Anchor nonvanishing Lemma~\ref{lem:step-001-anchor}, the Borel and Lebesgue
nullity Lemma~\ref{lem:step-002-lebesgue-null}, the Nullity under arbitrary full joint laws
Proposition~\ref{prop:step-002-law-null}, and
Lemma~\ref{lem:step-003-measurable-domains}, for every admissible law, positive-length interval, and
measurable legal pivot partition,
\begin{equation}
\label{eq:appendix-display-092}
\mathcal R_I
\subseteq
Z_\infty(I)\cup\bigcup_{n=1}^{\infty}A_n,
\qquad
\mu(Z_\infty(I))=0.
\end{equation}
Every nonpersistent root coefficient enters an actual finite chart in the original coefficient cube. The
statement includes endpoint and cube-boundary roots, tangent and multiple roots, nonpersistent coefficients
with finite or infinite root sets, and persistent or identically-zero affine combinations.


\end{lemma}

\begin{proof}
Lemma~\ref{lem:step-002-lebesgue-null} gives Borel measurability of the exact
persistent locus, and the probability-zero statement is exactly
Proposition~\ref{prop:step-002-law-null}, both instantiated with the same \(F_0,F,I,\mu\) and original
coefficient space. These two results are the only ones used to remove persistent or
identically-zero affine combinations; they are not assumed absent and no independent continuation or root
theorem is invoked.

Take \(a\in\mathcal R_I\setminus Z_\infty(I)\). By the definition of \(\mathcal R_I\), choose one root
\(\theta\in I\). The exact partition has a unique \(j\) with \(\theta\in E_j\), and legality gives
\(F_j(\theta)\neq0\). Let \(\beta=a_{-j}\) be the list of the original nonpivot coordinates. Since
\(a\in[-R,R]^N\),
\begin{equation}
\label{eq:appendix-display-093}
\beta\in[-R,R]^{N-1}.
\end{equation}
Solving the actual affine equation in the actual \(j\)-th coefficient gives
\begin{equation}
\label{eq:appendix-display-094}
\begin{aligned}
a_j
&=-\frac{F_0(\theta)}{F_j(\theta)}
-\sum_{i\neq j}a_i\frac{F_i(\theta)}{F_j(\theta)}\\
&=T_j(\theta,\beta).
\end{aligned}
\end{equation}
Thus \(\lvert T_j(\theta,\beta)\rvert=\lvert a_j\rvert\leq R\). Choose an integer
\begin{equation}
\label{eq:appendix-display-095}
n\geq\frac1{\lvert F_j(\theta)\rvert}.
\end{equation}
Then \(\theta\in E_{j,n}\), \((\theta,\beta)\in D_{j,n}\), and the exact original-coordinate insertion gives
\begin{equation}
\label{eq:appendix-display-096}
\Psi_j(\theta,\beta)=a.
\end{equation}
Therefore \(a\in A_{j,n}\subseteq A_n\), proving the inclusion.

This argument selects only one root and never differentiates it. It applies unchanged to tangent or multiple
roots and to a nonpersistent coefficient with infinitely many roots. Weak cube inequalities retain faces and
corners. An included endpoint is a literal member of \(I\) and of its partition cell; an excluded endpoint is
not part of the event. A pivot may be arbitrarily small, because the chosen root needs only one finite
coefficient-dependent level. For \(N=1\), beta is the empty tuple and the same equation reads
\(a_1=T_1(\theta)\). Persistent and identically-zero cases remain present in
\(Z_\infty(I)\) and are removed only by
Proposition~\ref{prop:step-002-law-null}. 
\end{proof}

\begin{proposition}[Exhausted affine pivot-sweep inequality]
\label{prop:step-003-pivot-sweep}

Under Assumptions~\ref{assump:parameter-regime}, \ref{assump:balcan-common-chain},
\ref{assump:anchored-derivative-closure}, and \ref{assump:cube-density-laws}, the Anchor
nonvanishing Lemma~\ref{lem:step-001-anchor}, the Borel and Lebesgue nullity
Lemma~\ref{lem:step-002-lebesgue-null}, the Nullity under arbitrary full joint laws
Proposition~\ref{prop:step-002-law-null},
Proposition~\ref{prop:step-003-finite-area}, and
Lemma~\ref{lem:step-003-root-coverage}, for every
\(\mu\in\mathcal D_{N,R,\kappa}\), every interval \(I\subseteq\Theta\) with
\(\lvert I\rvert>0\), and every Lebesgue-measurable legal partition
\begin{equation}
\label{eq:appendix-display-097}
I=\bigsqcup_{j=1}^N E_j,\qquad F_j\neq0\text{ on }E_j,
\end{equation}
one has, in the extended nonnegative reals,
\begin{equation}
\label{eq:appendix-display-098}
\begin{aligned}
\Pr_{\alpha\sim\mu}
[\exists\theta\in I:\phi_\alpha(\theta)=0]
&\leq
\kappa\sum_{j=1}^N
\int_{E_j}\int_{[-R,R]^{N-1}}
\mathbf 1\{\lvert T_j(\theta,\beta)\rvert\leq R\}
\lvert\partial_\theta T_j(\theta,\beta)\rvert
\,d\beta\,d\theta\\
&\leq
\kappa\sum_{j=1}^N
\int_{E_j}\int_{[-R,R]^{N-1}}
\lvert\partial_\theta T_j(\theta,\beta)\rvert
\,d\beta\,d\theta.
\end{aligned}
\end{equation}
The constants are literal. The conclusion includes \(N=1\), empty cells, zero Jacobians, pivots without a
uniform margin, every endpoint convention, all root multiplicities and fiber cardinalities, and
identically-zero affine combinations.


\end{proposition}

\begin{proof}
Lemma~\ref{lem:step-003-root-coverage}, completed-measure subadditivity, and the
persistent-root nullity give
\begin{equation}
\label{eq:appendix-display-099}
\begin{aligned}
\Pr_{\alpha\sim\mu}
[\exists\theta\in I:\phi_\alpha(\theta)=0]
&=\mu(\mathcal R_I)\\
&\leq
\mu(Z_\infty(I))
+\mu\left(\bigcup_{n=1}^{\infty}A_n\right)\\
&=
\mu\left(\bigcup_{n=1}^{\infty}A_n\right).
\end{aligned}
\end{equation}
The first equality uses only the primitive support statement
\(\mu([-R,R]^N)=1\), since \(\mathcal R_I\) is the full root event restricted to that support cube.
The measurable sets \(A_n\) increase by
Proposition~\ref{prop:step-003-finite-area}. Continuity from below and the finite-level bound yield
\begin{equation}
\label{eq:appendix-display-100}
\begin{aligned}
\mu\left(\bigcup_{n=1}^{\infty}A_n\right)
&=\lim_{n\to\infty}\mu(A_n)\\
&\leq
\kappa\lim_{n\to\infty}
\sum_{j=1}^N
\int_{E_{j,n}}\int_{[-R,R]^{N-1}}
\mathbf 1\{\lvert T_j\rvert\leq R\}
\lvert\partial_\theta T_j\rvert
\,d\beta\,d\theta.
\end{aligned}
\end{equation}

For fixed \(j\), extend the relevant nonnegative chart integrand by zero outside \(E_j\). By
Lemma~\ref{lem:step-003-measurable-domains},
\(\mathbf 1_{E_{j,n}}\uparrow\mathbf 1_{E_j}\). Monotone convergence, without any integrability
assumption, gives
\begin{equation}
\label{eq:appendix-display-101}
\begin{aligned}
&\lim_{n\to\infty}
\int_{E_{j,n}}\int_{[-R,R]^{N-1}}
\mathbf 1\{\lvert T_j\rvert\leq R\}
\lvert\partial_\theta T_j\rvert
\,d\beta\,d\theta\\
&\qquad=
\int_{E_j}\int_{[-R,R]^{N-1}}
\mathbf 1\{\lvert T_j\rvert\leq R\}
\lvert\partial_\theta T_j\rvert
\,d\beta\,d\theta.
\end{aligned}
\end{equation}
There are exactly \(N<\infty\) chart terms in the proposition, so direct finite addition commutes with the limit and
introduces no extra factor. This proves the indicator inequality, even if one or both sides are \(+\infty\).

Finally, pointwise on every legal chart domain,
\begin{equation}
\label{eq:appendix-display-102}
0\leq
\mathbf 1\{\lvert T_j(\theta,\beta)\rvert\leq R\}
\lvert\partial_\theta T_j(\theta,\beta)\rvert
\leq
\lvert\partial_\theta T_j(\theta,\beta)\rvert.
\end{equation}
Tonelli and finite addition give the indicator-dropped inequality with coefficient one. No term other than the
indicator is dropped.

For \(N=1\), the inner measure is
\(\lambda_0([-R,R]^0)=1\), so this is the exact one-dimensional area inequality. In every dimension,
\begin{equation}
\label{eq:appendix-display-103}
\lambda_{N-1}([-R,R]^{N-1})=(2R)^{N-1};
\end{equation}
the chart theorem leaves that beta integral intact, so this literal factor is available whenever a later
application bounds an integrand uniformly, without this step asserting any later specialization. Empty cells
give zero integrals, while divergent mass near a vanishing pivot is allowed by the extended-real statement.

\end{proof}


\begin{proof}[Synthesis of the measurable pivot sweep]
Fix $\mu\in\mathcal D_{N,R,\kappa}$, a positive-length interval
$I\subseteq\Theta$ with its chosen endpoint convention, and a measurable
legal pivot partition.  Lemma~\ref{lem:step-003-measurable-domains} makes
the root event and finite chart domains measurable and gives
$E_{j,n}\uparrow E_j$.  Lemma~\ref{lem:step-003-finite-chart} supplies the
original-coordinate root identity and
\begin{equation}
\label{eq:appendix-display-104}
|\det D\Psi_j|=|\partial_\theta T_j|
\end{equation}
on each finite domain.  Proposition~\ref{prop:step-003-finite-area} applies
the equal-dimensional area formula with full multiplicity and the single
joint-density bound $f_\mu\leq\kappa$.  Lemma~\ref{lem:step-003-root-coverage}
then covers every nonpersistent root and removes persistent coefficients by
the zero-probability result of Proposition~\ref{prop:step-002-law-null}.

Continuity from below for the coefficient images and monotone convergence
for the nonnegative Jacobian masses, as carried out in
Proposition~\ref{prop:step-003-pivot-sweep}, yield
\begin{equation}
\label{eq:appendix-display-105}
\Pr_{\alpha\sim\mu}[\exists\theta\in I:\phi_\alpha(\theta)=0]
\leq
\kappa\sum_{j=1}^N\int_{E_j}\int_{[-R,R]^{N-1}}
\mathbf1\{|T_j|\leq R\}|\partial_\theta T_j|\,d\beta\,d\theta.
\end{equation}
The pointwise inequality
\begin{equation}
\label{eq:appendix-display-106}
\mathbf1\{|T_j|\leq R\}|\partial_\theta T_j|
\leq|\partial_\theta T_j|
\end{equation}
gives the indicator-dropped form with coefficient one.  The same argument
covers tangent and multiple roots, infinite fibers, zero Jacobians, endpoints,
empty cells, and $N=1$, without a pivot margin or an independence assumption.
\end{proof}

\subsection{Coordinate-Free Section Parametrization}\label{app:step-004}

\begin{lemma}[Exact fixed-section parametrization and graph Jacobian]
\label{lem:step-004-s2-section-parametrization}

Under Assumption~\ref{assump:anchored-derivative-closure}, fix
\(\theta\in\Theta\) and \(j\in\{1,\ldots,N\}\) such that \(F_j(\theta)\neq0\). Define the exact
fixed-section beta domain

\begin{equation}
\label{eq:appendix-display-107}
\mathcal B_{j,\theta}
:=\{\beta\in[-R,R]^{N-1}:
\lvert T_j(\theta,\beta)\rvert\leq R\}.
\end{equation}

Then \(\mathcal B_{j,\theta}\) is Borel and

\begin{equation}
\label{eq:appendix-display-108}
\Psi_j(\theta,\cdot):\mathcal B_{j,\theta}
\longrightarrow H_\theta\cap[-R,R]^N
\end{equation}

is a bijection in the original coefficient coordinates. Its Euclidean
\((N-1)\)-dimensional Jacobian is constant in beta and equals

\begin{equation}
\label{eq:appendix-display-109}
J_{N-1}\Psi_j(\theta,\cdot)
=\frac{\lVert F(\theta)\rVert_2}{\lvert F_j(\theta)\rvert}.
\end{equation}

Consequently, for every nonnegative Borel function \(h\) on the section,

\begin{equation}
\label{eq:appendix-display-110}
\int_{H_\theta\cap[-R,R]^N}h(a)\,d\mathcal H^{N-1}(a)
=
\frac{\lVert F(\theta)\rVert_2}{\lvert F_j(\theta)\rvert}
\int_{\mathcal B_{j,\theta}}
h(\Psi_j(\theta,\beta))\,d\beta.
\end{equation}

The assertions remain literal for an empty section and for \(N=1\), where \(\mathcal H^0\) counts the
unique point of a nonempty section.


\end{lemma}

\begin{proof}
At fixed \(\theta\), \(T_j(\theta,\cdot)\) is affine and hence continuous.
Therefore \(\mathcal B_{j,\theta}\) is the intersection of the closed beta cube with the inverse image of the
closed interval \([-R,R]\), so it is closed and Borel.

For every \(\beta\in\mathcal B_{j,\theta}\), the setting-defined chart uses the original coefficient
coordinates:

\begin{equation}
\label{eq:appendix-display-111}
(\Psi_j(\theta,\beta))_j=T_j(\theta,\beta),
\qquad
(\Psi_j(\theta,\beta))_i=\beta_i\quad(i\neq j).
\end{equation}

All nonpivot coordinates lie in \([-R,R]\), and the defining weak inequality for
\(\mathcal B_{j,\theta}\) puts the pivot coordinate in the same closed interval. Direct substitution into
the affine equation gives

\begin{equation}
\label{eq:appendix-display-112}
\begin{aligned}
F_0(\theta)+\langle\Psi_j(\theta,\beta),F(\theta)\rangle
&=F_0(\theta)+F_j(\theta)T_j(\theta,\beta)
+\sum_{i\neq j}\beta_iF_i(\theta)\\
&=0.
\end{aligned}
\end{equation}

Thus the chart image is contained in \(H_\theta\cap[-R,R]^N\).

Conversely, take \(a\in H_\theta\cap[-R,R]^N\) and set \(\beta_i=a_i\) for \(i\neq j\). Since the pivot
is nonzero, the section equation has exactly one solution in the pivot coordinate:

\begin{equation}
\label{eq:appendix-display-113}
a_j
=-\frac{F_0(\theta)}{F_j(\theta)}
-\sum_{i\neq j}a_i\frac{F_i(\theta)}{F_j(\theta)}
=T_j(\theta,\beta).
\end{equation}

The cube membership of \(a\) implies \(\beta\in\mathcal B_{j,\theta}\), and then
\(a=\Psi_j(\theta,\beta)\). This proves surjectivity. Because the map copies all nonpivot coordinates, two
equal images have equal beta tuples, proving injectivity. No coefficient transformation or extra random
coordinate has been introduced; \(F_0\) remains the same deterministic offset.

Assume first \(N\geq2\), and order the nonpivot indices as \(i_1,\ldots,i_{N-1}\). Differentiating only with
respect to beta gives the columns

\begin{equation}
\label{eq:appendix-display-114}
c_i
:=\partial_{\beta_i}\Psi_j(\theta,\beta)
=e_i-\frac{F_i(\theta)}{F_j(\theta)}e_j,
\qquad i\neq j.
\end{equation}

Let \(v\in\mathbb R^{N-1}\) have coordinates \(v_i=F_i(\theta)/F_j(\theta)\) for \(i\neq j\). The Gram
matrix of the derivative columns is

\begin{equation}
\label{eq:appendix-display-115}
G=(\langle c_i,c_k\rangle)_{i,k\neq j}
=I_{N-1}+vv^{\mathsf T}.
\end{equation}

If \(v=0\), \(G=I_{N-1}\). If \(v\neq0\), \(G\) acts as the identity on \(v^\perp\) and sends \(v\) to
\((1+\lVert v\rVert_2^2)v\). Hence in both cases

\begin{equation}
\label{eq:appendix-display-116}
\det G
=1+\lVert v\rVert_2^2
=1+\sum_{i\neq j}\frac{F_i(\theta)^2}{F_j(\theta)^2}
=\frac{\lVert F(\theta)\rVert_2^2}{F_j(\theta)^2}.
\end{equation}

The Euclidean graph Jacobian is the nonnegative square root of the Gram determinant, so

\begin{equation}
\label{eq:appendix-display-117}
J_{N-1}\Psi_j(\theta,\cdot)
=\sqrt{\det G}
=\frac{\lVert F(\theta)\rVert_2}{\lvert F_j(\theta)\rvert}.
\end{equation}

This computation explicitly retains either sign of the pivot. Applying the restated injective affine area
formula on the measurable domain \(\mathcal B_{j,\theta}\) proves the integral identity asserted by the lemma. Because
the parametrization is affine, its graph Jacobian is constant in beta.

If \(N=1\), then \(j=1\), the beta space is \(\mathbb R^0=\{()\}\), and the empty Gram determinant is one.
Moreover

\begin{equation}
\label{eq:appendix-display-118}
\frac{\lVert F(\theta)\rVert_2}{\lvert F_1(\theta)\rvert}=1.
\end{equation}

The hyperplane \(H_\theta\subseteq\mathbb R\) is the singleton
\(\{-F_0(\theta)/F_1(\theta)\}\). It meets the cube exactly when
\(\lvert T_1(\theta,())\rvert\leq R\). In the nonempty case both \(\lambda_0\) on the beta point and
\(\mathcal H^0\) on the unique section point have mass one; in the empty case both integrals are zero. This
also proves the asserted dimension-zero formula. Cube faces and corners are retained throughout because every
cube inequality is weak. 
\end{proof}

\begin{lemma}[Exact normal-density identity and pivot cancellation]
\label{lem:step-004-s2-normal-density}

Under Assumption~\ref{assump:anchored-derivative-closure} and
Lemma~\ref{lem:step-004-s2-section-parametrization}, fix \(\theta\in\Theta\) and a legal pivot
\(F_j(\theta)\neq0\). For every \(\beta\in\mathbb R^{N-1}\), fixed-beta differentiation of the same chart
root identity gives

\begin{equation}
\label{eq:appendix-display-119}
\partial_\theta T_j(\theta,\beta)
=-
\frac{F_0'(\theta)
+\langle\Psi_j(\theta,\beta),F'(\theta)\rangle}
{F_j(\theta)}.
\end{equation}

On \(\mathcal B_{j,\theta}\), the following is an equality of nonnegative measures under the bijection
\(a=\Psi_j(\theta,\beta)\):

\begin{equation}
\label{eq:appendix-display-120}
\lvert\partial_\theta T_j(\theta,\beta)\rvert\,d\beta
=
\frac{\lvert F_0'(\theta)+\langle a,F'(\theta)\rangle\rvert}
{\lVert F(\theta)\rVert_2}
\,d\mathcal H^{N-1}(a).
\end{equation}

Precisely, for every nonnegative Borel function \(h\) on the section,

\begin{equation}
\label{eq:appendix-display-121}
\begin{aligned}
&\int_{\mathcal B_{j,\theta}}
h(\Psi_j(\theta,\beta))
\lvert\partial_\theta T_j(\theta,\beta)\rvert\,d\beta\\
&\qquad=
\int_{H_\theta\cap[-R,R]^N}
h(a)
\frac{\lvert F_0'(\theta)+\langle a,F'(\theta)\rangle\rvert}
{\lVert F(\theta)\rVert_2}
\,d\mathcal H^{N-1}(a).
\end{aligned}
\end{equation}

The identities include either pivot sign, a zero numerator, zero chart velocity, an empty section, cube faces
and corners, and \(N=1\).


\end{lemma}

\begin{proof}
The setting-defined chart identity, already verified directly in
Lemma~\ref{lem:step-004-s2-section-parametrization}, is

\begin{equation}
\label{eq:appendix-display-122}
F_0(\theta)+F_j(\theta)T_j(\theta,\beta)
+\sum_{i\neq j}\beta_iF_i(\theta)=0.
\end{equation}

Hold beta fixed and differentiate with respect to \(\theta\) on the open set where \(F_j\neq0\). The
setting-defined feature functions and the quotient chart are continuously differentiable there. The product
rule gives

\begin{equation}
\label{eq:appendix-display-123}
F_0'(\theta)
+F_j'(\theta)T_j(\theta,\beta)
+F_j(\theta)\partial_\theta T_j(\theta,\beta)
+\sum_{i\neq j}\beta_iF_i'(\theta)=0.
\end{equation}

Because \(\Psi_j\) inserts \(T_j\) in the original \(j\)-th coordinate and beta in every other original
coordinate,

\begin{equation}
\label{eq:appendix-display-124}
F_j'(\theta)T_j(\theta,\beta)
+\sum_{i\neq j}\beta_iF_i'(\theta)
=\langle\Psi_j(\theta,\beta),F'(\theta)\rangle.
\end{equation}

Dividing by the legal nonzero pivot yields exactly

\begin{equation}
\label{eq:appendix-display-125}
\partial_\theta T_j(\theta,\beta)
=-
\frac{F_0'(\theta)+\langle\Psi_j(\theta,\beta),F'(\theta)\rangle}
{F_j(\theta)}.
\end{equation}

This is the derivative of the same algebraic chart root identity, not a root theorem and not an implicit-root
argument. No root is selected as a function of theta, and no simple-root or transversality condition appears.

Taking absolute values retains both pivot signs:

\begin{equation}
\label{eq:appendix-display-126}
\lvert\partial_\theta T_j(\theta,\beta)\rvert
=
\frac{\lvert F_0'(\theta)
+\langle\Psi_j(\theta,\beta),F'(\theta)\rangle\rvert}
{\lvert F_j(\theta)\rvert}.
\end{equation}

Lemma~\ref{lem:step-004-s2-section-parametrization} gives, under the same bijection,

\begin{equation}
\label{eq:appendix-display-127}
d\mathcal H^{N-1}(a)
=\frac{\lVert F(\theta)\rVert_2}{\lvert F_j(\theta)\rvert}\,d\beta,
\qquad
d\beta
=\frac{\lvert F_j(\theta)\rvert}{\lVert F(\theta)\rVert_2}
\,d\mathcal H^{N-1}(a).
\end{equation}

Multiplying these two literal identities cancels the same absolute pivot factor exactly and gives

\begin{equation}
\label{eq:appendix-display-128}
\lvert\partial_\theta T_j\rvert\,d\beta
=
\frac{\lvert F_0'+\langle a,F'\rangle\rvert}{\lVert F\rVert_2}
\,d\mathcal H^{N-1}(a).
\end{equation}

Applying the affine change-of-variables formula from
Lemma~\ref{lem:step-004-s2-section-parametrization} with the nonnegative test function

\begin{equation}
\label{eq:appendix-display-129}
h(a)\frac{\lvert F_0'(\theta)+\langle a,F'(\theta)\rangle\rvert}
{\lVert F(\theta)\rVert_2}
\end{equation}

makes this differential notation a rigorous integral equality. In particular, taking \(h\equiv1\)
gives the exact fixed-section identity used in Lemma~\ref{lem:step-004-s2-measurable-partition}.

If the numerator is zero at a point, the signed derivative identity makes
\(\partial_\theta T_j=0\) at its unique beta preimage, so both densities vanish. If the whole chart velocity
vanishes, both measures are zero. If the section is empty, both measures are supported on the empty set.
Touching cube faces or corners does not change the affine area formula. For \(N=1\), the Hausdorff Jacobian is
one, \(\lVert F\rVert_2=\lvert F_1\rvert\), and \(\mathcal H^0\) evaluates the right-hand density at the unique
point of a nonempty section; hence the same formulas remain literal. 
\end{proof}

\begin{lemma}[Measurable section mass and exact partition summation]
\label{lem:step-004-s2-measurable-partition}

Under Assumption~\ref{assump:anchored-derivative-closure} and
Lemma~\ref{lem:step-004-s2-normal-density}, fix an interval \(I\subseteq\Theta\) with
\(\lvert I\rvert>0\), retaining its literal endpoint convention, and a Lebesgue-measurable legal partition

\begin{equation}
\label{eq:appendix-display-130}
I=\bigsqcup_{j=1}^N E_j.
\end{equation}

For \(\theta\in E_j\), define the chart mass

\begin{equation}
\label{eq:appendix-display-131}
\mathcal W_j(\theta)
:=\int_{[-R,R]^{N-1}}
\mathbf1\{\lvert T_j(\theta,\beta)\rvert\leq R\}
\lvert\partial_\theta T_j(\theta,\beta)\rvert\,d\beta,
\end{equation}

and define the coordinate-free section mass

\begin{equation}
\label{eq:appendix-display-132}
\mathcal V(\theta)
:=\int_{H_\theta\cap[-R,R]^N}
\frac{\lvert F_0'(\theta)+\langle a,F'(\theta)\rangle\rvert}
{\lVert F(\theta)\rVert_2}
\,d\mathcal H^{N-1}(a).
\end{equation}

Then the cellwise functions are Lebesgue measurable, \(\mathcal V\) is Lebesgue measurable on \(I\), and

\begin{equation}
\label{eq:appendix-display-133}
\mathcal W_j(\theta)=\mathcal V(\theta)
\qquad(\theta\in E_j).
\end{equation}

Consequently, in \([0,\infty]\),

\begin{equation}
\label{eq:appendix-display-134}
\sum_{j=1}^N\int_{E_j}\mathcal W_j(\theta)\,d\theta
=\int_I\mathcal V(\theta)\,d\theta.
\end{equation}

This equality holds for every measurable legal partition, including disconnected or empty cells, and its
right-hand side is independent of the legal pivot assignment.


\end{lemma}

\begin{proof}
For each \(j\), let

\begin{equation}
\label{eq:appendix-display-135}
\Omega_j:=\{\theta\in\Theta:F_j(\theta)\neq0\}.
\end{equation}

The setting-defined feature functions are continuous, so \(\Omega_j\) is relatively open. On
\(\Omega_j\times\mathbb R^{N-1}\), the quotient formula for \(T_j\) and its fixed-beta theta derivative are
continuous. Define on the entire product \(\Theta\times[-R,R]^{N-1}\)

\begin{equation}
\label{eq:appendix-display-136}
g_j(\theta,\beta)
:=
\begin{cases}
\mathbf1\{\lvert T_j(\theta,\beta)\rvert\leq R\}
\lvert\partial_\theta T_j(\theta,\beta)\rvert,
&\theta\in\Omega_j,\\
0,&\theta\notin\Omega_j.
\end{cases}
\end{equation}

This piecewise function is Borel and nonnegative: it is a Borel function on the Borel open set
\(\Omega_j\times[-R,R]^{N-1}\) and is zero on its Borel complement. Tonelli's measurability conclusion gives
the Borel function

\begin{equation}
\label{eq:appendix-display-137}
\widehat{\mathcal W}_j(\theta)
:=\int_{[-R,R]^{N-1}}g_j(\theta,\beta)\,d\beta
\quad(\theta\in\Theta),
\end{equation}

with values in \([0,\infty]\). Since legality gives \(E_j\subseteq\Omega_j\), the restriction of
\(\widehat{\mathcal W}_j\) to \(E_j\) is exactly \(\mathcal W_j\). It is therefore Lebesgue measurable even
when \(E_j\) is not Borel.

For each \(\theta\in E_j\), the beta points selected by the indicator are exactly
\(\mathcal B_{j,\theta}\). Taking \(h\equiv1\) in
Lemma~\ref{lem:step-004-s2-normal-density} gives the pointwise identity

\begin{equation}
\label{eq:appendix-display-138}
\mathcal W_j(\theta)=\mathcal V(\theta).
\end{equation}

Thus, pointwise on \(I\),

\begin{equation}
\label{eq:appendix-display-139}
\mathcal V(\theta)
=\sum_{j=1}^N\mathbf1_{E_j}(\theta)
\widehat{\mathcal W}_j(\theta).
\end{equation}

The right-hand side is a finite sum of nonnegative Lebesgue-measurable functions, so this identity proves
measurability of the moving-section mass \(\mathcal V\) without invoking a separate moving-hyperplane,
coarea, or root theorem.

Finite additivity for the exact disjoint partition now yields

\begin{equation}
\label{eq:appendix-display-140}
\begin{aligned}
\sum_{j=1}^N\int_{E_j}\mathcal W_j(\theta)\,d\theta
&=\sum_{j=1}^N\int_{E_j}\mathcal V(\theta)\,d\theta\\
&=\int_I\mathcal V(\theta)\,d\theta.
\end{aligned}
\end{equation}

Every term is nonnegative, so the calculation is valid if the common value is \(+\infty\); no expression of
the form \(\infty-\infty\) occurs. The cells are disjoint, so different legal charts never double charge a
single theta in the sum. Although two pivot coordinates may both be legal at the same theta, the exact
fixed-section equality shows that either chosen chart represents the same \(\mathcal V(\theta)\). Empty cells
contribute zero. Included interval endpoints remain literal members of their cells and excluded endpoints do
not enter; no closure of \(I\) is substituted.

For \(N=1\), the sole beta integral is integration over \(\mathbb R^0\) with mass one, the partition has the
single legal cell \(E_1=I\), and the pointwise equality is exactly the \(\mathcal H^0\) evaluation established
in Lemma~\ref{lem:step-004-s2-normal-density}. 
\end{proof}

\begin{proposition}[Coordinate-free first affine swept-area inequality]
\label{prop:step-004-s2-affine-swept-area}

Under Assumption~\ref{assump:anchored-derivative-closure}, the Exhausted affine pivot-sweep
Proposition~\ref{prop:step-003-pivot-sweep}, and
Lemma~\ref{lem:step-004-s2-measurable-partition}, for every
\(\mu\in\mathcal D_{N,R,\kappa}\) and every interval \(I\subseteq\Theta\) with
\(\lvert I\rvert>0\), retaining the interval's literal endpoint convention,

\begin{equation}
\label{eq:appendix-display-141}
\begin{aligned}
\Pr_{\alpha\sim\mu}[\exists\theta\in I:\phi_\alpha(\theta)=0]
&\leq
\kappa\int_I\int_{H_\theta\cap[-R,R]^N}
\frac{\lvert F_0'(\theta)+\langle a,F'(\theta)\rangle\rvert}
{\lVert F(\theta)\rVert_2}
\,d\mathcal H^{N-1}(a)\,d\theta.
\end{aligned}
\end{equation}

The inequality is in \([0,\infty]\), keeps the original coefficient coordinates and deterministic offset,
and has the literal factor \(\kappa\). It applies to every arbitrary correlated full joint law in
\(\mathcal D_{N,R,\kappa}\), and includes all endpoint, root-multiplicity, cube-boundary, empty-section,
zero-velocity, and \(N=1\) cases covered by its stated inputs.


\end{proposition}

\begin{proof}
Fix the deterministic setting instance. Then fix an arbitrary
\(\mu\in\mathcal D_{N,R,\kappa}\) and, after the law, an arbitrary interval
\(I\subseteq\Theta\) with \(\lvert I\rvert>0\) and its literal endpoint convention.

At least one measurable legal partition exists directly from the primitive anchor: take

\begin{equation}
\label{eq:appendix-display-142}
E_{j_*}=I,
\qquad
E_j=\varnothing\quad(j\neq j_*),
\end{equation}

because \(F_{j_*}\equiv1\). More generally, fix any measurable legal partition. The indicator form of the
Proposition~\ref{prop:step-003-pivot-sweep} gives, without any new event or root argument,

\begin{equation}
\label{eq:appendix-display-143}
\begin{aligned}
\Pr_{\alpha\sim\mu}[\exists\theta\in I:\phi_\alpha(\theta)=0]
&\leq
\kappa\sum_{j=1}^N\int_{E_j}\int_{[-R,R]^{N-1}}
\mathbf1\{\lvert T_j(\theta,\beta)\rvert\leq R\}
\lvert\partial_\theta T_j(\theta,\beta)\rvert
\,d\beta\,d\theta.
\end{aligned}
\end{equation}

Lemma~\ref{lem:step-004-s2-measurable-partition} identifies the entire chart sum, exactly and in the extended
nonnegative reals, with

\begin{equation}
\label{eq:appendix-display-144}
\int_I\int_{H_\theta\cap[-R,R]^N}
\frac{\lvert F_0'(\theta)+\langle a,F'(\theta)\rangle\rvert}
{\lVert F(\theta)\rVert_2}
\,d\mathcal H^{N-1}(a)\,d\theta.
\end{equation}

Substitution proves the claimed inequality with no remainder, pivot factor, chart-count factor, or change to
\(\kappa\). Proposition~\ref{prop:step-003-pivot-sweep} already treats the original full joint density, so the deterministic
fixed-section conversion neither conditions on nor marginalizes a coordinate and preserves arbitrary
correlation. It also deletes no coefficient, parameter endpoint, root, cube face, or corner. Tangent, multiple,
and persistent roots remain covered exactly as in that proposition; this proof has invoked no second
root theorem. Because \(\mu\) and then \(I\) were arbitrary, the quantifier order is preserved.

\end{proof}


\begin{proof}[Section-formula synthesis]
Fix the deterministic instance, then a law, a positive-length interval, and
a measurable legal pivot partition.  For every $\theta\in E_j$,
Lemma~\ref{lem:step-004-s2-section-parametrization} gives a bijection from
\begin{equation}
\label{eq:appendix-display-145}
\{\beta\in[-R,R]^{N-1}:|T_j(\theta,\beta)|\leq R\}
\end{equation}
onto $H_\theta\cap[-R,R]^N$ with graph Jacobian
$\|F(\theta)\|_2/|F_j(\theta)|$.  On the same chart,
Lemma~\ref{lem:step-004-s2-normal-density} gives
\begin{equation}
\label{eq:appendix-display-146}
\partial_\theta T_j(\theta,\beta)
=-
\frac{F_0'(\theta)+\langle\Psi_j(\theta,\beta),F'(\theta)\rangle}
{F_j(\theta)}.
\end{equation}
The pivot denominator therefore cancels exactly:
\begin{equation}
\label{eq:appendix-display-147}
|\partial_\theta T_j|\,d\beta
=
\frac{|F_0'+\langle a,F'\rangle|}{\|F\|_2}
\,d\mathcal H^{N-1}(a).
\end{equation}
Lemma~\ref{lem:step-004-s2-measurable-partition} proves measurability and
sums this equality over the finite partition in the extended nonnegative
reals.  Substitution into the indicator form of
Proposition~\ref{prop:step-003-pivot-sweep} is precisely
Proposition~\ref{prop:step-004-s2-affine-swept-area}.  It preserves the
original coefficient coordinates, the single correlated joint law, literal
$\kappa$, all endpoint conventions, empty cells, either pivot sign, zero
velocity, cube boundaries, and the $N=1$ convention.
\end{proof}

\subsection{Translated Cube Sections}\label{app:step-005}

\begin{theorem}[Brunn--Minkowski inequality]\label{thm:brunn-minkowski}
Let $d\geq1$, let $X,Y$ be nonempty compact subsets of a $d$-dimensional Euclidean affine space, and let
$0\leq\lambda\leq1$. Then
\begin{equation}
\label{eq:appendix-display-148}
\operatorname{vol}_d((1-\lambda)X+\lambda Y)^{1/d}
\geq(1-\lambda)\operatorname{vol}_d(X)^{1/d}
+\lambda\operatorname{vol}_d(Y)^{1/d}.
\end{equation}
The statement permits either input to have zero $d$-volume \citep{gardner2002brunn}.
\end{theorem}

\begin{theorem}[Ball's central cube-section theorem]\label{thm:ball-cube-slicing}
For $n\geq2$, let $Q_n=[-1/2,1/2]^n$. For every Euclidean hyperplane $L$ through the origin,
\begin{equation}
\label{eq:appendix-display-149}
\mathcal H^{n-1}(Q_n\cap L)\leq\sqrt2.
\end{equation}
This is the main cube-slicing theorem of \citet{ball1986cube}.
\end{theorem}

\begin{lemma}[Central maximality of parallel cube sections]
\label{lem:step-005-central-maximality}
Under Assumption~\ref{assump:parameter-regime}, if \(N\geq2\), \(u\in\mathbb R^N\) satisfies \(\|u\|_2=1\), and

\begin{equation}
\label{eq:appendix-display-150}
v(t):=\mathcal H^{N-1}\bigl([-R,R]^N\cap(u^\perp+tu)\bigr),
\qquad t\in\mathbb R,
\end{equation}

then the function \(t\mapsto v(t)^{1/(N-1)}\) is concave on the closed interval on which the sections are nonempty, \(v(-t)=v(t)\), and

\begin{equation}
\label{eq:appendix-display-151}
v(t)\leq v(0)\qquad\text{for every }t\in\mathbb R.
\end{equation}


\end{lemma}

\begin{proof}
Put \(K=[-R,R]^N\), set \(d=N-1\), and identify each affine slice isometrically with a subset of \(u^\perp\) by defining

\begin{equation}
\label{eq:appendix-display-152}
A_t:=\{x\in u^\perp:x+tu\in K\}.
\end{equation}

The translation \(x\mapsto x+tu\) is a Euclidean isometry from \(A_t\) onto \(K\cap(u^\perp+tu)\). Therefore

\begin{equation}
\label{eq:appendix-display-153}
v(t)=\operatorname{vol}_d(A_t).
\end{equation}

The projection of \(K\) onto the line spanned by \(u\) is

\begin{equation}
\label{eq:appendix-display-154}
\bigl\{\langle x,u\rangle:x\in K\bigr\}
=\left[-R\sum_{i=1}^N|u_i|,\ R\sum_{i=1}^N|u_i|\right].
\end{equation}

Indeed, the two extrema follow by choosing each cube coordinate with the sign opposite to, or equal to, the sign of \(u_i\), and convexity fills the interval between them. Thus this is exactly the closed interval of nonempty sections; in particular, its endpoints are retained.

Fix \(s,t\) in this interval and \(0\leq\lambda\leq1\). If \(x_s\in A_s\) and \(x_t\in A_t\), convexity of \(K\) gives

\begin{equation}
\label{eq:appendix-display-155}
(1-\lambda)(x_s+su)+\lambda(x_t+tu)
=\bigl((1-\lambda)x_s+\lambda x_t\bigr)
 +\bigl((1-\lambda)s+\lambda t\bigr)u\in K.
\end{equation}

Because \((1-\lambda)x_s+\lambda x_t\in u^\perp\), this proves

\begin{equation}
\label{eq:appendix-display-156}
(1-\lambda)A_s+\lambda A_t
\subseteq A_{(1-\lambda)s+\lambda t}.
\end{equation}

Theorem~\ref{thm:brunn-minkowski} and monotonicity of Euclidean volume now yield

\begin{equation}
\label{eq:appendix-display-157}
v((1-\lambda)s+\lambda t)^{1/d}
\geq(1-\lambda)v(s)^{1/d}+\lambda v(t)^{1/d}.
\end{equation}

This argument remains valid when an endpoint slice has zero \(d\)-volume; Brunn--Minkowski does not require positive input volume. Hence \(v^{1/d}\) is concave on the full closed interval of nonempty sections, including tangent slices supported on faces, lower-dimensional faces, edges, or corners.

Central symmetry \(K=-K\) gives \(A_{-t}=-A_t\), so \(v(-t)=v(t)\). Applying concavity to the midpoint \(0=(t+(-t))/2\) gives, for every \(t\) in the support interval,

\begin{equation}
\label{eq:appendix-display-158}
v(0)^{1/d}
\geq\frac12v(t)^{1/d}+\frac12v(-t)^{1/d}
=v(t)^{1/d}.
\end{equation}

Since \(d\geq1\), raising both sides to the power \(d\) gives \(v(t)\leq v(0)\). For \(t\) outside the support interval the translated section is empty, so \(v(t)=0\leq v(0)\). This proves the assertion for all real offsets, without extending the concave function by zero outside an endpoint where a facet section may have positive measure.
\end{proof}

\begin{lemma}[Euclidean scaling of Ball's cube-slicing bound]
\label{lem:step-005-ball-scaling}
Under Assumption~\ref{assump:parameter-regime}, if \(N\geq2\) and \(u\in\mathbb R^N\) satisfies \(\|u\|_2=1\), then

\begin{equation}
\label{eq:appendix-display-159}
\mathcal H^{N-1}\bigl([-R,R]^N\cap u^\perp\bigr)
\leq\sqrt2(2R)^{N-1}.
\end{equation}


\end{lemma}

\begin{proof}
Let \(Q_N=[-1/2,1/2]^N\), which has \(N\)-dimensional volume one. Theorem~\ref{thm:ball-cube-slicing}, in the same Euclidean metric, gives

\begin{equation}
\label{eq:appendix-display-160}
\mathcal H^{N-1}(Q_N\cap u^\perp)\leq\sqrt2.
\end{equation}

Assumption~\ref{assump:parameter-regime} gives \(R>0\), and scalar dilation by \(2R\) satisfies the literal set identity

\begin{equation}
\label{eq:appendix-display-161}
(2R)(Q_N\cap u^\perp)
=([-R,R]^N)\cap u^\perp.
\end{equation}

For every subset \(E\) of a Euclidean \((N-1)\)-plane and every \(c>0\), Hausdorff measure obeys the exact metric scaling identity

\begin{equation}
\label{eq:appendix-display-162}
\mathcal H^{N-1}(cE)=c^{N-1}\mathcal H^{N-1}(E).
\end{equation}

Taking \(c=2R\) proves

\begin{equation}
\label{eq:appendix-display-163}
\mathcal H^{N-1}\bigl([-R,R]^N\cap u^\perp\bigr)
=(2R)^{N-1}\mathcal H^{N-1}(Q_N\cap u^\perp)
\leq\sqrt2(2R)^{N-1}.
\end{equation}

No normalization factor is hidden: both the source and target use Euclidean Hausdorff measure, and a scalar dilation multiplies every tangent length by exactly \(2R\).
\end{proof}

\begin{proposition}[Uniform translated cube-section certificate]
\label{prop:step-005-translated-section-certificate}
Under Assumption~\ref{assump:parameter-regime} and, when \(N\geq2\), Lemmas~\ref{lem:step-005-central-maximality} and \ref{lem:step-005-ball-scaling}, for every Euclidean unit vector \(u\in\mathbb R^N\) and every \(t\in\mathbb R\),

\begin{equation}
\label{eq:appendix-display-164}
\mathcal H^{N-1}\bigl([-R,R]^N\cap(u^\perp+tu)\bigr)
\leq
\mathcal H^{N-1}\bigl([-R,R]^N\cap u^\perp\bigr)
\leq\sqrt2(2R)^{N-1}.
\end{equation}

Moreover, if a setting-defined affine root section has \(F(\theta)\ne0\), then with

\begin{equation}
\label{eq:appendix-display-165}
u_\theta:=\frac{F(\theta)}{\|F(\theta)\|_2},
\qquad
t_\theta:=-\frac{F_0(\theta)}{\|F(\theta)\|_2},
\end{equation}

one has the exact same-normal identity

\begin{equation}
\label{eq:appendix-display-166}
H_\theta=u_\theta^\perp+t_\theta u_\theta
\end{equation}

and therefore

\begin{equation}
\label{eq:appendix-display-167}
\mathcal H^{N-1}\bigl(H_\theta\cap[-R,R]^N\bigr)
\leq
\mathcal H^{N-1}\bigl(F(\theta)^\perp\cap[-R,R]^N\bigr)
\leq\sqrt2(2R)^{N-1}.
\end{equation}


\end{proposition}

\begin{proof}
First suppose \(N\geq2\). Lemma~\ref{lem:step-005-central-maximality} bounds every translated section by its parallel central section, including support endpoints and empty sections. Lemma~\ref{lem:step-005-ball-scaling} bounds that central section uniformly over the Euclidean unit normal \(u\). Their composition proves the asserted two inequalities.

Now suppose \(N=1\). A Euclidean unit normal is \(u=1\) or \(u=-1\), and \(u^\perp=\{0\}\). Thus \(u^\perp+tu=\{tu\}\). Its intersection with \([-R,R]\) is either one point or empty, while the central section is the singleton \(\{0\}\). Since \(\mathcal H^0\) counts a singleton as one and the empty set as zero,

\begin{equation}
\label{eq:appendix-display-168}
\mathcal H^0([-R,R]\cap\{tu\})
\leq1
=\mathcal H^0([-R,R]\cap\{0\})
\leq\sqrt2(2R)^0.
\end{equation}

This proves both inequalities literally in dimension one.

It remains to verify that the actual affine section is the same object, rather than a substituted central section. For \(F(\theta)\ne0\), division of

\begin{equation}
\label{eq:appendix-display-169}
F_0(\theta)+\langle a,F(\theta)\rangle=0
\end{equation}

by \(\|F(\theta)\|_2\) shows that \(a\in H_\theta\) exactly when

\begin{equation}
\label{eq:appendix-display-170}
\langle a,u_\theta\rangle=t_\theta.
\end{equation}

Every \(a\in\mathbb R^N\) has the orthogonal decomposition

\begin{equation}
\label{eq:appendix-display-171}
a=\bigl(a-\langle a,u_\theta\rangle u_\theta\bigr)
 +\langle a,u_\theta\rangle u_\theta,
\end{equation}

so this level set is precisely \(u_\theta^\perp+t_\theta u_\theta\). Its parallel central section is precisely \(u_\theta^\perp=F(\theta)^\perp\). In the actual branch, the setting anchor gives \(F_{j_*}(\theta)=1\), hence \(F(\theta)\ne0\) at every \(\theta\). In particular, when \(N=1\), \(F(\theta)\) is a nonzero scalar, \(H_\theta=\{-F_0(\theta)/F(\theta)\}\), and its intersection with \([-R,R]\) is a singleton or empty, as proved in the preceding dimension-one calculation.

The proof uses no regularity of \(t_\theta\) as \(\theta\) varies. A central offset, a noncentral offset, tangency to a cube face, intersection through a face, edge, or corner, a coordinate-aligned normal, a diagonal normal, a support endpoint, and a zero-measure or empty section are all already covered by the quantified geometric argument.
\end{proof}


\begin{proof}[Synthesis of the translated-section bound]
For $N\geq2$, Lemma~\ref{lem:step-005-central-maximality}, through the
stated Brunn--Minkowski theorem, bounds every translated section by the
parallel central section.  Lemma~\ref{lem:step-005-ball-scaling} applies
Ball's central cube-section theorem with the exact Euclidean Hausdorff
scaling from the unit-volume cube to $[-R,R]^N$.  Proposition~\ref{prop:step-005-translated-section-certificate}
also proves the $N=1$ case and identifies the actual affine section using
\begin{equation}
\label{eq:appendix-display-172}
u_\theta=\frac{F(\theta)}{\|F(\theta)\|_2},
\qquad
t_\theta=-\frac{F_0(\theta)}{\|F(\theta)\|_2}.
\end{equation}
Consequently, for every $N\geq1$, every $R>0$, and every actual section,
\begin{equation}
\label{eq:appendix-display-173}
\mathcal H^{N-1}(H_\theta\cap[-R,R]^N)
\leq
\mathcal H^{N-1}(F(\theta)^\perp\cap[-R,R]^N)
\leq\sqrt2(2R)^{N-1}.
\end{equation}
The conclusion includes every orientation and offset, empty sections, support
endpoints, cube faces and corners, and the stated zero-dimensional convention.
\end{proof}

\subsection{Affine Root-Section Velocity}\label{app:step-006}

\begin{lemma}[Root-section Euclidean coupling]
\label{lem:step-006-s2-root-coupling}

Under Assumptions~\ref{assump:parameter-regime} and
\ref{assump:anchored-derivative-closure}, if \(\theta\in\Theta\) and
\(a\in H_\theta\cap[-R,R]^N\), then

\begin{equation}
\lvert F_0(\theta)\rvert
=\lvert\langle a,F(\theta)\rangle\rvert
\leq R\sqrt N\,\lVert F(\theta)\rVert_2,
\label{eq:step-006-root-f0}
\end{equation}

\begin{equation}
\lVert\widetilde F(\theta)\rVert_2
\leq\sqrt{1+NR^2}\,\lVert F(\theta)\rVert_2,
\label{eq:step-006-augmented-bound}
\end{equation}

and

\begin{equation}
\lVert(1,a)\rVert_2\leq\sqrt{1+NR^2}.
\label{eq:step-006-coefficient-bound}
\end{equation}


\end{lemma}

\begin{proof}
Membership in the actual affine root section means, by the setting definition of
\(H_\theta\),

\begin{equation}
\label{eq:appendix-display-177}
F_0(\theta)+\langle a,F(\theta)\rangle=0.
\end{equation}

Thus, pointwise for this identical \(a\) and \(\theta\),

\begin{equation}
F_0(\theta)=-\langle a,F(\theta)\rangle,
\qquad
\lvert F_0(\theta)\rvert=\lvert\langle a,F(\theta)\rangle\rvert.
\label{eq:step-006-root-identity}
\end{equation}

No bound on \(F_0\) away from the root section has been introduced. Cube membership gives, coordinate by
coordinate, \(\lvert a_i\rvert\leq R\), and hence

\begin{equation}
\lVert a\rVert_2^2=\sum_{i=1}^N a_i^2\leq NR^2,
\qquad
\lVert a\rVert_2\leq R\sqrt N.
\label{eq:step-006-cube-norm}
\end{equation}

Applying Euclidean Cauchy--Schwarz to the right-hand side of
\eqref{eq:step-006-root-identity} and then using
\eqref{eq:step-006-cube-norm} proves
\eqref{eq:step-006-root-f0}:

\begin{equation}
\label{eq:appendix-display-180}
\lvert F_0(\theta)\rvert
=\lvert\langle a,F(\theta)\rangle\rvert
\leq\lVert a\rVert_2\lVert F(\theta)\rVert_2
\leq R\sqrt N\,\lVert F(\theta)\rVert_2.
\end{equation}

Squaring this derived bound, and using the setting's actual augmented feature vector rather than a surrogate,
gives

\begin{equation}
\label{eq:appendix-display-181}
\begin{aligned}
\lVert\widetilde F(\theta)\rVert_2^2
&=\lvert F_0(\theta)\rvert^2+\lVert F(\theta)\rVert_2^2\\
&\leq NR^2\lVert F(\theta)\rVert_2^2+\lVert F(\theta)\rVert_2^2\\
&=(1+NR^2)\lVert F(\theta)\rVert_2^2.
\end{aligned}
\end{equation}

Both sides are nonnegative, so taking square roots proves
\eqref{eq:step-006-augmented-bound}. Finally, the same cube calculation gives

\begin{equation}
\label{eq:appendix-display-182}
\lVert(1,a)\rVert_2^2=1+\lVert a\rVert_2^2\leq1+NR^2,
\end{equation}

and taking the nonnegative square root proves
\eqref{eq:step-006-coefficient-bound}. Every constant in
\eqref{eq:step-006-root-f0}--\eqref{eq:step-006-coefficient-bound} has been derived from the actual
root equation and cube membership. 
\end{proof}

\begin{proposition}[Exact affine root-section velocity]
\label{prop:step-006-s2-affine-velocity}

Under Assumptions~\ref{assump:parameter-regime} and
\ref{assump:anchored-derivative-closure}, Lemma~\ref{lem:step-006-s2-root-coupling}, and
Lemma~\ref{lem:step-001-height}, if \(\theta\in\Theta\) and
\(a\in H_\theta\cap[-R,R]^N\), then

\begin{equation}
F_0'(\theta)+\langle a,F'(\theta)\rangle
=\langle(1,a),B(\theta)\widetilde F(\theta)\rangle,
\label{eq:step-006-derivative-identity}
\end{equation}

and

\begin{equation}
\frac{\lvert F_0'(\theta)+\langle a,F'(\theta)\rangle\rvert}
{\lVert F(\theta)\rVert_2}
\leq(1+NR^2)\widehat\Lambda_{B,T}.
\label{eq:step-006-velocity-bound}
\end{equation}

More precisely, the coefficient and feature square-root factors in
\eqref{eq:step-006-velocity-bound} multiply exactly as

\begin{equation}
\sqrt{1+NR^2}\,\sqrt{1+NR^2}=1+NR^2.
\label{eq:step-006-factor-product}
\end{equation}


\end{proposition}

\begin{proof}
The vector \(a\) is held fixed in this pointwise velocity evaluation; no
\(\theta\)-dependent section parametrization and no derivative of \(a\) is introduced. The exact primitive
closure identity is

\begin{equation}
\label{eq:appendix-display-186}
\widetilde F'(\theta)
=\bigl(F_0'(\theta),F'(\theta)\bigr)
=B(\theta)\widetilde F(\theta).
\end{equation}

Taking its Euclidean inner product with the identical augmented coefficient vector \((1,a)\) gives

\begin{equation}
\label{eq:appendix-display-187}
\begin{aligned}
F_0'(\theta)+\langle a,F'(\theta)\rangle
&=\langle(1,a),\bigl(F_0'(\theta),F'(\theta)\bigr)\rangle\\
&=\langle(1,a),\widetilde F'(\theta)\rangle\\
&=\langle(1,a),B(\theta)\widetilde F(\theta)\rangle.
\end{aligned}
\end{equation}

This proves the identity \eqref{eq:step-006-derivative-identity} on the
actual objects. Euclidean Cauchy--Schwarz and the definition of
the induced operator norm then give

\begin{equation}
\begin{aligned}
\lvert F_0'(\theta)+\langle a,F'(\theta)\rangle\rvert
&=\lvert\langle(1,a),B(\theta)\widetilde F(\theta)\rangle\rvert\\
&\leq\lVert(1,a)\rVert_2\,
       \lVert B(\theta)\widetilde F(\theta)\rVert_2\\
&\leq\lVert(1,a)\rVert_2\,
       \lVert B(\theta)\rVert_{\mathrm{op}}\,
       \lVert\widetilde F(\theta)\rVert_2.
\end{aligned}
\label{eq:step-006-operator-chain}
\end{equation}

Assumption~\ref{assump:anchored-derivative-closure} supplies the legal normalization

\begin{equation}
\lVert F(\theta)\rVert_2\geq\lvert F_{j_*}(\theta)\rvert=1.
\label{eq:step-006-anchor-normalization}
\end{equation}

Lemma~\ref{lem:step-001-height} yields

\begin{equation}
\lVert B(\theta)\rVert_{\mathrm{op}}
\leq\sup_{\vartheta\in\Theta}\lVert B(\vartheta)\rVert_{\mathrm{op}}
\leq\widehat\Lambda_{B,T}.
\label{eq:step-006-matrix-bound}
\end{equation}

Divide \eqref{eq:step-006-operator-chain} by the positive quantity in
\eqref{eq:step-006-anchor-normalization}, and substitute
\eqref{eq:step-006-augmented-bound},
\eqref{eq:step-006-coefficient-bound}, and
\eqref{eq:step-006-matrix-bound} without dropping or absorbing
any term. The two separate square-root factors remain visible:

\begin{equation}
\begin{aligned}
\frac{\lvert F_0'(\theta)+\langle a,F'(\theta)\rangle\rvert}
     {\lVert F(\theta)\rVert_2}
&\leq
\frac{
\underbrace{\sqrt{1+NR^2}}_{\text{from }\lVert(1,a)\rVert_2}
\,\widehat\Lambda_{B,T}\,
\underbrace{\sqrt{1+NR^2}\lVert F(\theta)\rVert_2}_{\text{from }\lVert\widetilde F(\theta)\rVert_2}}
{\lVert F(\theta)\rVert_2}\\
&=\bigl(\sqrt{1+NR^2}\bigr)
   \bigl(\sqrt{1+NR^2}\bigr)\widehat\Lambda_{B,T}\\
&=(1+NR^2)\widehat\Lambda_{B,T}.
\end{aligned}
\label{eq:step-006-velocity-calculation}
\end{equation}

The cancellation in the first equality of
\eqref{eq:step-006-velocity-calculation} is exact because the anchor proves
the denominator is nonzero. Equation~\eqref{eq:step-006-velocity-calculation}
proves \eqref{eq:step-006-velocity-bound} and
\eqref{eq:step-006-factor-product}, with no residual or hidden factor. 
\end{proof}

\begin{proposition}[Boundary and zero-certificate consistency]
\label{prop:step-006-s2-boundary}

Under Assumptions~\ref{assump:parameter-regime} and
\ref{assump:anchored-derivative-closure}, use
Lemma~\ref{lem:step-006-s2-root-coupling},
Proposition~\ref{prop:step-006-s2-affine-velocity}, and
Lemma~\ref{lem:step-001-height}. Their conclusions remain valid when
\(N=1\), \(F_0\equiv0\), \(a=0\), the root section is empty,
\(\theta\) is an endpoint of \(\Theta\), \(R>0\), or
\(\widehat\Lambda_{B,T}=0\). In the zero-certificate case, for every
\(\theta\in\Theta\),

\begin{equation}
\label{eq:appendix-display-192}
B(\theta)=0,
\qquad
\widetilde F'(\theta)=0,
\end{equation}

and for every \(a\in H_\theta\cap[-R,R]^N\),

\begin{equation}
\frac{\lvert F_0'(\theta)+\langle a,F'(\theta)\rangle\rvert}
{\lVert F(\theta)\rVert_2}=0.
\label{eq:step-006-zero-velocity}
\end{equation}


\end{proposition}

\begin{proof}
Each requested case is a direct specialization of the same pointwise identities.

1. If \(N=1\), cube membership gives \(\lVert a\rVert_2=\lvert a\rvert\leq R\), so the two factors are
   \(\sqrt{1+R^2}\), and their product is exactly \(1+R^2\). The anchor is the sole feature coordinate and
   still gives \(\lVert F(\theta)\rVert_2\geq1\). No positive-dimensional section geometry is used.

2. If \(F_0\equiv0\), then at every root-section point
   \(\langle a,F(\theta)\rangle=0\) and
   \(\lVert\widetilde F(\theta)\rVert_2=\lVert F(\theta)\rVert_2\) exactly. Also \(F_0'=0\), so
   \eqref{eq:step-006-derivative-identity}
   specializes to
\begin{equation}
\label{eq:appendix-indented-002}
   \langle a,F'(\theta)\rangle
   =\langle(1,a),B(\theta)\widetilde F(\theta)\rangle.
\end{equation}
   The general velocity bound remains valid. If only \(F_0(\theta)=0\) at the tested point, the root and
   augmented-norm conclusions remain exact there; no derivative conclusion about \(F_0'\) is inferred.

3. If \(a=0\) belongs to the tested section, the actual root equation gives \(F_0(\theta)=0\), while
   \(\lVert(1,a)\rVert_2=1\) and
   \(\lVert\widetilde F(\theta)\rVert_2=\lVert F(\theta)\rVert_2\). Membership at one point does not imply
   \(F_0'(\theta)=0\); the derivative numerator is controlled by the exact
   closure calculation from \eqref{eq:step-006-derivative-identity} through
   \eqref{eq:step-006-velocity-calculation}.

4. If \(H_\theta\cap[-R,R]^N\) is empty, the assertion quantified over its elements is vacuous. No
   normalization, section measure, or event is evaluated on a nonexistent coefficient point.

5. If \(\theta\) is an endpoint of the compact interval \(\Theta\), the exact closure identity remains valid
   because it is assumed on the open interval \(U\supseteq\Theta\). Thus all derivatives and every pointwise
   equality in the proposition is the restriction of an identity valid on an open neighborhood of the endpoint.

6. The primitive regime requires \(R>0\). Therefore \(R\sqrt N\), \(1+NR^2\), and
   \(\sqrt{1+NR^2}\) are finite, with the square-root factor strictly positive. The proof never divides by
   \(R\) and needs no additional lower bound or limiting convention for it.

7. If \(\widehat\Lambda_{B,T}=0\), the certificate and nonnegativity of the operator norm give, for
   every \(\theta\in\Theta\),
\begin{equation}
\label{eq:appendix-indented-003}
   0\leq\lVert B(\theta)\rVert_{\mathrm{op}}
   \leq\sup_{\vartheta\in\Theta}\lVert B(\vartheta)\rVert_{\mathrm{op}}
   \leq0.
\end{equation}
   Hence \(\lVert B(\theta)\rVert_{\mathrm{op}}=0\), so \(B(\theta)=0\). The exact primitive closure, rather
   than a law-level argument, now gives
\begin{equation}
\label{eq:appendix-indented-004}
   \widetilde F'(\theta)=B(\theta)\widetilde F(\theta)=0.
\end{equation}
   The primitive anchor still gives \(\lVert F(\theta)\rVert_2\geq1\), so
   the normalization in \eqref{eq:step-006-zero-velocity} is legal,
   while its numerator is exactly zero. This proves
   \eqref{eq:step-006-zero-velocity}. No root-event probability conclusion is asserted or
   needed in this local velocity step.

These checks use the identical pointwise objects and introduce no section-size, density, transversality,
root-simplicity, pivot, or probabilistic claim. 
\end{proof}


\begin{proof}[Synthesis of the affine normal-velocity bound]
Fix $\theta\in\Theta$ and $a\in H_\theta\cap[-R,R]^N$.
Lemma~\ref{lem:step-006-s2-root-coupling} gives
\begin{equation}
\label{eq:appendix-display-194}
|F_0(\theta)|=|\langle a,F(\theta)\rangle|
\leq R\sqrt N\,\|F(\theta)\|_2,
\qquad
\|\widetilde F(\theta)\|_2
\leq\sqrt{1+NR^2}\,\|F(\theta)\|_2.
\end{equation}
Proposition~\ref{prop:step-006-s2-affine-velocity} applies the closure
identity to the identical objects and obtains
\begin{equation}
\label{eq:appendix-display-195}
F_0'(\theta)+\langle a,F'(\theta)\rangle
=\langle(1,a),B(\theta)\widetilde F(\theta)\rangle.
\end{equation}
Together with $\|(1,a)\|_2\leq\sqrt{1+NR^2}$ and
Lemma~\ref{lem:step-001-height}, this yields
\begin{equation}
\label{eq:appendix-display-196}
\frac{|F_0'(\theta)+\langle a,F'(\theta)\rangle|}
{\|F(\theta)\|_2}
\leq(1+NR^2)\widehat\Lambda_{B,T}.
\end{equation}
Proposition~\ref{prop:step-006-s2-boundary} verifies endpoints, empty
sections, $a=0$, vanishing $F_0$, and zero certificate.  Since the pair
$(\theta,a)$ was arbitrary, this is the required bound on every actual affine
root section, without a section-size, density, or probability assertion.
\end{proof}

\subsection{General Affine Rate}\label{app:step-007}

\begin{lemma}[Fixed-law affine probability rate with literal constants]
\label{lem:step-007-s2-fixed-law-rate}
Under Assumption~\ref{assump:cube-density-laws}, the setting parameter definitions, and Propositions~\ref{prop:step-004-s2-affine-swept-area},
\ref{prop:step-005-translated-section-certificate}, and
\ref{prop:step-006-s2-affine-velocity}, if
\(\mu\in\mathcal D_{N,R,\kappa}\) and \(I\subseteq\Theta\) is an interval with
\(\lvert I\rvert>0\), then

\begin{equation}
\begin{aligned}
\Pr_{\alpha\sim\mu}[\exists\theta\in I:\phi_\alpha(\theta)=0]
&\leq
\kappa\int_I\int_{H_\theta\cap[-R,R]^N}
\frac{\lvert F_0'(\theta)+\langle a,F'(\theta)\rangle\rvert}
{\lVert F(\theta)\rVert_2}
\,d\mathcal H^{N-1}(a)\,d\theta\\
&\leq
\kappa\sqrt2(2R)^{N-1}(1+NR^2)
\widehat\Lambda_{B,T}\lvert I\rvert\\
&=
\frac{A(1+NR^2)\widehat\Lambda_{B,T}}{\sqrt2R}
\lvert I\rvert.
\end{aligned}
\label{eq:step-007-rate-chain}
\end{equation}

This holds for every arbitrarily correlated admissible law, with the endpoint and
zero-dimensional-section conventions.


\end{lemma}

\begin{proof}
Fix the deterministic presentation data. Next fix an arbitrary
\(\mu\in\mathcal D_{N,R,\kappa}\), and after the law fix an arbitrary interval
\(I\subseteq\Theta\) with \(\lvert I\rvert>0\). Proposition~\ref{prop:step-004-s2-affine-swept-area} gives the first line of
\eqref{eq:step-007-rate-chain} directly:

\begin{equation}
\Pr_{\alpha\sim\mu}[\exists\theta\in I:\phi_\alpha(\theta)=0]
\leq
\kappa\int_I\int_{H_\theta\cap[-R,R]^N}
\frac{\lvert F_0'(\theta)+\langle a,F'(\theta)\rangle\rvert}
{\lVert F(\theta)\rVert_2}
\,d\mathcal H^{N-1}(a)\,d\theta.
\label{eq:step-007-area-bound}
\end{equation}

The factor \(\kappa\) in \eqref{eq:step-007-area-bound} is already the full-joint-density domination. It is neither repeated nor
converted into a marginal bound.

For every fixed \(\theta\in I\) and every
\(a\in H_\theta\cap[-R,R]^N\), Proposition~\ref{prop:step-006-s2-affine-velocity} applies to this identical section point and gives

\begin{equation}
0\leq
\frac{\lvert F_0'(\theta)+\langle a,F'(\theta)\rangle\rvert}
{\lVert F(\theta)\rVert_2}
\leq(1+NR^2)\widehat\Lambda_{B,T}.
\label{eq:step-007-velocity-bound}
\end{equation}

Integrating \eqref{eq:step-007-velocity-bound} on the actual section,
including the empty-section case, yields

\begin{equation}
\begin{aligned}
&\int_{H_\theta\cap[-R,R]^N}
\frac{\lvert F_0'(\theta)+\langle a,F'(\theta)\rangle\rvert}
{\lVert F(\theta)\rVert_2}
\,d\mathcal H^{N-1}(a)\\
&\qquad\leq
(1+NR^2)\widehat\Lambda_{B,T}\,
\mathcal H^{N-1}(H_\theta\cap[-R,R]^N).
\end{aligned}
\label{eq:step-007-section-integrand-bound}
\end{equation}

Proposition~\ref{prop:step-005-translated-section-certificate} applies to the same
\(H_\theta\), coefficient cube, Euclidean metric, and Hausdorff convention, so

\begin{equation}
\mathcal H^{N-1}(H_\theta\cap[-R,R]^N)
\leq\sqrt2(2R)^{N-1}.
\label{eq:step-007-section-measure}
\end{equation}

Substituting \eqref{eq:step-007-section-measure} into
\eqref{eq:step-007-section-integrand-bound}, then integrating the resulting
constant over the literal interval \(I\), gives

\begin{equation}
\begin{aligned}
&\kappa\int_I\int_{H_\theta\cap[-R,R]^N}
\frac{\lvert F_0'(\theta)+\langle a,F'(\theta)\rangle\rvert}
{\lVert F(\theta)\rVert_2}
\,d\mathcal H^{N-1}(a)\,d\theta\\
&\qquad\leq
\kappa\int_I
\sqrt2(2R)^{N-1}(1+NR^2)\widehat\Lambda_{B,T}
\,d\theta\\
&\qquad=
\kappa\sqrt2(2R)^{N-1}(1+NR^2)
\widehat\Lambda_{B,T}\lvert I\rvert.
\end{aligned}
\label{eq:step-007-integrated-bound}
\end{equation}

There is no chart sum in
\eqref{eq:step-007-section-integrand-bound}--\eqref{eq:step-007-integrated-bound},
hence no chart-count factor. There is no residual, surrogate section,
second law domination, or hidden absorption. The factor in
\eqref{eq:step-007-integrated-bound} is the direct product of the actual
translated-section measure and the actual pointwise normal velocity.

It remains only to verify the final equality in
\eqref{eq:step-007-rate-chain}. Assumption~\ref{assump:parameter-regime} gives
\(R>0\), and the setting defines \(A=(2R)^N\kappa\). Therefore

\begin{equation}
\begin{aligned}
\kappa\sqrt2(2R)^{N-1}
&=\frac{\sqrt2}{2R}\,\kappa(2R)^N\\
&=\frac{\sqrt2}{2R}\,A\\
&=\frac{A}{\sqrt2R},
\end{aligned}
\label{eq:step-007-coefficient-algebra}
\end{equation}

where \(\sqrt2/(2R)=1/(\sqrt2R)\) follows from \(2=(\sqrt2)^2\) and \(R>0\). Multiplying
\eqref{eq:step-007-coefficient-algebra} by the unchanged
nonnegative factors
\((1+NR^2)\widehat\Lambda_{B,T}\lvert I\rvert\) proves the last line of
\eqref{eq:step-007-rate-chain}.

For \(N=1\), \eqref{eq:step-007-area-bound} uses the \(\mathcal H^0\)
counting convention and zero-dimensional beta mass one;
\eqref{eq:step-007-section-measure} reads \(0\) or
\(1\leq\sqrt2(2R)^0\); and \eqref{eq:step-007-velocity-bound} has the literal factor
\((1+R^2)\widehat\Lambda_{B,T}\). Empty sections contribute zero to
\eqref{eq:step-007-section-integrand-bound}. Included interval endpoints are
inherited from \eqref{eq:step-007-area-bound}, while changing endpoint
conventions does not alter \eqref{eq:step-007-integrated-bound}. No lower bound on
\(\lvert I\rvert\) beyond strict positivity appears, so the argument applies unchanged to arbitrarily short
positive-length intervals. Because the law and interval were arbitrary,
\eqref{eq:step-007-rate-chain} holds separately for every such
pair, without independence or a law/interval union bound. 
\end{proof}

\begin{lemma}[Exact interval-then-law capacity conversion]
\label{lem:step-007-s2-capacity-conversion}
Under Assumption~\ref{assump:cube-density-laws}, the setting parameter definitions,
Propositions~\ref{prop:step-004-s2-affine-swept-area},
\ref{prop:step-005-translated-section-certificate}, and
\ref{prop:step-006-s2-affine-velocity}, and
Lemma~\ref{lem:step-007-s2-fixed-law-rate},

\begin{equation}
C^{\mathrm{aff}}_{\mathcal D}(F_0,F;\Theta)
\leq
\frac{A(1+NR^2)\widehat\Lambda_{B,T}}{\sqrt2R}.
\label{eq:step-007-capacity-bound}
\end{equation}

The conversion first takes the setting-defined inner supremum over positive-length intervals for a fixed law
and then the outer supremum over admissible laws.


\end{lemma}

\begin{proof}
Fix an arbitrary \(\mu\in\mathcal D_{N,R,\kappa}\). For every interval
\(I\subseteq\Theta\) with \(\lvert I\rvert>0\),
Lemma~\ref{lem:step-007-s2-fixed-law-rate} and positivity of \(\lvert I\rvert\) give

\begin{equation}
\frac{
\Pr_{\alpha\sim\mu}[\exists\theta\in I:\phi_\alpha(\theta)=0]
}{\lvert I\rvert}
\leq
\frac{A(1+NR^2)\widehat\Lambda_{B,T}}{\sqrt2R}.
\label{eq:step-007-fixed-interval-ratio}
\end{equation}

The right-hand side is deterministic and independent of \(I\). Taking exactly the inner interval supremum in
the setting definition yields, for this fixed \(\mu\),

\begin{equation}
\sup_{\substack{I\subseteq\Theta\ \mathrm{interval}\\ \lvert I\rvert>0}}
\frac{
\Pr_{\alpha\sim\mu}[\exists\theta\in I:\phi_\alpha(\theta)=0]
}{\lvert I\rvert}
\leq
\frac{A(1+NR^2)\widehat\Lambda_{B,T}}{\sqrt2R}.
\label{eq:step-007-fixed-law-sup}
\end{equation}

Now take the outer supremum over \(\mu\in\mathcal D\). By the setting definition,

\begin{equation}
\begin{aligned}
C^{\mathrm{aff}}_{\mathcal D}(F_0,F;\Theta)
&=
\sup_{\mu\in\mathcal D}
\sup_{\substack{I\subseteq\Theta\ \mathrm{interval}\\ \lvert I\rvert>0}}
\frac{
\Pr_{\alpha\sim\mu}[\exists\theta\in I:\phi_\alpha(\theta)=0]
}{\lvert I\rvert}\\
&\leq
\frac{A(1+NR^2)\widehat\Lambda_{B,T}}{\sqrt2R}.
\end{aligned}
\label{eq:step-007-capacity-sup}
\end{equation}

Equations~\eqref{eq:step-007-fixed-interval-ratio}--\eqref{eq:step-007-capacity-sup}
divide only by the positive scalar \(\lvert I\rvert\). They do not form a simultaneous
random event, exchange random and deterministic suprema, or apply any law/interval union bound. Ordinary
probability for each fixed law is preserved. 
\end{proof}

\begin{proposition}[General affine coefficient-sweep rate and zero certificate]
\label{prop:step-007-s2-general-affine-rate}
Under Assumption~\ref{assump:cube-density-laws}, the setting parameter definitions, Propositions~\ref{prop:step-004-s2-affine-swept-area},
\ref{prop:step-005-translated-section-certificate}, and
\ref{prop:step-006-s2-affine-velocity}, and
Lemmas~\ref{lem:step-007-s2-fixed-law-rate} and
\ref{lem:step-007-s2-capacity-conversion}, every
\(\mu\in\mathcal D_{N,R,\kappa}\) and every interval
\(I\subseteq\Theta\) with \(\lvert I\rvert>0\) satisfy

\begin{equation}
\begin{aligned}
\Pr_{\alpha\sim\mu}[\exists\theta\in I:\phi_\alpha(\theta)=0]
&\leq
\kappa\int_I\int_{H_\theta\cap[-R,R]^N}
\frac{\lvert F_0'(\theta)+\langle a,F'(\theta)\rangle\rvert}
{\lVert F(\theta)\rVert_2}
\,d\mathcal H^{N-1}(a)\,d\theta\\
&\leq
\kappa\sqrt2(2R)^{N-1}(1+NR^2)
\widehat\Lambda_{B,T}\lvert I\rvert\\
&=
\frac{A(1+NR^2)\widehat\Lambda_{B,T}}{\sqrt2R}
\lvert I\rvert,
\end{aligned}
\label{eq:step-007-final-rate}
\end{equation}

and

\begin{equation}
C^{\mathrm{aff}}_{\mathcal D}(F_0,F;\Theta)
\leq
\frac{A(1+NR^2)\widehat\Lambda_{B,T}}{\sqrt2R}.
\label{eq:step-007-final-capacity}
\end{equation}

If \(\widehat\Lambda_{B,T}=0\), then for every such law and interval,

\begin{equation}
\Pr_{\alpha\sim\mu}[\exists\theta\in I:\phi_\alpha(\theta)=0]=0.
\label{eq:step-007-zero-probability}
\end{equation}

These conclusions retain arbitrary coefficient correlation, ordinary probability, uniformity over every
positive-length interval, the \(N=1\) zero-dimensional section convention, empty root sections, literal
interval endpoints, and the exact dependence on \(A,N,R,\kappa,\widehat\Lambda_{B,T}\), and
\(\lvert I\rvert\). There are no hidden constants. Once the supplied certificate is fixed, additional
dependence on \(q,M,\Delta\) is exactly degree zero.


\end{proposition}

\begin{proof}
Equation~\eqref{eq:step-007-final-rate} is
Lemma~\ref{lem:step-007-s2-fixed-law-rate}, and
\eqref{eq:step-007-final-capacity} is
Lemma~\ref{lem:step-007-s2-capacity-conversion}. Their proofs use the identical section, coefficient,
feature tuple, matrix certificate, Euclidean norm, and Hausdorff measure convention used by the three
propositions, so no change of objects or constants occurs.

For the zero-certificate case, substitute \(\widehat\Lambda_{B,T}=0\) into
\eqref{eq:step-007-final-rate}. Since
\(A=(2R)^N\kappa\) is finite, \(R>0\), and probabilities are nonnegative,

\begin{equation}
0
\leq
\Pr_{\alpha\sim\mu}[\exists\theta\in I:\phi_\alpha(\theta)=0]
\leq
\frac{A(1+NR^2)\cdot0}{\sqrt2R}\lvert I\rvert
=0.
\label{eq:step-007-zero-chain}
\end{equation}

Thus \eqref{eq:step-007-zero-probability} is an exact probability conclusion
for each law and interval. It is not inferred from a vacuous empty-section
statement: even if an actual root section is nonempty,
Proposition~\ref{prop:step-006-s2-affine-velocity} makes the nonnegative
integrand in \eqref{eq:step-007-area-bound} equal to zero when
\(\widehat\Lambda_{B,T}=0\), and Proposition~\ref{prop:step-004-s2-affine-swept-area} then bounds the actual root event by that zero integral.
No separate persistent-root or affine-hyperplane nullity result is imported.

The \(N=1\), empty-section, endpoint, and arbitrarily short positive-length interval cases were checked in
Lemma~\ref{lem:step-007-s2-fixed-law-rate}; positivity of interval length is the only condition used in the
capacity division. This completes every clause of the proposition. 
\end{proof}


\begin{proof}[Synthesis of the general affine rate]
Proposition~\ref{prop:step-004-s2-affine-swept-area} supplies the
normal-velocity integral with coefficient $\kappa$.  On the same sections,
Propositions~\ref{prop:step-006-s2-affine-velocity} and
\ref{prop:step-005-translated-section-certificate} supply, respectively,
\begin{equation}
\label{eq:appendix-display-212}
\frac{|F_0'+\langle a,F'\rangle|}{\|F\|_2}
\leq(1+NR^2)\widehat\Lambda_{B,T},
\qquad
\mathcal H^{N-1}(H_\theta\cap[-R,R]^N)
\leq\sqrt2(2R)^{N-1}.
\end{equation}
Their literal product in Lemma~\ref{lem:step-007-s2-fixed-law-rate} is
\begin{equation}
\label{eq:appendix-display-213}
\kappa\,[\sqrt2(2R)^{N-1}]
[(1+NR^2)\widehat\Lambda_{B,T}]\,|I|,
\end{equation}
and the only algebraic conversion is
\begin{equation}
\label{eq:appendix-display-214}
\kappa\sqrt2(2R)^{N-1}
=\frac{\sqrt2}{2R}\kappa(2R)^N
=\frac{A}{\sqrt2R}.
\end{equation}
Lemma~\ref{lem:step-007-s2-capacity-conversion} divides only by positive
$|I|$, takes the interval supremum for a fixed law, and only then takes the
law supremum.  Proposition~\ref{prop:step-007-s2-general-affine-rate}
combines these conclusions and proves the zero-certificate branch.  Thus no
chart-count factor, second density domination, probability conversion,
confidence parameter, independence condition, pivot margin, transversality,
or hidden constant enters the rate.
\end{proof}

\subsection{Homogeneous Projective Rate}\label{app:step-008}

\begin{lemma}[Exact homogeneous radial cancellation on the actual section]
\label{lem:step-008-s2-radial-cancellation}

Under Assumptions~\ref{assump:parameter-regime}, \ref{assump:balcan-common-chain}, and
\ref{assump:anchored-derivative-closure}, Lemma~\ref{lem:step-001-anchor}, and the exact specialization \(F_0\equiv0\), define

\begin{equation}
\label{eq:appendix-display-215}
r(\theta):=\lVert F(\theta)\rVert_2,
\qquad
\gamma_F(\theta):=\frac{F(\theta)}{r(\theta)}.
\end{equation}

Then, for every \(\theta\in\Theta\), \(r(\theta)\geq1\), both definitions and all differentiations are legal,
and the actual setting-defined coefficient section is

\begin{equation}
\label{eq:appendix-display-216}
H_\theta
=\{a\in\mathbb R^N:\langle a,F(\theta)\rangle=0\}
=F(\theta)^\perp
=\gamma_F(\theta)^\perp.
\end{equation}

For every \(a\in H_\theta\),

\begin{equation}
\label{eq:appendix-display-217}
\frac{\lvert\langle a,F'(\theta)\rangle\rvert}
{\lVert F(\theta)\rVert_2}
=\lvert\langle a,\gamma_F'(\theta)\rangle\rvert.
\end{equation}

This is an equality on the same actual coefficient section, with no bound on \(r\) beyond anchor legality and no
bound on \(r'\).


\end{lemma}

\begin{proof}
Lemma~\ref{lem:step-001-anchor} gives

\begin{equation}
\label{eq:appendix-display-218}
F_{j_*}(\theta)=1,
\qquad
r(\theta)=\lVert F(\theta)\rVert_2\geq1
\end{equation}

for every \(\theta\in\Theta\), and gives differentiability of \(\gamma_F\). Thus normalization and division by
\(r(\theta)\) are legal, including at interval endpoints by restriction from \(U\).

Because \(F_0\equiv0\), the setting definition of the actual section gives

\begin{equation}
\label{eq:appendix-display-219}
\begin{aligned}
H_\theta
&=\{a:F_0(\theta)+\langle a,F(\theta)\rangle=0\}\\
&=\{a:\langle a,F(\theta)\rangle=0\}\\
&=\{a:r(\theta)\langle a,\gamma_F(\theta)\rangle=0\}\\
&=\gamma_F(\theta)^\perp.
\end{aligned}
\end{equation}

This also proves \(H_\theta=F(\theta)^\perp\). In particular, for the same \(a\) appearing in the actual
section integral,

\begin{equation}
\label{eq:appendix-display-220}
\langle a,\gamma_F(\theta)\rangle=0.
\end{equation}

Differentiate the object identity \(F=r\gamma_F\):

\begin{equation}
\label{eq:appendix-display-221}
F'=r'\gamma_F+r\gamma_F'.
\end{equation}

Taking the inner product with that same \(a\in H_\theta\), dividing by the positive
\(r=\lVert F\rVert_2\), and using actual-section orthogonality gives

\begin{equation}
\label{eq:appendix-display-222}
\frac{\langle a,F'\rangle}{\lVert F\rVert_2}
=\frac{r'}r\langle a,\gamma_F\rangle+\langle a,\gamma_F'\rangle
=\langle a,\gamma_F'\rangle.
\end{equation}

Taking absolute values proves the asserted equality. The term involving \(r'\) vanishes algebraically; it is
not bounded, integrated, absorbed, or replaced by an amplitude condition. 
\end{proof}

\begin{lemma}[Literal homogeneous coefficient algebra]
\label{lem:step-008-s2-literal-algebra}

Under Assumption~\ref{assump:parameter-regime} and the setting definition
\(A=(2R)^N\kappa\),

\begin{equation}
\label{eq:appendix-display-223}
\kappa R\sqrt N\,\sqrt{2}(2R)^{N-1}
=A\sqrt{\frac N2}.
\end{equation}


\end{lemma}

\begin{proof}
Assumption~\ref{assump:parameter-regime} gives \(R>0\), \(N\geq1\), and
\(\kappa>0\). Hence every power and square root in the proposition is legal, and

\begin{equation}
\label{eq:appendix-display-224}
R(2R)^{N-1}=\frac{(2R)^N}{2}.
\end{equation}

Therefore

\begin{equation}
\label{eq:appendix-display-225}
\begin{aligned}
\kappa R\sqrt N\,\sqrt{2}(2R)^{N-1}
&=\kappa(2R)^N\frac{\sqrt N\,\sqrt{2}}{2}\\
&=\kappa(2R)^N\sqrt{\frac N2}\\
&=A\sqrt{\frac N2}.
\end{aligned}
\end{equation}

Every equality is literal; no term is dominated or hidden. 
\end{proof}

\begin{proposition}[Stationary projective branch is one fixed law-null central hyperplane]
\label{prop:step-008-s2-stationary-projective}

Under Assumptions~\ref{assump:parameter-regime}, \ref{assump:cube-density-laws}, and
\ref{assump:anchored-derivative-closure}, Lemma~\ref{lem:step-001-anchor}, and the exact
specialization \(F_0\equiv0\), suppose

\begin{equation}
\label{eq:appendix-display-226}
\Gamma_{\mathrm{proj}}(F)=0.
\end{equation}

Then there is a fixed unit vector \(\gamma_0\in\mathbb R^N\) such that
\(\gamma_F(\theta)=\gamma_0\) for every \(\theta\in\Theta\). For every nonempty interval
\(I\subseteq\Theta\),

\begin{equation}
\label{eq:appendix-display-227}
\left\{a\in\mathbb R^N:
\exists\theta\in I,\ \langle a,F(\theta)\rangle=0\right\}
=\gamma_0^\perp.
\end{equation}

Thus the union root event is one fixed proper nonempty central hyperplane, not an empty section, and every
\(\mu\in\mathcal D_{N,R,\kappa}\) satisfies

\begin{equation}
\label{eq:appendix-display-228}
\Pr_{\alpha\sim\mu}
[\exists\theta\in I:\langle\alpha,F(\theta)\rangle=0]=0.
\end{equation}

For \(N=1\), the same statement reads \(\gamma_0^\perp=\{0\}\), and the conclusion remains ordinary
probability zero under every admissible one-dimensional density.


\end{proposition}

\begin{proof}
By definition,

\begin{equation}
\label{eq:appendix-display-229}
\Gamma_{\mathrm{proj}}(F)
=\sup_{\theta\in\Theta}\lVert\gamma_F'(\theta)\rVert_2.
\end{equation}

All terms under the supremum are nonnegative. Hence
\(\Gamma_{\mathrm{proj}}(F)=0\) implies
\(\gamma_F'(\theta)=0\) for every \(\theta\in\Theta\). The set \(\Theta\) is a connected interval, and
\(\gamma_F\) is differentiable by Lemma~\ref{lem:step-001-anchor}. The coordinatewise
zero-derivative criterion therefore gives a fixed vector \(\gamma_0\) with
\(\gamma_F(\theta)=\gamma_0\) on all of \(\Theta\). Since every \(\gamma_F(\theta)\) is normalized,
\(\lVert\gamma_0\rVert_2=1\).

Lemma~\ref{lem:step-001-anchor} also gives
\(r(\theta)=\lVert F(\theta)\rVert_2\geq1\). Thus, for every \(\theta\),

\begin{equation}
\label{eq:appendix-display-230}
F(\theta)=r(\theta)\gamma_0,
\qquad
\langle a,F(\theta)\rangle=0
\ \Longleftrightarrow\
\langle a,\gamma_0\rangle=0.
\end{equation}

Every positive-length interval is nonempty, so taking the union over \(\theta\in I\) gives exactly
\(\gamma_0^\perp\). Because \(\gamma_0\neq0\), this is a proper hyperplane. Because \(0\in\gamma_0^\perp\),
it is central and nonempty. The conclusion depends only on the projective direction, so even an arbitrary
positive differentiable radial factor would leave this root set unchanged.

A proper hyperplane is \(N\)-dimensional Lebesgue-null. Assumption~\ref{assump:cube-density-laws} supplies one
full joint density \(f_\mu\), so

\begin{equation}
\label{eq:appendix-display-231}
\mu(\gamma_0^\perp)
=\int_{\gamma_0^\perp}f_\mu(a)\,d\operatorname{Leb}^N(a)
=0.
\end{equation}

This conclusion uses neither independence nor marginal densities. When \(N=1\), a unit
\(\gamma_0\) has either normal orientation and its orthogonal hyperplane is the singleton \(\{0\}\).
Its intersection with \([-R,R]\) is nonempty and has \(\mathcal H^0\)-mass one, but it has
one-dimensional Lebesgue measure zero, so every admissible density still assigns it probability zero.

\end{proof}

\begin{proposition}[Sharp pairwise homogeneous interval rate]
\label{prop:step-008-s2-pairwise-homogeneous-rate}

Under Assumptions~\ref{assump:parameter-regime}, \ref{assump:balcan-common-chain},
\ref{assump:cube-density-laws}, and \ref{assump:anchored-derivative-closure}, the current
Propositions~\ref{prop:step-001-projective},
\ref{prop:step-004-s2-affine-swept-area}, and
\ref{prop:step-005-translated-section-certificate}, Lemma~\ref{lem:step-001-anchor}, Lemmas~\ref{lem:step-008-s2-radial-cancellation} and
\ref{lem:step-008-s2-literal-algebra}, and
Proposition~\ref{prop:step-008-s2-stationary-projective}, if \(F_0\equiv0\), then for every
\(\mu\in\mathcal D_{N,R,\kappa}\) and every interval \(I\subseteq\Theta\) with
\(\lvert I\rvert>0\),

\begin{equation}
\label{eq:appendix-display-232}
\begin{aligned}
\Pr_{\alpha\sim\mu}
[\exists\theta\in I:\langle\alpha,F(\theta)\rangle=0]
&\leq
A\sqrt{\frac N2}\,
\Gamma_{\mathrm{proj}}(F)\lvert I\rvert\\
&\leq
A\sqrt{\frac N2}\,
\widehat\Lambda_{B,T}\lvert I\rvert.
\end{aligned}
\end{equation}

This is ordinary probability for the arbitrary, possibly correlated, full-joint-density law. The coefficients
and both inequalities are literal, with no clipping at one.


\end{proposition}

\begin{proof}
Fix the complete deterministic instance, including \(F_0\equiv0\), before selecting
a law or interval. Now fix an arbitrary \(\mu\in\mathcal D_{N,R,\kappa}\), and only afterward fix an arbitrary
interval \(I\subseteq\Theta\) with \(\lvert I\rvert>0\), retaining its literal endpoint convention.

Because \(F_0\equiv0\), one has \(F_0'\equiv0\) and
\(\phi_\alpha(\theta)=\langle\alpha,F(\theta)\rangle\). Apply only the coordinate-free first
swept-area inequality, Proposition~\ref{prop:step-004-s2-affine-swept-area}:

\begin{equation}
\label{eq:appendix-display-233}
\begin{aligned}
\Pr_{\alpha\sim\mu}
[\exists\theta\in I:\langle\alpha,F(\theta)\rangle=0]
&\leq
\kappa\int_I\int_{H_\theta\cap[-R,R]^N}
\frac{\lvert\langle a,F'(\theta)\rangle\rvert}
{\lVert F(\theta)\rVert_2}
\,d\mathcal H^{N-1}(a)\,d\theta.
\end{aligned}
\end{equation}

Lemma~\ref{lem:step-008-s2-radial-cancellation} identifies the same actual section as
\(H_\theta=\gamma_F(\theta)^\perp\) and replaces the integrand by equality, not by an estimate:

\begin{equation}
\label{eq:appendix-display-234}
\begin{aligned}
\Pr_{\alpha\sim\mu}
[\exists\theta\in I:\langle\alpha,F(\theta)\rangle=0]
&\leq
\kappa\int_I\int_{\gamma_F(\theta)^\perp\cap[-R,R]^N}
\lvert\langle a,\gamma_F'(\theta)\rangle\rvert
\,d\mathcal H^{N-1}(a)\,d\theta.
\end{aligned}
\end{equation}

For every \(a\in[-R,R]^N\),

\begin{equation}
\label{eq:appendix-display-235}
\lVert a\rVert_2^2
=\sum_{i=1}^N a_i^2
\leq NR^2,
\qquad
\lVert a\rVert_2\leq R\sqrt N.
\end{equation}

Euclidean Cauchy--Schwarz therefore gives, on the actual central section,

\begin{equation}
\label{eq:appendix-display-236}
\lvert\langle a,\gamma_F'(\theta)\rangle\rvert
\leq
R\sqrt N\,\lVert\gamma_F'(\theta)\rVert_2.
\end{equation}

The integrands are nonnegative. Hence Proposition~\ref{prop:step-005-translated-section-certificate}, applied with the actual Euclidean unit normal
\(u=\gamma_F(\theta)\) and \(t=0\), yields

\begin{equation}
\label{eq:appendix-display-237}
\begin{aligned}
&\int_{\gamma_F(\theta)^\perp\cap[-R,R]^N}
\lvert\langle a,\gamma_F'(\theta)\rangle\rvert
\,d\mathcal H^{N-1}(a)\\
&\qquad\leq
R\sqrt N\,\lVert\gamma_F'(\theta)\rVert_2\,
\mathcal H^{N-1}
\bigl(\gamma_F(\theta)^\perp\cap[-R,R]^N\bigr)\\
&\qquad\leq
R\sqrt N\,\sqrt{2}(2R)^{N-1}
\lVert\gamma_F'(\theta)\rVert_2.
\end{aligned}
\end{equation}

This calculation is literal also for \(N=1\): the actual section is the singleton \(\{0\}\),
\(\mathcal H^0\) counts it once, and the integrand at \(a=0\) is zero. Thus no positive point mass is inferred
from the counting measure.

Integrating the pointwise inequality and using the definition of projective speed gives

\begin{equation}
\label{eq:appendix-display-238}
\begin{aligned}
\Pr_{\alpha\sim\mu}
[\exists\theta\in I:\langle\alpha,F(\theta)\rangle=0]
&\leq
\kappa R\sqrt N\,\sqrt{2}(2R)^{N-1}
\int_I\lVert\gamma_F'(\theta)\rVert_2\,d\theta\\
&\leq
\kappa R\sqrt N\,\sqrt{2}(2R)^{N-1}
\Gamma_{\mathrm{proj}}(F)\lvert I\rvert.
\end{aligned}
\end{equation}

No interval-length lower bound other than \(\lvert I\rvert>0\) is used. If
\(\Gamma_{\mathrm{proj}}(F)=0\), Proposition~\ref{prop:step-008-s2-stationary-projective} independently
identifies the event as one fixed proper central hyperplane and proves its probability is zero; the
nonnegative-integral derivation in \eqref{eq:appendix-display-238} also remains valid without dividing by
\(\Gamma_{\mathrm{proj}}(F)\).

Apply the literal identity in Lemma~\ref{lem:step-008-s2-literal-algebra}:

\begin{equation}
\label{eq:appendix-display-239}
\Pr_{\alpha\sim\mu}
[\exists\theta\in I:\langle\alpha,F(\theta)\rangle=0]
\leq
A\sqrt{\frac N2}\,
\Gamma_{\mathrm{proj}}(F)\lvert I\rvert.
\end{equation}

Finally, apply the exact named projective certificate,
Proposition~\ref{prop:step-001-projective},

\begin{equation}
\label{eq:appendix-display-240}
\Gamma_{\mathrm{proj}}(F)\leq\widehat\Lambda_{B,T},
\end{equation}

to obtain

\begin{equation}
\label{eq:appendix-display-241}
\Pr_{\alpha\sim\mu}
[\exists\theta\in I:\langle\alpha,F(\theta)\rangle=0]
\leq
A\sqrt{\frac N2}\,
\widehat\Lambda_{B,T}\lvert I\rvert.
\end{equation}

The full-joint-density factor \(\kappa\) entered only once, inside the first swept-area inequality.
There is no conditioning, marginalization, independence reduction, union bound, second probability conversion,
chart-count factor, or auxiliary tolerance. If either stated right-hand side exceeds one, it remains a valid
upper bound and is not clipped. 
\end{proof}

\begin{proposition}[Defining-supremum closure for the homogeneous capacity]
\label{prop:step-008-s2-pf-closure}

Under Assumptions~\ref{assump:parameter-regime}, \ref{assump:balcan-common-chain},
\ref{assump:cube-density-laws}, and \ref{assump:anchored-derivative-closure},
Proposition~\ref{prop:step-008-s2-pairwise-homogeneous-rate}, and the exact specialization
\(F_0\equiv0\),

\begin{equation}
\label{eq:appendix-display-242}
C^{\mathrm{Pf}}_{\mathcal D}(F;\Theta)
\leq
A\sqrt{\frac N2}\,\widehat\Lambda_{B,T}.
\end{equation}

The interval supremum is taken for each fixed law before the outer law supremum, exactly as in the setting
definition.


\end{proposition}

\begin{proof}
Fix \(\mu\in\mathcal D\). For every interval \(I\subseteq\Theta\) with
\(\lvert I\rvert>0\), Proposition~\ref{prop:step-008-s2-pairwise-homogeneous-rate} gives

\begin{equation}
\label{eq:appendix-display-243}
\Pr_{\alpha\sim\mu}
[\exists\theta\in I:\langle\alpha,F(\theta)\rangle=0]
\leq
A\sqrt{\frac N2}\,
\widehat\Lambda_{B,T}\lvert I\rvert.
\end{equation}

Only now divide by the strictly positive \(\lvert I\rvert\):

\begin{equation}
\label{eq:appendix-display-244}
\frac{
\Pr_{\alpha\sim\mu}
[\exists\theta\in I:\langle\alpha,F(\theta)\rangle=0]
}{\lvert I\rvert}
\leq
A\sqrt{\frac N2}\,\widehat\Lambda_{B,T}.
\end{equation}

The right-hand side is deterministic and independent of \(I\). Therefore, still for this fixed law,

\begin{equation}
\label{eq:appendix-display-245}
\sup_{\substack{I\subseteq\Theta\ \mathrm{interval}\\ \lvert I\rvert>0}}
\frac{
\Pr_{\alpha\sim\mu}
[\exists\theta\in I:\langle\alpha,F(\theta)\rangle=0]
}{\lvert I\rvert}
\leq
A\sqrt{\frac N2}\,\widehat\Lambda_{B,T}.
\end{equation}

The same deterministic inequality holds for every \(\mu\in\mathcal D\). Taking the outer law supremum only
after the interval supremum gives

\begin{equation}
\label{eq:appendix-display-246}
\begin{aligned}
C^{\mathrm{Pf}}_{\mathcal D}(F;\Theta)
&=
\sup_{\mu\in\mathcal D}
\sup_{\substack{I\subseteq\Theta\ \mathrm{interval}\\ \lvert I\rvert>0}}
\frac{
\Pr_{\alpha\sim\mu}
[\exists\theta\in I:\langle\alpha,F(\theta)\rangle=0]
}{\lvert I\rvert}\\
&\leq
A\sqrt{\frac N2}\,\widehat\Lambda_{B,T}.
\end{aligned}
\end{equation}

No zero-length interval is divided by, and arbitrarily short positive-length intervals remain within the
supremum. 
\end{proof}


\begin{proof}[Synthesis of the homogeneous projective rate]
Fix the homogeneous deterministic presentation before selecting a law and
interval.  Lemma~\ref{lem:step-001-anchor} makes
$r=\|F\|_2\geq1$, and Lemma~\ref{lem:step-008-s2-radial-cancellation}
gives on the actual section
\begin{equation}
\label{eq:appendix-display-247}
H_\theta=F(\theta)^\perp=\gamma_F(\theta)^\perp,
\qquad
\frac{|\langle a,F'(\theta)\rangle|}{\|F(\theta)\|_2}
=|\langle a,\gamma_F'(\theta)\rangle|.
\end{equation}
Proposition~\ref{prop:step-008-s2-stationary-projective} handles
$\Gamma_{\mathrm{proj}}(F)=0$ as one fixed proper law-null central
hyperplane. Apply Cauchy--Schwarz and $\|a\|_2\leq R\sqrt N$ in
Proposition~\ref{prop:step-004-s2-affine-swept-area}. The section estimate
is Proposition~\ref{prop:step-005-translated-section-certificate}. Together
they give the coefficient $\kappa R\sqrt N\sqrt2(2R)^{N-1}$, which
Lemma~\ref{lem:step-008-s2-literal-algebra} identifies exactly with
$A\sqrt{N/2}$.  Proposition~\ref{prop:step-008-s2-pairwise-homogeneous-rate}
therefore gives
\begin{equation}
\label{eq:appendix-display-248}
\Pr_{\alpha\sim\mu}
[\exists\theta\in I:\langle\alpha,F(\theta)\rangle=0]
\leq A\sqrt{\frac N2}\Gamma_{\mathrm{proj}}(F)|I|
\leq A\sqrt{\frac N2}\widehat\Lambda_{B,T}|I|,
\end{equation}
where the last inequality is Proposition~\ref{prop:step-001-projective}.
Finally, Proposition~\ref{prop:step-008-s2-pf-closure} preserves the
interval-then-law supremum order and proves
\begin{equation}
\label{eq:appendix-display-249}
C^{\mathrm{Pf}}_{\mathcal D}(F;\Theta)
\leq A\sqrt{\frac N2}\widehat\Lambda_{B,T}.
\end{equation}
\end{proof}

\subsection{Exact Monic Presentation}\label{app:step-009}

\begin{proposition}[Augmented Monomial Presentation] \label{prop:step-009-monic-presentation}
Under Assumptions~\ref{assump:parameter-regime}, \ref{assump:balcan-common-chain}, and the explicit monic specialization of Assumption~\ref{assump:anchored-derivative-closure}, let \(d\geq1\), let \(J\subset\mathbb R\) be a bounded interval, and let \(\alpha=(\alpha_0,\ldots,\alpha_{d-1})\in[-R,R]^d\). If
\begin{equation}
\label{eq:appendix-display-250}
F_0(\theta)=\theta^d,
\qquad
F_{k+1}(\theta)=\theta^k\quad(0\leq k\leq d-1),
\end{equation}
then, with the ordered augmented tuple
\begin{equation}
\label{eq:appendix-display-251}
\widetilde F=(F_0,F_1,\ldots,F_d)
=(\theta^d,1,\theta,\ldots,\theta^{d-1}),
\end{equation}
one has pointwise on \(\mathbb R\)
\begin{equation}
\label{eq:appendix-display-252}
F_0(\theta)+\langle\alpha,F(\theta)\rangle
=\theta^d+\sum_{k=0}^{d-1}\alpha_k\theta^k
=p_\alpha(\theta).
\end{equation}
The Balcan descriptor values are exactly \(q=0\), \(M=0\), \(\Delta=d\), \(N=d\), and \(A=(2R)^d\kappa\), and the literal anchor is \(F_1=1\). The coefficient vector contains exactly the \(d\) lower coefficients; the leading coefficient one is deterministic and is not a coordinate of \(\alpha\).


\end{proposition}

\begin{proof}
Take the empty common chain. By the setting's explicit convention for an empty chain,
\begin{equation}
\label{eq:appendix-display-253}
q=0,
\qquad
M=0.
\end{equation}
The output polynomials in the one scalar parameter are
\begin{equation}
\label{eq:appendix-display-254}
Q_0(\theta)=\theta^d,
\qquad
Q_{k+1}(\theta)=\theta^k
\quad(0\leq k\leq d-1).
\end{equation}
Their coordinate index map is
\begin{equation}
\label{eq:appendix-display-255}
k\in\{0,\ldots,d-1\}
\longmapsto
i=k+1\in\{1,\ldots,d\}.
\end{equation}
Thus there are \(d\) random-feature coordinates, so \(N=d\), and
\begin{equation}
\label{eq:appendix-display-256}
\Delta
=\max\left\{\deg Q_0,\max_{0\leq k\leq d-1}\deg Q_{k+1}\right\}
=\max\{d,d-1\}
=d.
\end{equation}
This calculation also covers \(d=1\), when the second maximum is the single value \(0\). By the setting definition of \(A\),
\begin{equation}
\label{eq:appendix-display-257}
A=(2R)^N\kappa=(2R)^d\kappa.
\end{equation}
The first random feature corresponds to \(k=0\), hence
\begin{equation}
\label{eq:appendix-display-258}
F_1(\theta)=\theta^0=1
\quad\text{for every }\theta,
\end{equation}
which verifies the literal anchor required by Assumption~\ref{assump:anchored-derivative-closure}.

Finally, using the index map from Proposition~\ref{prop:step-009-monic-presentation} coefficient by coefficient,
\begin{equation}
\label{eq:appendix-display-259}
\begin{aligned}
F_0(\theta)+\langle\alpha,F(\theta)\rangle
&=F_0(\theta)+\sum_{i=1}^{d}\alpha_{i-1}F_i(\theta)\\
&=\theta^d+\sum_{i=1}^{d}\alpha_{i-1}\theta^{i-1}\\
&=\theta^d+\sum_{k=0}^{d-1}\alpha_k\theta^k
=p_\alpha(\theta).
\end{aligned}
\end{equation}
This is a pointwise identity for every real \(\theta\). The domain \([-R,R]^d\) is the actual lower-coefficient cube. There is no coordinate \(\alpha_d\), and the deterministic term \(F_0=\theta^d\) is not part of the random feature inner product.
\end{proof}

\begin{lemma}[Exact Derivative-Shift Certificate] \label{lem:step-009-derivative-shift}
Under Assumptions~\ref{assump:parameter-regime}, \ref{assump:balcan-common-chain}, and \ref{assump:anchored-derivative-closure}, let \(d\geq1\) and let \(\widetilde F\) be the ordered tuple in Proposition~\ref{prop:step-009-monic-presentation}. Define a constant matrix \(B\in\mathbb R^{(d+1)\times(d+1)}\), with rows and columns indexed by \(0,1,\ldots,d\), by declaring its only nonzero entries to be
\begin{equation}
\label{eq:appendix-display-260}
B_{0,d}=d,
\qquad
B_{k+1,k}=k\quad(1\leq k\leq d-1).
\end{equation}
Then \(\widetilde F'=B\widetilde F\) coefficient by coefficient, \(m=0\), and
\begin{equation}
\label{eq:appendix-display-261}
\widehat\Lambda_{B,T}
=\left(\sum_{k=1}^{d}k^2\right)^{1/2}.
\end{equation}
For \(d=1\), \(B\) is the \(2\)-by-\(2\) matrix with the sole nonzero entry \(B_{0,1}=1\).


\end{lemma}

\begin{proof}
The row actions of \(B\) follow the ordered coordinate map exactly.

Row \(0\) has only column \(d\) nonzero. Since
\(F_d(\theta)=\theta^{d-1}\),
\begin{equation}
\label{eq:appendix-indented-005}
  (B\widetilde F)_0=dF_d=d\theta^{d-1}=F_0'(\theta).
\end{equation}
Row \(1\) is zero because it is not among the declared nonzero rows. Hence
\begin{equation}
\label{eq:appendix-indented-006}
  (B\widetilde F)_1=0=(1)'=F_1'(\theta).
\end{equation}
Every remaining row is uniquely \(r=k+1\) for an integer
\(1\leq k\leq d-1\). Its only nonzero column is \(s=k\), and the feature
index map gives \(F_k(\theta)=\theta^{k-1}\). Therefore
\begin{equation}
\label{eq:appendix-indented-007}
  (B\widetilde F)_{k+1}
  =kF_k
  =k\theta^{k-1}
  =(\theta^k)'
  =F_{k+1}'(\theta).
\end{equation}

Rows \(0\), \(1\), and \(k+1\) for \(1\leq k\leq d-1\) are exactly all rows \(0,1,\ldots,d\). Thus every coordinate agrees and
\begin{equation}
\label{eq:appendix-display-262}
\widetilde F'(\theta)=B\widetilde F(\theta)
\quad\text{for every }\theta\in\mathbb R.
\end{equation}
This verifies the closure identity of Assumption~\ref{assump:anchored-derivative-closure}; it was not imported as an unverified property of the monic tuple.

Because every entry of \(B\) is constant, its setting representation is
\begin{equation}
\label{eq:appendix-display-263}
B_{rs}(\theta)=b_{rs,0},
\qquad
b_{rs,0}=B_{rs},
\end{equation}
with no coefficient of positive degree. Hence the polynomial matrix degree is exactly \(m=0\). Since \(T_*^0=1\), the setting-defined coefficient height satisfies
\begin{equation}
\label{eq:appendix-display-264}
\begin{aligned}
\widehat\Lambda_{B,T}^{2}
&=\sum_{r=0}^{d}\sum_{s=0}^{d}
  \left(\sum_{\ell=0}^{0}\lvert b_{rs,\ell}\rvert T_*^\ell\right)^2\\
&=\sum_{r=0}^{d}\sum_{s=0}^{d}\lvert B_{rs}\rvert^2\\
&=\lVert B\rVert_{\mathrm F}^{2}\\
&=d^2+\sum_{k=1}^{d-1}k^2\\
&=\sum_{k=1}^{d}k^2.
\end{aligned}
\end{equation}
Taking the nonnegative square root proves the asserted certificate. In particular, neither \(T\) nor the location of \(J\) appears.

When \(d=1\), the range \(1\leq k\leq d-1\) is empty and
\begin{equation}
\label{eq:appendix-display-265}
B=\begin{pmatrix}0&1\\0&0\end{pmatrix}.
\end{equation}
Its row actions are \((B\widetilde F)_0=F_1=1=(\theta)'\) and \((B\widetilde F)_1=0=(1)'\), while the height calculation gives \(\widehat\Lambda_{B,T}^2=1=\sum_{k=1}^{1}k^2\).
\end{proof}

\begin{proposition}[Two-Pivot Coefficient Charts] \label{prop:step-009-two-pivot-charts}
Under Assumptions~\ref{assump:parameter-regime}, \ref{assump:balcan-common-chain}, and \ref{assump:anchored-derivative-closure}, let \(d\geq2\), let \(J\subset\mathbb R\) be a bounded interval, and use the tuple and coefficient vector of Proposition~\ref{prop:step-009-monic-presentation}. Then
\begin{equation}
\label{eq:appendix-display-266}
E_1=J\cap\{\lvert\theta\rvert\leq1\},
\qquad
E_d=J\cap\{\lvert\theta\rvert>1\}
\end{equation}
is a measurable disjoint partition of \(J\), with \(\theta=\pm1\) assigned to \(E_1\). The pivot \(\alpha_0\) is legal on \(E_1\), the pivot \(\alpha_{d-1}\) is legal on \(E_d\), and the exact pivot-coordinate maps are
\begin{equation}
\label{eq:appendix-display-267}
T_1(\theta,\beta)
=-\theta^d-\sum_{k=1}^{d-1}\beta_k\theta^k,
\qquad
\beta=(\beta_1,\ldots,\beta_{d-1})\in[-R,R]^{d-1},
\end{equation}
and
\begin{equation}
\label{eq:appendix-display-268}
T_d(\theta,\beta)
=-\theta-\sum_{k=0}^{d-2}\beta_k\theta^{k-d+1},
\qquad
\beta=(\beta_0,\ldots,\beta_{d-2})\in[-R,R]^{d-1}.
\end{equation}
Inserting the pivot values into their respective coefficient coordinates gives a pointwise zero of the same polynomial \(p_\alpha\).


\end{proposition}

\begin{proof}
Every interval is Borel. The set \(\{\lvert\theta\rvert\leq1\}\) is closed and the set \(\{\lvert\theta\rvert>1\}\) is open, so \(E_1\) and \(E_d\) are measurable. Their defining inequalities are complementary and disjoint, hence
\begin{equation}
\label{eq:appendix-display-269}
E_1\cap E_d=\varnothing,
\qquad
E_1\cup E_d=J.
\end{equation}
Both boundary points \(1\) and \(-1\) satisfy the weak inner inequality and therefore belong to \(E_1\) whenever they belong to \(J\).

The coefficient-to-feature index map of Proposition~\ref{prop:step-009-monic-presentation} sends \(\alpha_0\) to \(F_1=1\). Thus the inner pivot denominator is exactly one on all of \(E_1\). Solving
\begin{equation}
\label{eq:appendix-display-270}
0=\theta^d+\alpha_0+\sum_{k=1}^{d-1}\alpha_k\theta^k
\end{equation}
for \(\alpha_0\), and writing the remaining coefficient \(\alpha_k\) as \(\beta_k\), gives
\begin{equation}
\label{eq:appendix-display-271}
T_1(\theta,\beta)
=-\theta^d-\sum_{k=1}^{d-1}\beta_k\theta^k.
\end{equation}
The full coefficient insertion map is
\begin{equation}
\label{eq:appendix-display-272}
\Psi_1(\theta,\beta)
=\bigl(T_1(\theta,\beta),\beta_1,\ldots,\beta_{d-1}\bigr)
\in\mathbb R^d,
\end{equation}
and direct substitution verifies
\begin{equation}
\label{eq:appendix-display-273}
p_{\Psi_1(\theta,\beta)}(\theta)
=\theta^d+T_1(\theta,\beta)
+\sum_{k=1}^{d-1}\beta_k\theta^k
=0.
\end{equation}

On \(E_d\), the outer pivot \(\alpha_{d-1}\) multiplies
\begin{equation}
\label{eq:appendix-display-274}
F_d(\theta)=\theta^{d-1}.
\end{equation}
Because \(d\geq2\) and \(\lvert\theta\rvert>1\),
\(\lvert F_d(\theta)\rvert=\lvert\theta\rvert^{d-1}>1\), so this pivot is nonzero. Solving
\begin{equation}
\label{eq:appendix-display-275}
0=\theta^d+\sum_{k=0}^{d-2}\alpha_k\theta^k
+\alpha_{d-1}\theta^{d-1}
\end{equation}
for \(\alpha_{d-1}\), with \(\beta_k=\alpha_k\) for \(0\leq k\leq d-2\), gives
\begin{equation}
\label{eq:appendix-display-276}
\begin{aligned}
T_d(\theta,\beta)
&=-\frac{\theta^d}{\theta^{d-1}}
  -\sum_{k=0}^{d-2}\beta_k\frac{\theta^k}{\theta^{d-1}}\\
&=-\theta-\sum_{k=0}^{d-2}\beta_k\theta^{k-d+1}.
\end{aligned}
\end{equation}
All divisions are legal because \(\theta\neq0\) on \(E_d\). The full insertion map is
\begin{equation}
\label{eq:appendix-display-277}
\Psi_d(\theta,\beta)
=\bigl(\beta_0,\ldots,\beta_{d-2},T_d(\theta,\beta)\bigr)
\in\mathbb R^d.
\end{equation}
Multiplying the pivot formula by \(\theta^{d-1}\) yields
\begin{equation}
\label{eq:appendix-display-278}
T_d(\theta,\beta)\theta^{d-1}
=-\theta^d-\sum_{k=0}^{d-2}\beta_k\theta^k,
\end{equation}
and consequently
\begin{equation}
\label{eq:appendix-display-279}
p_{\Psi_d(\theta,\beta)}(\theta)
=\theta^d+\sum_{k=0}^{d-2}\beta_k\theta^k
+T_d(\theta,\beta)\theta^{d-1}
=0.
\end{equation}
The nonpivot input is the actual cube \([-R,R]^{d-1}\). The inserted pivot value may or may not lie in \([-R,R]\); the downstream chart indicator, not this deterministic specialization, handles that separate condition.
\end{proof}

\begin{lemma}[Inner-Region Chart Velocity] \label{lem:step-009-inner-velocity}
Under Assumptions~\ref{assump:parameter-regime}, \ref{assump:balcan-common-chain}, and \ref{assump:anchored-derivative-closure}, and Proposition~\ref{prop:step-009-two-pivot-charts}, let \(d\geq2\), \(\theta\in E_1\), and \(\beta=(\beta_1,\ldots,\beta_{d-1})\in[-R,R]^{d-1}\). Then
\begin{equation}
\label{eq:appendix-display-280}
\lvert\partial_\theta T_1(\theta,\beta)\rvert
\leq d+R\sum_{k=1}^{d-1}k
=d+\frac{Rd(d-1)}2.
\end{equation}


\end{lemma}

\begin{proof}
Term-by-term differentiation of the finite inner chart gives
\begin{equation}
\label{eq:appendix-display-281}
\partial_\theta T_1(\theta,\beta)
=-d\theta^{d-1}-\sum_{k=1}^{d-1}k\beta_k\theta^{k-1}.
\end{equation}
Here every exponent \(d-1\) and \(k-1\) is nonnegative. Because \(\theta\in E_1\) implies \(\lvert\theta\rvert\leq1\), and because \(\lvert\beta_k\rvert\leq R\), the triangle inequality gives every intermediate term explicitly:
\begin{equation}
\label{eq:appendix-display-282}
\begin{aligned}
\lvert\partial_\theta T_1(\theta,\beta)\rvert
&\leq d\lvert\theta\rvert^{d-1}
  +\sum_{k=1}^{d-1}k\lvert\beta_k\rvert\lvert\theta\rvert^{k-1}\\
&\leq d
  +R\sum_{k=1}^{d-1}k.
\end{aligned}
\end{equation}
At \(\theta=0\), the \(k=1\) factor is \(\lvert\theta\rvert^0=1\), while every term with a positive exponent vanishes, so the same calculation remains literal.

For completeness, the finite arithmetic sum follows by pairing the list \(1,\ldots,d-1\) with its reverse. Since \(k\mapsto d-k\) permutes \(\{1,\ldots,d-1\}\),
\begin{equation}
\label{eq:appendix-display-283}
\begin{aligned}
2\sum_{k=1}^{d-1}k
&=\sum_{k=1}^{d-1}k+\sum_{k=1}^{d-1}(d-k)\\
&=\sum_{k=1}^{d-1}\bigl(k+d-k\bigr)\\
&=\sum_{k=1}^{d-1}d\\
&=d(d-1).
\end{aligned}
\end{equation}
Therefore \(\sum_{k=1}^{d-1}k=d(d-1)/2\), which proves
\begin{equation}
\label{eq:appendix-display-284}
\lvert\partial_\theta T_1\rvert
\leq d+\frac{Rd(d-1)}2.
\end{equation}
\end{proof}

\begin{lemma}[Outer-Region Chart Velocity] \label{lem:step-009-outer-velocity}
Under Assumptions~\ref{assump:parameter-regime}, \ref{assump:balcan-common-chain}, and \ref{assump:anchored-derivative-closure}, and Proposition~\ref{prop:step-009-two-pivot-charts}, let \(d\geq2\), \(\theta\in E_d\), and \(\beta=(\beta_0,\ldots,\beta_{d-2})\in[-R,R]^{d-1}\). Then
\begin{equation}
\label{eq:appendix-display-285}
\lvert\partial_\theta T_d(\theta,\beta)\rvert
\leq1+R\sum_{m=1}^{d-1}\frac{m}{\lvert\theta\rvert^{m+1}}
\leq1+\frac{Rd(d-1)}2
\leq d+\frac{Rd(d-1)}2.
\end{equation}


\end{lemma}

\begin{proof}
Every power in the outer chart is defined because \(\lvert\theta\rvert>1\). Direct differentiation gives
\begin{equation}
\label{eq:appendix-display-286}
\begin{aligned}
\partial_\theta T_d(\theta,\beta)
&=-1-\sum_{k=0}^{d-2}\beta_k(k-d+1)\theta^{k-d}\\
&=-1+\sum_{k=0}^{d-2}(d-1-k)\beta_k\theta^{k-d}.
\end{aligned}
\end{equation}
Using \(\lvert\beta_k\rvert\leq R\) and the triangle inequality,
\begin{equation}
\label{eq:appendix-display-287}
\lvert\partial_\theta T_d(\theta,\beta)\rvert
\leq1+R\sum_{k=0}^{d-2}(d-1-k)\lvert\theta\rvert^{k-d}.
\end{equation}
Now set
\begin{equation}
\label{eq:appendix-display-288}
m=d-1-k.
\end{equation}
As \(k\) increases from \(0\) to \(d-2\), \(m\) decreases from \(d-1\) to \(1\), and
\begin{equation}
\label{eq:appendix-display-289}
k=d-1-m,
\qquad
k-d=-m-1.
\end{equation}
Reversing this finite order gives the exact reindexing
\begin{equation}
\label{eq:appendix-display-290}
\begin{aligned}
\sum_{k=0}^{d-2}(d-1-k)\lvert\theta\rvert^{k-d}
&=\sum_{m=1}^{d-1}m\lvert\theta\rvert^{-m-1}\\
&=\sum_{m=1}^{d-1}\frac{m}{\lvert\theta\rvert^{m+1}}.
\end{aligned}
\end{equation}
This is the only use of negative powers. Since \(\lvert\theta\rvert>1\), every \(m\geq1\) satisfies
\begin{equation}
\label{eq:appendix-display-291}
0<\frac1{\lvert\theta\rvert^{m+1}}<1.
\end{equation}
Therefore
\begin{equation}
\label{eq:appendix-display-292}
\lvert\partial_\theta T_d(\theta,\beta)\rvert
\leq1+R\sum_{m=1}^{d-1}m.
\end{equation}
Again displaying the entire finite-sum calculation, the reverse-list pairing gives
\begin{equation}
\label{eq:appendix-display-293}
2\sum_{m=1}^{d-1}m
=\sum_{m=1}^{d-1}\bigl(m+d-m\bigr)
=\sum_{m=1}^{d-1}d
=d(d-1),
\end{equation}
so
\begin{equation}
\label{eq:appendix-display-294}
\lvert\partial_\theta T_d(\theta,\beta)\rvert
\leq1+\frac{Rd(d-1)}2.
\end{equation}
Finally, \(d\geq2\) implies \(1\leq d\), and hence
\begin{equation}
\label{eq:appendix-display-295}
1+\frac{Rd(d-1)}2
\leq d+\frac{Rd(d-1)}2.
\end{equation}
For negative outer \(\theta\), integer powers may change sign inside the derivative, but
\(\lvert\theta^{k-d}\rvert=\lvert\theta\rvert^{k-d}\); this absolute-value calculation is unchanged. No upper bound on \(\lvert\theta\rvert\) was used.
\end{proof}

\begin{proposition}[Global Linear-Case Chart] \label{prop:step-009-linear-case}
Under Assumptions~\ref{assump:parameter-regime}, \ref{assump:balcan-common-chain}, and \ref{assump:anchored-derivative-closure}, and Proposition~\ref{prop:step-009-monic-presentation} specialized to \(d=1\), let \(J\subset\mathbb R\) be bounded and let \(\alpha_0\in[-R,R]\). Then, with the explicit zero-dimensional convention
\begin{equation}
\label{eq:appendix-display-296}
[-R,R]^0=\{()\},
\end{equation}
\begin{equation}
\label{eq:appendix-display-297}
\widetilde F(\theta)=(\theta,1),
\qquad
p_{\alpha_0}(\theta)=\theta+\alpha_0,
\end{equation}
the sole coefficient \(\alpha_0\) is a legal pivot on every \(\theta\in J\), the zero-dimensional-beta chart is
\begin{equation}
\label{eq:appendix-display-298}
T_1(\theta)=-\theta,
\end{equation}
and
\begin{equation}
\label{eq:appendix-display-299}
\lvert\partial_\theta T_1(\theta)\rvert=1
=d+\frac{Rd(d-1)}2.
\end{equation}


\end{proposition}

\begin{proof}
For \(d=1\), Proposition~\ref{prop:step-009-monic-presentation} gives the augmented tuple \((F_0,F_1)=(\theta,1)\), the actual coefficient vector \((\alpha_0)\in[-R,R]^1\), and
\begin{equation}
\label{eq:appendix-display-300}
p_{\alpha_0}(\theta)=\theta+\alpha_0.
\end{equation}
The only random feature is the constant \(F_1=1\), so it is nonzero on all of \(J\). Solving the root equation gives \(\alpha_0=-\theta\).

The nonpivot coordinate set has dimension \(d-1=0\). Under the standard product convention,
\begin{equation}
\label{eq:appendix-display-301}
[-R,R]^0=\{()\},
\end{equation}
so there is one empty beta tuple and no beta sum. Thus the full coefficient chart is simply
\begin{equation}
\label{eq:appendix-display-302}
\Psi_1(\theta)=\bigl(T_1(\theta)\bigr)=(-\theta),
\end{equation}
and pointwise
\begin{equation}
\label{eq:appendix-display-303}
p_{\Psi_1(\theta)}(\theta)=\theta-\theta=0.
\end{equation}
Direct differentiation gives \(\partial_\theta T_1=-1\) and hence \(\lvert\partial_\theta T_1\rvert=1\). Since \(d(d-1)/2=0\) at \(d=1\), this is exactly the common cap \(d+Rd(d-1)/2=1\). Lemma~\ref{lem:step-009-derivative-shift} separately verifies that the matrix is \(2\)-by-\(2\), has only \(B_{0,1}=1\), and has height one.
\end{proof}

\begin{lemma}[Location-Free Boundary and Index Closure] \label{lem:step-009-boundary-index-closure}
Under Assumptions~\ref{assump:parameter-regime}, \ref{assump:balcan-common-chain}, and \ref{assump:anchored-derivative-closure}, let \(d\geq1\), let \(J\subset\mathbb R\) be any bounded interval, and use Proposition~\ref{prop:step-009-monic-presentation}, Lemma~\ref{lem:step-009-derivative-shift}, Proposition~\ref{prop:step-009-two-pivot-charts} when \(d\geq2\), Lemmas~\ref{lem:step-009-inner-velocity} and \ref{lem:step-009-outer-velocity} when \(d\geq2\), and Proposition~\ref{prop:step-009-linear-case} when \(d=1\). Then all chart, certificate, and object conclusions remain valid at \(\theta=0\), \(\theta=1\), \(\theta=-1\), negative outer \(\theta\), empty partition pieces, \(d=1\), and \(d=2\). They hold for every location of \(J\), use exactly \(d\) lower-coordinate inputs and a \((d+1)\)-square augmented matrix, and contain neither a singular/randomized leading coordinate nor a bound depending on \(\sup_{\theta\in J}\lvert\theta\rvert\).


\end{lemma}

\begin{proof}
At \(\theta=0\), the point is in \(E_1\) for \(d\geq2\), and the pivot is the literal constant \(F_1=1\). The inner formula contains only nonnegative powers, and Lemma~\ref{lem:step-009-inner-velocity} explicitly treats the exponent-zero term. At \(\theta=1\) and \(\theta=-1\), the weak inequality \(\lvert\theta\rvert\leq1\) again assigns the point to \(E_1\), so no negative power is evaluated at the transition.

For negative \(\theta\) with \(\lvert\theta\rvert>1\), the outer pivot \(\theta^{d-1}\) is nonzero. Proposition~\ref{prop:step-009-two-pivot-charts} derives the chart algebraically at that same signed \(\theta\), while Lemma~\ref{lem:step-009-outer-velocity} takes absolute values only after differentiation and therefore retains all signs correctly.

Either \(E_1\) or \(E_d\) may be empty. The set identities
\begin{equation}
\label{eq:appendix-display-304}
E_1\cap E_d=\varnothing,
\qquad
E_1\cup E_d=J
\end{equation}
remain valid, and a chart assertion on an empty cell is vacuous. If \(J\) itself is empty, both cells are empty and every pointwise assertion is vacuous. If \(J\) lies wholly in one region, only that region's chart is active.

The case \(d=1\) is closed by Proposition~\ref{prop:step-009-linear-case}. For the first genuinely two-pivot case \(d=2\), every index can be inspected explicitly:
\begin{equation}
\label{eq:appendix-display-305}
\widetilde F=(\theta^2,1,\theta),
\qquad
B=\begin{pmatrix}
0&0&2\\
0&0&0\\
0&1&0
\end{pmatrix},
\qquad
\widehat\Lambda_{B,T}=\sqrt{2^2+1^2}=\sqrt5.
\end{equation}
The two charts and their derivatives are
\begin{equation}
\label{eq:appendix-display-306}
T_1(\theta,\beta_1)=-\theta^2-\beta_1\theta,
\qquad
\partial_\theta T_1=-2\theta-\beta_1,
\end{equation}
and
\begin{equation}
\label{eq:appendix-display-307}
T_2(\theta,\beta_0)=-\theta-\beta_0\theta^{-1},
\qquad
\partial_\theta T_2=-1+\beta_0\theta^{-2}.
\end{equation}
Consequently, on their respective regions,
\begin{equation}
\label{eq:appendix-display-308}
\lvert\partial_\theta T_1\rvert\leq2+R,
\qquad
\lvert\partial_\theta T_2\rvert
\leq1+\frac{R}{\lvert\theta\rvert^2}
\leq1+R
\leq2+R.
\end{equation}
Thus the index conventions and bounds agree at \(d=2\) without a missing row, feature, or coefficient.

For any bounded interval \(J\), Assumption~\ref{assump:parameter-regime} can be instantiated on a compact \(\Theta\supseteq J\) inside some \([-T,T]\). Lemma~\ref{lem:step-009-derivative-shift} has \(m=0\), so \(T_*^0=1\) eliminates \(T\) from the certificate. The inner velocity uses only \(\lvert\theta\rvert\leq1\), and the outer velocity uses reciprocal powers bounded by one on \(\lvert\theta\rvert>1\). Hence the bounds do not use \(\sup_{\theta\in J}\lvert\theta\rvert\), even when the location of \(J\) varies without a common bound across problem instances.

Finally, Proposition~\ref{prop:step-009-monic-presentation} maps coefficient indices \(k=0,\ldots,d-1\) bijectively to feature indices \(1,\ldots,d\), and Lemma~\ref{lem:step-009-derivative-shift} indexes the augmented matrix by \(0,\ldots,d\). Thus the coefficient cube is exactly \([-R,R]^d\), each beta cube has dimension \(d-1\), and the matrix has size \((d+1)\)-by-\((d+1)\). The leading coefficient one belongs to deterministic \(F_0=\theta^d\); no extra random coordinate, augmented density, or singular law is introduced.

\end{proof}


\begin{proof}[Synthesis of the exact monic presentation]
For every $d\geq1$, Proposition~\ref{prop:step-009-monic-presentation}
gives
\begin{equation}
\label{eq:appendix-display-309}
\widetilde F=(\theta^d,1,\theta,\ldots,\theta^{d-1}),
\qquad
F_0+\langle\alpha,F\rangle=p_\alpha,
\end{equation}
with $q=M=0$, $\Delta=N=d$, $A=(2R)^d\kappa$, $F_1=1$, and
exactly $\alpha=(\alpha_0,\ldots,\alpha_{d-1})\in[-R,R]^d$ random.
Lemma~\ref{lem:step-009-derivative-shift} gives $m=0$, proves
$\widetilde F'=B\widetilde F$ coefficient by coefficient, and yields
\begin{equation}
\label{eq:appendix-display-310}
B_{0,d}=d,
\qquad B_{k+1,k}=k\quad(1\leq k\leq d-1),
\qquad
\widehat\Lambda_{B,T}=\left(\sum_{k=1}^d k^2\right)^{1/2}.
\end{equation}
For $d\geq2$, Proposition~\ref{prop:step-009-two-pivot-charts} gives the
measurable disjoint inner and outer cells, their legal pivots, and
\begin{equation}
\label{eq:appendix-display-311}
T_1=-\theta^d-\sum_{k=1}^{d-1}\beta_k\theta^k,
\qquad
T_d=-\theta-\sum_{k=0}^{d-2}\beta_k\theta^{k-d+1}.
\end{equation}
Lemmas~\ref{lem:step-009-inner-velocity} and
\ref{lem:step-009-outer-velocity} give
\begin{equation}
\label{eq:appendix-display-312}
|\partial_\theta T_1|
\leq d+\frac{Rd(d-1)}2
\end{equation}
and
\begin{equation}
\label{eq:appendix-display-313}
|\partial_\theta T_d|
\leq1+R\sum_{r=1}^{d-1}\frac{r}{|\theta|^{r+1}}
\leq1+\frac{Rd(d-1)}2
\leq d+\frac{Rd(d-1)}2.
\end{equation}
Proposition~\ref{prop:step-009-linear-case} gives the same cap for $d=1$,
and Lemma~\ref{lem:step-009-boundary-index-closure} covers both signs,
$|\theta|=1$, empty cells, $d=2$, arbitrary bounded interval location, and
all index conventions.  The leading coefficient one remains deterministic;
no additional random coordinate or density is introduced.
\end{proof}

\subsection{Affine-Monic Probability Recovery}\label{app:step-010}

\begin{lemma}[Exact monic chart measure accounting]
\label{lem:step-010-measure-accounting}
Under Proposition~\ref{prop:step-009-two-pivot-charts} when \(d\geq2\) and Proposition~\ref{prop:step-009-linear-case} when \(d=1\), if \(d\geq1\), \(R>0\), and \(J\subset\mathbb R\) is a bounded interval, then every chart's nonpivot cube has exact volume
\begin{equation}
\label{eq:appendix-display-314}
\lambda_{d-1}([-R,R]^{d-1})=(2R)^{d-1},
\end{equation}
where
\begin{equation}
\label{eq:appendix-display-315}
\lambda_0([-R,R]^0)=1=(2R)^0
\end{equation}
when \(d=1\). Moreover, the active cells are measurable and disjoint and have total Lebesgue length \(\lvert J\rvert\). For \(d\geq2\),
\begin{equation}
\label{eq:appendix-display-316}
\lvert E_1\rvert+\lvert E_d\rvert=\lvert J\rvert,
\end{equation}
and the intermediate cells \(E_j=\varnothing\), \(2\leq j\leq d-1\), have zero length; for \(d=1\), \(E_1=J\) and \(\lvert E_1\rvert=\lvert J\rvert\).


\end{lemma}

\begin{proof}
For \(d\geq2\), either chart has exactly \(d-1\) nonpivot coordinates, each ranging over \([-R,R]\). Product measure gives
\begin{equation}
\label{eq:appendix-display-317}
\lambda_{d-1}([-R,R]^{d-1})
=\prod_{r=1}^{d-1}\lambda_1([-R,R])
=\prod_{r=1}^{d-1}2R
=(2R)^{d-1}.
\end{equation}

Proposition~\ref{prop:step-009-two-pivot-charts} gives
\begin{equation}
\label{eq:appendix-display-318}
E_1=J\cap\{\lvert\theta\rvert\leq1\},
\qquad
E_d=J\cap\{\lvert\theta\rvert>1\}.
\end{equation}
The cells are measurable, disjoint, and have union \(J\); \(\theta=\pm1\) belongs only to \(E_1\). Adding the empty intermediate cells yields a legal partition indexed by every \(j=1,\ldots,d\), and finite additivity gives
\begin{equation}
\label{eq:appendix-display-319}
\sum_{j=1}^{d}\lvert E_j\rvert
=\lvert E_1\rvert+\lvert E_d\rvert
=\lvert J\rvert.
\end{equation}
This remains exact if \(J\), \(E_1\), or \(E_d\) is empty, or if \(J\) lies entirely in one region.

For \(d=1\), Proposition~\ref{prop:step-009-linear-case} gives one empty beta tuple and the sole cell \(E_1=J\). By the zero-dimensional product convention,
\begin{equation}
\label{eq:appendix-display-320}
[-R,R]^0=\{()\},
\qquad
\lambda_0(\{()\})=1=(2R)^0,
\end{equation}
and \(\lvert E_1\rvert=\lvert J\rvert\). Thus no dimension introduces a multiplier for the number of active charts. 
\end{proof}

\begin{proposition}[Exact positive-length affine-monic sweep]
\label{prop:step-010-positive-length}
Under Assumption~\ref{assump:cube-density-laws}, Proposition~\ref{prop:step-003-pivot-sweep}, Proposition~\ref{prop:step-009-monic-presentation}, Lemma~\ref{lem:step-009-derivative-shift}, Proposition~\ref{prop:step-009-two-pivot-charts}, Lemmas~\ref{lem:step-009-inner-velocity} and \ref{lem:step-009-outer-velocity}, Proposition~\ref{prop:step-009-linear-case}, Lemma~\ref{lem:step-009-boundary-index-closure}, and Lemma~\ref{lem:step-010-measure-accounting}, let \(d\geq1\), let \(J\subset\mathbb R\) be a bounded interval with \(\lvert J\rvert>0\), and let
\begin{equation}
\label{eq:appendix-display-321}
\alpha=(\alpha_0,\ldots,\alpha_{d-1})\sim\mu,
\end{equation}
where \(\mu\) is any Borel full joint-density law supported on \([-R,R]^d\) with \(f_\mu\leq\kappa\) almost everywhere. Then, in ordinary probability,
\begin{equation}
\label{eq:appendix-display-322}
\Pr_{\alpha\sim\mu}[\exists\theta\in J:p_\alpha(\theta)=0]
\leq
\kappa(2R)^{d-1}
\left(d+\frac{Rd(d-1)}2\right)\lvert J\rvert.
\end{equation}


\end{proposition}

\begin{proof}
Proposition~\ref{prop:step-009-monic-presentation} gives \(N=d\), the original coefficient order \((\alpha_0,\ldots,\alpha_{d-1})\), and the exact pointwise identity
\begin{equation}
\label{eq:appendix-display-323}
F_0(\theta)+\langle\alpha,F(\theta)\rangle
=p_\alpha(\theta).
\end{equation}
Thus the event in Proposition~\ref{prop:step-003-pivot-sweep} is exactly the event in this proposition. The leading coefficient remains the deterministic coefficient of \(F_0=\theta^d\).

Lemma~\ref{lem:step-009-boundary-index-closure} permits a nondegenerate compact \(\Theta\supseteq J\) inside some \([-T,T]\). Proposition~\ref{prop:step-009-monic-presentation} and Lemma~\ref{lem:step-009-derivative-shift} supply on that containing interval the exact presentation, anchor, and derivative closure required by Proposition~\ref{prop:step-003-pivot-sweep}. Because the matrix is constant and \(m=0\), this instantiation introduces no \(T\)-dependent conclusion.

Assumption~\ref{assump:cube-density-laws} supplies one full joint density \(f_\mu\leq\kappa\) on the original cube. Proposition~\ref{prop:step-003-pivot-sweep} already performs direct joint-density domination. This specialization neither factors \(f_\mu\) nor takes a marginal or conditional density.

First suppose \(d\geq2\). Use the cells \(E_1,E_d\) and define \(E_j=\varnothing\) for \(2\leq j\leq d-1\). Proposition~\ref{prop:step-009-two-pivot-charts} makes this a measurable legal partition for the \(d\) features. Applying both chart inequalities gives
\begin{equation}
\label{eq:appendix-display-324}
\begin{aligned}
\Pr_{\alpha\sim\mu}[\exists\theta\in J:p_\alpha(\theta)=0]
&\leq
\kappa\int_{E_1}\int_{[-R,R]^{d-1}}
\mathbf 1\{\lvert T_1(\theta,\beta)\rvert\leq R\}
\lvert\partial_\theta T_1(\theta,\beta)\rvert\,d\beta\,d\theta\\
&\quad+
\kappa\int_{E_d}\int_{[-R,R]^{d-1}}
\mathbf 1\{\lvert T_d(\theta,\beta)\rvert\leq R\}
\lvert\partial_\theta T_d(\theta,\beta)\rvert\,d\beta\,d\theta\\
&\leq
\kappa\int_{E_1}\int_{[-R,R]^{d-1}}
\lvert\partial_\theta T_1(\theta,\beta)\rvert\,d\beta\,d\theta\\
&\quad+
\kappa\int_{E_d}\int_{[-R,R]^{d-1}}
\lvert\partial_\theta T_d(\theta,\beta)\rvert\,d\beta\,d\theta.
\end{aligned}
\end{equation}
No other chart contributes because its cell is empty.

Set only as inherited shorthand
\begin{equation}
\label{eq:appendix-display-325}
V_d:=d+\frac{Rd(d-1)}2.
\end{equation}
Lemmas~\ref{lem:step-009-inner-velocity} and \ref{lem:step-009-outer-velocity} give, on their complete chart domains,
\begin{equation}
\label{eq:appendix-display-326}
\lvert\partial_\theta T_1\rvert\leq V_d,
\qquad
\lvert\partial_\theta T_d\rvert
\leq1+\frac{Rd(d-1)}2
\leq V_d.
\end{equation}
The outer-chart bound in \eqref{eq:appendix-display-326} is valid for both
signs of \(\theta\). Lemma~\ref{lem:step-009-boundary-index-closure} assigns
\(0,\pm1\) to the legal inner branch and confirms that empty cells and
arbitrary interval locations require no alteration.
Lemma~\ref{lem:step-010-measure-accounting} now yields
\begin{equation}
\label{eq:appendix-display-327}
\begin{aligned}
\Pr_{\alpha\sim\mu}[\exists\theta\in J:p_\alpha(\theta)=0]
&\leq
\kappa V_d(2R)^{d-1}\lvert E_1\rvert
+\kappa V_d(2R)^{d-1}\lvert E_d\rvert\\
&=
\kappa(2R)^{d-1}V_d
(\lvert E_1\rvert+\lvert E_d\rvert)\\
&=
\kappa(2R)^{d-1}V_d\lvert J\rvert.
\end{aligned}
\end{equation}
The two chart integrals charge disjoint parameter cells; they do not produce a factor of two.

Now suppose \(d=1\). Proposition~\ref{prop:step-009-linear-case} supplies the sole legal cell \(E_1=J\), the chart \(T_1=-\theta\), and
\begin{equation}
\label{eq:appendix-display-328}
\lvert\partial_\theta T_1\rvert=1=V_1.
\end{equation}
The \(N=1\) case of Proposition~\ref{prop:step-003-pivot-sweep} and Lemma~\ref{lem:step-010-measure-accounting} give
\begin{equation}
\label{eq:appendix-display-329}
\begin{aligned}
\Pr_{\alpha_0\sim\mu}[\exists\theta\in J:p_{\alpha_0}(\theta)=0]
&\leq
\kappa\int_J\int_{[-R,R]^0}
\mathbf 1\{\lvert T_1(\theta)\rvert\leq R\}
\lvert\partial_\theta T_1(\theta)\rvert\,d\lambda_0\,d\theta\\
&\leq
\kappa\int_J\int_{[-R,R]^0}1\,d\lambda_0\,d\theta\\
&=\kappa\lvert J\rvert\\
&=\kappa(2R)^0
\left(1+\frac{R\cdot1\cdot0}{2}\right)\lvert J\rvert.
\end{aligned}
\end{equation}

In both branches, Proposition~\ref{prop:step-009-two-pivot-charts} or Proposition~\ref{prop:step-009-linear-case} inserts the pivot into the same original coefficient vector and gives \(p_{\Psi_j(\theta,\beta)}(\theta)=0\). Thus the chart vector represents the identical polynomial root event. 
\end{proof}

\begin{lemma}[Degenerate bounded intervals are law-null]
\label{lem:step-010-degenerate-interval}
Under Assumption~\ref{assump:cube-density-laws} and Proposition~\ref{prop:step-009-monic-presentation}, let \(d\geq1\), let \(J\subset\mathbb R\) be a bounded interval with \(\lvert J\rvert=0\), and let \(\alpha\sim\mu\) have any Borel full joint density supported on \([-R,R]^d\) with \(f_\mu\leq\kappa\) almost everywhere. Then
\begin{equation}
\label{eq:appendix-display-330}
\Pr_{\alpha\sim\mu}[\exists\theta\in J:p_\alpha(\theta)=0]=0.
\end{equation}
Consequently,
\begin{equation}
\label{eq:appendix-display-331}
\Pr_{\alpha\sim\mu}[\exists\theta\in J:p_\alpha(\theta)=0]
\leq
\kappa(2R)^{d-1}
\left(d+\frac{Rd(d-1)}2\right)\lvert J\rvert
=0.
\end{equation}


\end{lemma}

\begin{proof}
If \(J=\varnothing\), the event is empty. Otherwise, a real interval of length zero is a singleton, say \(J=\{\theta_0\}\). Proposition~\ref{prop:step-009-monic-presentation} identifies the event with
\begin{equation}
\label{eq:appendix-display-332}
\left\{\alpha\in[-R,R]^d:
\alpha_0=-\theta_0^d-\sum_{k=1}^{d-1}\alpha_k\theta_0^k\right\}.
\end{equation}
This is a closed affine-hyperplane slice because the coefficient of \(\alpha_0\) is exactly one. For each fixed \((\alpha_1,\ldots,\alpha_{d-1})\), its section in the \(\alpha_0\)-coordinate is empty or a singleton and therefore has one-dimensional Lebesgue measure zero. Tonelli gives zero \(d\)-dimensional Lebesgue measure. When \(d=1\), the nonpivot tuple is empty and the same set is the singleton \(\{\alpha_0=-\theta_0\}\cap[-R,R]\), so the argument remains literal.

Applying the full joint-density cap directly,
\begin{equation}
\label{eq:appendix-display-333}
\Pr_{\alpha\sim\mu}[p_\alpha(\theta_0)=0]
=
\int_{\{\alpha:p_\alpha(\theta_0)=0\}}f_\mu(\alpha)\,d\lambda_d(\alpha)
\leq
\kappa\lambda_d(\{\alpha:p_\alpha(\theta_0)=0\})
=0.
\end{equation}
No independence, marginal or conditional density, polynomial-root theorem, or randomized leading coordinate is used. Since \(\lvert J\rvert=0\), the target right-hand side is also zero. 
\end{proof}

\begin{proposition}[Complete affine-monic presentation, certificate, and root-probability wrapper]
\label{prop:step-010-complete-affine-monic-wrapper}
Under Assumption~\ref{assump:cube-density-laws} in dimension \(d\), Proposition~\ref{prop:step-003-pivot-sweep}, Proposition~\ref{prop:step-009-monic-presentation}, Lemma~\ref{lem:step-009-derivative-shift}, Proposition~\ref{prop:step-009-two-pivot-charts}, Lemmas~\ref{lem:step-009-inner-velocity} and \ref{lem:step-009-outer-velocity}, Proposition~\ref{prop:step-009-linear-case}, Lemma~\ref{lem:step-009-boundary-index-closure}, Lemma~\ref{lem:step-010-measure-accounting}, Proposition~\ref{prop:step-010-positive-length}, and Lemma~\ref{lem:step-010-degenerate-interval}, let \(d\geq1\), \(R>0\), \(0<\kappa<\infty\), and let \(J\subset\mathbb R\) be any bounded interval. Let
\begin{equation}
\label{eq:appendix-display-334}
\alpha=(\alpha_0,\ldots,\alpha_{d-1})\sim\mu,
\end{equation}
where \(\mu\) is any Borel probability law on \(\mathbb R^d\) with one full joint density \(f_\mu\), supported on \([-R,R]^d\), satisfying \(f_\mu\leq\kappa\) almost everywhere; its coordinates may be arbitrarily correlated. Then all conclusions below hold simultaneously.

1. \textbf{Actual augmented presentation and random object.} The tuple and polynomial are
\begin{equation}
\label{eq:appendix-indented-008}
   (F_0,F_1,\ldots,F_d)
   =(\theta^d,1,\theta,\ldots,\theta^{d-1}),
   \qquad
   F=(1,\theta,\ldots,\theta^{d-1}),
\end{equation}
\begin{equation}
\label{eq:appendix-indented-009}
   F_0(\theta)+\langle\alpha,F(\theta)\rangle
   =\theta^d+\sum_{k=0}^{d-1}\alpha_k\theta^k
   =p_\alpha(\theta)
   \quad\text{for every }\theta\in\mathbb R.
\end{equation}
   The vector \(\alpha\in[-R,R]^d\) consists of exactly the \(d\) lower coefficients. The leading coefficient one is deterministic, belongs to \(F_0=\theta^d\), and is not a random or singular additional coordinate. The exact presentation parameters are
\begin{equation}
\label{eq:appendix-indented-010}
   q=M=m=0,\qquad
   \Delta=N=d,\qquad
   A=(2R)^d\kappa.
\end{equation}

2. \textbf{Exact constant derivative-shift certificate.} Index the rows and
   columns by \(0,1,\ldots,d\). Define the constant \((d+1)\)-square matrix by
\begin{equation}
\label{eq:appendix-indented-011}
   B_{rs}=
   \begin{cases}
   d,&(r,s)=(0,d),\\
   k,&(r,s)=(k+1,k)\text{ for some }1\leq k\leq d-1,\\
   0,&\text{otherwise}.
   \end{cases}
\end{equation}
   Equivalently, its exactly nonzero entries are
\begin{equation}
\label{eq:appendix-indented-012}
   B_{0,d}=d,\qquad
   B_{k+1,k}=k\quad(1\leq k\leq d-1),
\end{equation}
   and
\begin{equation}
\label{eq:appendix-indented-013}
   \widetilde F'(\theta)=B\widetilde F(\theta)
   \quad\text{for every }\theta.
\end{equation}
   Its exact setting certificate is
\begin{equation}
\label{eq:appendix-indented-014}
   \widehat\Lambda_{B,T}
   =\left(\sum_{k=1}^{d}k^2\right)^{1/2}.
\end{equation}
   Because \(m=0\), this equality contains no interval-location factor and no hidden \(T\) factor. At \(d=1\),
\begin{equation}
\label{eq:appendix-indented-015}
   B=\begin{pmatrix}0&1\\0&0\end{pmatrix},
   \qquad
   \widehat\Lambda_{B,T}=1.
\end{equation}

3. \textbf{Legal charts and velocities for \(d\geq2\).} The cells
\begin{equation}
\label{eq:appendix-indented-016}
   E_1=J\cap\{\lvert\theta\rvert\leq1\},
   \qquad
   E_d=J\cap\{\lvert\theta\rvert>1\},
   \qquad
   E_j=\varnothing\quad(2\leq j\leq d-1)
\end{equation}
   form a measurable disjoint legal partition of \(J\), with
\begin{equation}
\label{eq:appendix-indented-017}
   \lvert E_1\rvert+\lvert E_d\rvert=\lvert J\rvert.
\end{equation}
   On \(E_1\), the legal pivot is the original coefficient \(\alpha_0\) multiplying \(F_1=1\), and
\begin{equation}
\label{eq:appendix-indented-018}
   T_1(\theta,\beta)
   =-\theta^d-\sum_{k=1}^{d-1}\beta_k\theta^k,
   \qquad
   \Psi_1(\theta,\beta)
   =(T_1(\theta,\beta),\beta_1,\ldots,\beta_{d-1}).
\end{equation}
   On \(E_d\), the legal pivot is the original coefficient \(\alpha_{d-1}\) multiplying \(F_d=\theta^{d-1}\neq0\), and
\begin{equation}
\label{eq:appendix-indented-019}
   T_d(\theta,\beta)
   =-\theta-\sum_{k=0}^{d-2}\beta_k\theta^{k-d+1},
   \qquad
   \Psi_d(\theta,\beta)
   =(\beta_0,\ldots,\beta_{d-2},T_d(\theta,\beta)).
\end{equation}
   In both cases \(\beta\in[-R,R]^{d-1}\), and
\begin{equation}
\label{eq:appendix-indented-020}
   \lambda_{d-1}([-R,R]^{d-1})=(2R)^{d-1}.
\end{equation}
   The insertion maps are in the original order \((\alpha_0,\ldots,\alpha_{d-1})\), and
\begin{equation}
\label{eq:appendix-indented-021}
   F_0(\theta)+\langle\Psi_1(\theta,\beta),F(\theta)\rangle
   =p_{\Psi_1(\theta,\beta)}(\theta)=0,
\end{equation}
\begin{equation}
\label{eq:appendix-indented-022}
   F_0(\theta)+\langle\Psi_d(\theta,\beta),F(\theta)\rangle
   =p_{\Psi_d(\theta,\beta)}(\theta)=0.
\end{equation}
   Thus either chart vector gives the identical target polynomial. The inserted vector belongs to the support cube exactly when its pivot also satisfies \(\lvert T_1\rvert\leq R\) or \(\lvert T_d\rvert\leq R\), respectively.

   With
\begin{equation}
\label{eq:appendix-indented-023}
   V_d:=d+\frac{Rd(d-1)}2,
\end{equation}
   the exact derivative formulas and bounds are
\begin{equation}
\label{eq:appendix-indented-024}
   \partial_\theta T_1
   =-d\theta^{d-1}-\sum_{k=1}^{d-1}k\beta_k\theta^{k-1},
   \qquad
   \lvert\partial_\theta T_1\rvert\leq V_d,
\end{equation}
\begin{equation}
\label{eq:appendix-indented-025}
   \partial_\theta T_d
   =-1+\sum_{k=0}^{d-2}(d-1-k)\beta_k\theta^{k-d},
\end{equation}
\begin{equation}
\label{eq:appendix-indented-026}
   \lvert\partial_\theta T_d\rvert
   \leq
   1+R\sum_{m=1}^{d-1}\frac{m}{\lvert\theta\rvert^{m+1}}
   \leq1+\frac{Rd(d-1)}2
   \leq V_d.
\end{equation}

4. \textbf{Separate linear chart.} At \(d=1\),
\begin{equation}
\label{eq:appendix-indented-027}
   (F_0,F_1)=(\theta,1),\qquad
   p_{\alpha_0}(\theta)=\theta+\alpha_0,
   \qquad
   E_1=J,
\end{equation}
   and the beta space is genuinely zero-dimensional:
\begin{equation}
\label{eq:appendix-indented-028}
   [-R,R]^0=\{()\},\qquad
   \lambda_0([-R,R]^0)=1=(2R)^0.
\end{equation}
   The sole legal chart is
\begin{equation}
\label{eq:appendix-indented-029}
   T_1(\theta)=-\theta,\qquad
   \Psi_1(\theta)=(-\theta),\qquad
   p_{\Psi_1(\theta)}(\theta)=0,
\end{equation}
   and
\begin{equation}
\label{eq:appendix-indented-030}
   \lvert\partial_\theta T_1\rvert=1=V_1.
\end{equation}
   No beta coordinate is introduced.

5. \textbf{Boundary, dimension, and location closure.} The inner chart handles \(\theta=0,\theta=1,\theta=-1\); the outer chart handles both signs of \(\theta\) whenever \(\lvert\theta\rvert>1\). Empty \(J\), empty active cells, and intervals lying in only one region require no change. At \(d=2\), the generic formulas specialize to
\begin{equation}
\label{eq:appendix-indented-031}
   (F_0,F_1,F_2)=(\theta^2,1,\theta),
   \qquad
   B=\begin{pmatrix}0&0&2\\0&0&0\\0&1&0\end{pmatrix},
   \qquad
   \widehat\Lambda_{B,T}=\sqrt5,
\end{equation}
\begin{equation}
\label{eq:appendix-indented-032}
   T_1=-\theta^2-\beta_1\theta,
   \qquad
   T_2=-\theta-\beta_0\theta^{-1},
   \qquad
   \lvert\partial_\theta T_1\rvert\leq2+R,
   \qquad
   \lvert\partial_\theta T_2\rvert\leq1+R\leq2+R.
\end{equation}
   Every coefficient vector has dimension \(d\), every beta cube has dimension \(d-1\), and \(B\) has dimension \(d+1\). None of the presentation, certificate, chart, velocity, or probability conclusions depends on the location of \(J\).

6. \textbf{Exact ordinary-probability bridge for every bounded interval.} For the same original law \(\mu\), with no independence assumption and for every bounded interval \(J\), including empty and singleton intervals,
\begin{equation}
\label{eq:appendix-indented-033}
   \boxed{
   \Pr_{\alpha\sim\mu}[\exists\theta\in J:p_\alpha(\theta)=0]
   \leq
   \kappa(2R)^{d-1}
   \left(d+\frac{Rd(d-1)}2\right)\lvert J\rvert.}
\end{equation}
   When \(\lvert J\rvert=0\), the event is empty or a proper affine-hyperplane slice and has probability zero under the same full joint density. When \(d=1\), this bound is exactly \(\kappa\lvert J\rvert\).


\end{proposition}

\begin{proof}
Proposition~\ref{prop:step-009-monic-presentation} gives clause 1 exactly. Lemma~\ref{lem:step-009-derivative-shift} gives clause 2 exactly, including the \(d=1\) matrix and the absence of a \(T\) factor. Proposition~\ref{prop:step-009-two-pivot-charts}, Lemmas~\ref{lem:step-009-inner-velocity} and \ref{lem:step-009-outer-velocity}, and Lemma~\ref{lem:step-010-measure-accounting} give clause 3. Proposition~\ref{prop:step-009-linear-case} and Lemma~\ref{lem:step-010-measure-accounting} give clause 4. Lemma~\ref{lem:step-009-boundary-index-closure} gives clause 5, including the explicit \(d=2\) specialization and arbitrary interval location.

For \(\lvert J\rvert>0\), Proposition~\ref{prop:step-010-positive-length} applies Proposition~\ref{prop:step-003-pivot-sweep} to these exact original-coordinate charts, uses the same full joint-density cap, and proves clause 6. For \(\lvert J\rvert=0\), Lemma~\ref{lem:step-010-degenerate-interval} proves clause 6 by affine-hyperplane nullity. These cases exhaust every bounded interval.

Every item in the statement is therefore either a conclusion of the named
results or the integration and nullity argument proved here. This completes
the proposition. 
\end{proof}


\begin{proof}[Synthesis of affine-monic probability recovery]
Proposition~\ref{prop:step-010-complete-affine-monic-wrapper} combines the
exact $d$-coordinate monic presentation, constant derivative-shift matrix,
certificate, legal chart cells, insertion maps, regional velocity bounds,
zero-dimensional $d=1$ chart, and all sign, boundary, dimension, empty-cell,
and interval-location cases.  For $|J|>0$ its probability clause follows
from Proposition~\ref{prop:step-010-positive-length}, which applies
Proposition~\ref{prop:step-003-pivot-sweep} to the original coefficient
coordinates and uses the exact beta-cube volume from
Lemma~\ref{lem:step-010-measure-accounting}.  For $|J|=0$,
Lemma~\ref{lem:step-010-degenerate-interval} gives probability zero under
the same full joint density.  Hence, for every bounded interval and every
possibly correlated admissible law,
\begin{equation}
\label{eq:appendix-display-335}
\Pr_{\alpha\sim\mu}[\exists\theta\in J:p_\alpha(\theta)=0]
\leq\kappa(2R)^{d-1}
\left(d+\frac{Rd(d-1)}2\right)|J|.
\end{equation}
This is one complete presentation, certificate, and probability result;
it uses neither an independent polynomial-root theorem nor a random leading
coordinate.
\end{proof}


\subsection{Counter-example 1 Scale Calculation}\label{app:step-011}

\begin{lemma}[Exact one-entry derivative certificate]
\label{lem:step-011-shear-certificate}

Under Assumptions~\ref{assump:parameter-regime}, \ref{assump:balcan-common-chain}, and
\ref{assump:anchored-derivative-closure}, suppose \(0<\delta\leq1\), \(\Theta=[-1,1]\), and, in the exact
coordinate ordering \(0,1,2\),

\begin{equation}
\label{eq:appendix-display-336}
\widetilde F(\theta)
=
\begin{pmatrix}
F_0(\theta)\\ F_1(\theta)\\ F_2(\theta)
\end{pmatrix}
=
\begin{pmatrix}
0\\ 1\\ \theta/\delta
\end{pmatrix}.
\end{equation}

Then

\begin{equation}
\label{eq:appendix-display-337}
q=0,\qquad M=0,\qquad \Delta=1,\qquad N=2,\qquad m=0,
\end{equation}

and the constant matrix

\begin{equation}
\label{eq:appendix-display-338}
B
=
\begin{pmatrix}
0&0&0\\
0&0&0\\
0&1/\delta&0
\end{pmatrix}
\end{equation}

has the sole nonzero entry \(B_{2,1}=1/\delta\), satisfies

\begin{equation}
\label{eq:appendix-display-339}
\widetilde F'(\theta)=B\widetilde F(\theta),
\end{equation}

and has exact static certificate

\begin{equation}
\label{eq:appendix-display-340}
\widehat\Lambda_{B,T}=\frac1\delta.
\end{equation}

This value is independent of the legal containing parameter \(T\) satisfying
\([-1,1]\subseteq[-T,T]\).


\end{lemma}

\begin{proof}
The chain is empty, so \(q=0\) and the setting convention gives \(M=0\). The three
output polynomials are \(Q_0=0\), \(Q_1=1\), and \(Q_2(\theta)=\theta/\delta\); their maximum total degree is
one, so \(\Delta=1\). There are two random-coordinate features, hence \(N=2\). This matrix is constant,
so \(m=0\), and direct multiplication in the stated row-column ordering gives

\begin{equation}
\label{eq:appendix-display-341}
B\widetilde F(\theta)
=
\begin{pmatrix}
0\\ 0\\ (1/\delta)F_1(\theta)
\end{pmatrix}
=
\begin{pmatrix}
0\\ 0\\ 1/\delta
\end{pmatrix}
=\widetilde F'(\theta).
\end{equation}

Thus this is the supplied anchored closure matrix, with \(F_1\equiv1\) as the anchor. Its coefficient list has
only \(b_{2,1,0}=1/\delta\). Since \(m=0\), \(T_*^0=1\), and therefore the setting definition gives

\begin{equation}
\label{eq:appendix-display-342}
\begin{aligned}
\widehat\Lambda_{B,T}
&=\left(
\sum_{r=0}^{2}\sum_{s=0}^{2}
\left(\sum_{\ell=0}^{0}\lvert b_{rs,\ell}\rvert T_*^\ell\right)^2
\right)^{1/2}\\
&=\left(\left\lvert\frac1\delta\right\rvert^2\right)^{1/2}
=\frac1\delta,
\end{aligned}
\end{equation}

where the last equality uses \(\delta>0\). This direct calculation agrees with the one-entry-shear
specialization in Proposition~\ref{prop:step-001-boundary}. Since the only power is \(T_*^0=1\), changing the
legal containing \(T\) changes neither the coefficient list nor the certificate. The certificate is therefore
a raw Euclidean coefficient-height calculation, not a value inferred from a later probability claim.

\end{proof}

\begin{lemma}[Euclidean projective speed]
\label{lem:step-011-projective-speed}

Under Assumptions~\ref{assump:parameter-regime},
\ref{assump:balcan-common-chain}, and
\ref{assump:anchored-derivative-closure}, and
Lemma~\ref{lem:step-011-shear-certificate}, let
\(F(\theta)=(1,\theta/\delta)\) on \(\Theta=[-1,1]\), where
\(0<\delta\leq1\), and set \(x=\theta/\delta\). Then

\begin{equation}
\label{eq:appendix-display-343}
\gamma_F(\theta)=\frac{(1,x)}{\sqrt{1+x^2}},
\qquad
\gamma_F'(\theta)
=\frac{(-x,1)}{\delta(1+x^2)^{3/2}},
\end{equation}

so

\begin{equation}
\label{eq:appendix-display-344}
\lVert\gamma_F'(\theta)\rVert_2
=\frac{1}{\delta(1+x^2)},
\qquad
\Gamma_{\mathrm{proj}}(F)=\frac1\delta.
\end{equation}

At \(\theta=0\),

\begin{equation}
\label{eq:appendix-display-345}
\gamma_F(0)=(1,0),
\qquad
\gamma_F'(0)=(0,1/\delta),
\qquad
\lVert\gamma_F'(0)\rVert_2=1/\delta.
\end{equation}


\end{lemma}

\begin{proof}
The Euclidean norm is

\begin{equation}
\label{eq:appendix-display-346}
\lVert F(\theta)\rVert_2=\sqrt{1+x^2}>0,
\end{equation}

so normalization is valid everywhere, including the endpoints. Since \(x'=1/\delta\), direct differentiation
gives

\begin{equation}
\label{eq:appendix-display-347}
\begin{aligned}
\gamma_F'(\theta)
&=\frac1\delta\frac{d}{dx}
\left((1,x)(1+x^2)^{-1/2}\right)\\
&=\frac1\delta
\left(
-\frac{x}{(1+x^2)^{3/2}},
\frac{1}{(1+x^2)^{3/2}}
\right).
\end{aligned}
\end{equation}

Taking the Euclidean norm, without changing the parameter coordinate, yields

\begin{equation}
\label{eq:appendix-display-348}
\begin{aligned}
\lVert\gamma_F'(\theta)\rVert_2
&=\frac1\delta
\frac{\sqrt{x^2+1}}{(1+x^2)^{3/2}}\\
&=\frac{1}{\delta(1+x^2)}.
\end{aligned}
\end{equation}

Because \(1+x^2\geq1\), this is at most \(1/\delta\). The point \(\theta=0\in[-1,1]\) has \(x=0\) and
attains equality, so the ordinary supremum is exactly \(1/\delta\). Substitution at \(x=0\) gives the stated
active vector value. The projective certificate from Proposition~\ref{prop:step-001-projective}
therefore holds with equality for this instance. 
\end{proof}

\begin{proposition}[Exact homogeneous upper coefficient]
\label{prop:step-011-upper-coefficient}

Under Assumptions~\ref{assump:parameter-regime}, \ref{assump:balcan-common-chain},
\ref{assump:cube-density-laws}, and \ref{assump:anchored-derivative-closure}, the Propositions~\ref{prop:step-008-s2-pairwise-homogeneous-rate} and
\ref{prop:step-008-s2-pf-closure}, and Lemmas~\ref{lem:step-011-shear-certificate} and
\ref{lem:step-011-projective-speed}, specialize to the following statement. For every possibly correlated law
\(\mu\in\mathcal D_{2,1,1/4}\) and every interval \(I\subseteq[-1,1]\) with \(\lvert I\rvert>0\),

\begin{equation}
\label{eq:appendix-display-349}
\Pr_{\alpha\sim\mu}
[\exists\theta\in I:\alpha_1+\alpha_2\theta/\delta=0]
\leq\frac1\delta\lvert I\rvert,
\end{equation}

and

\begin{equation}
\label{eq:appendix-display-350}
C^{\mathrm{Pf}}_{\mathcal D}(F;[-1,1])\leq\frac1\delta.
\end{equation}


\end{proposition}

\begin{proof}
For \(N=2\), \(R=1\), and \(\kappa=1/4\), the setting quantity is

\begin{equation}
\label{eq:appendix-display-351}
A=(2R)^N\kappa=2^2\cdot\frac14=1,
\qquad
A\sqrt{\frac N2}=1\cdot\sqrt{\frac22}=1.
\end{equation}

Proposition~\ref{prop:step-008-s2-pairwise-homogeneous-rate} applies to every law in the original class,
without an independence hypothesis, and gives

\begin{equation}
\label{eq:appendix-display-352}
\begin{aligned}
\Pr_{\alpha\sim\mu}
[\exists\theta\in I:\alpha_1+\alpha_2\theta/\delta=0]
&\leq
A\sqrt{\frac N2}\,\Gamma_{\mathrm{proj}}(F)\lvert I\rvert\\
&=\frac1\delta\lvert I\rvert.
\end{aligned}
\end{equation}

Here Lemma~\ref{lem:step-011-projective-speed} supplies the exact equality
\(\Gamma_{\mathrm{proj}}(F)=1/\delta\). The certificate version has the identical coefficient because
Lemma~\ref{lem:step-011-shear-certificate} gives
\(\widehat\Lambda_{B,T}=1/\delta\). Proposition~\ref{prop:step-008-s2-pf-closure}, with the same substitution,
gives the capacity bound. Neither bound is clipped at one. No feature of the
particular independent uniform law is used in this upper bound. 
\end{proof}

\begin{lemma}[Exact root-wedge equivalence]
\label{lem:step-011-root-wedges}

Under Assumptions~\ref{assump:parameter-regime} and \ref{assump:anchored-derivative-closure},
Lemma~\ref{lem:step-011-shear-certificate}, and \(0<\epsilon\leq\delta\leq1\), define

\begin{equation}
\label{eq:appendix-display-353}
t:=\frac\epsilon\delta\in(0,1],
\end{equation}

and the two closed subsets of \([-1,1]^2\)

\begin{equation}
\label{eq:appendix-display-354}
\begin{aligned}
W_+(t)
&:=\{(\alpha_1,\alpha_2):0\leq\alpha_2\leq1,
-t\alpha_2\leq\alpha_1\leq0\},\\
W_-(t)
&:=\{(\alpha_1,\alpha_2):-1\leq\alpha_2\leq0,
0\leq\alpha_1\leq-t\alpha_2\}.
\end{aligned}
\end{equation}

Then the full coefficient event is exactly

\begin{equation}
\label{eq:appendix-display-355}
\{\alpha\in[-1,1]^2:\exists\theta\in[0,\epsilon],
\ \alpha_1+\alpha_2\theta/\delta=0\}
=W_+(t)\cup W_-(t).
\end{equation}

Off the coefficient axes this is equivalently the union of both opposite-sign wedges

\begin{equation}
\label{eq:appendix-display-356}
\alpha_1\alpha_2<0,
\qquad
\lvert\alpha_1\rvert\leq t\lvert\alpha_2\rvert.
\end{equation}

The boundary \(\alpha_1=0\) gives the endpoint root \(\theta=0\), the equality boundary
\(\lvert\alpha_1\rvert=t\lvert\alpha_2\rvert\) gives the endpoint root \(\theta=\epsilon\), and the only root
with \(\alpha_2=0\) is the origin, where the function is identically zero.


\end{lemma}

\begin{proof}
Write \(u=\theta/\delta\). Since \(\delta>0\), the range
\(\theta\in[0,\epsilon]\) is exactly \(u\in[0,t]\), and the root equation becomes

\begin{equation}
\label{eq:appendix-display-357}
\alpha_1+\alpha_2u=0.
\end{equation}

If \(\alpha_2>0\), the unique candidate is \(u=-\alpha_1/\alpha_2\), and

\begin{equation}
\label{eq:appendix-display-358}
0\leq-\frac{\alpha_1}{\alpha_2}\leq t
\quad\Longleftrightarrow\quad
-t\alpha_2\leq\alpha_1\leq0.
\end{equation}

This is exactly the positive-\(\alpha_2\) part of \(W_+(t)\). If \(\alpha_2<0\), division reverses the
inequalities and gives

\begin{equation}
\label{eq:appendix-display-359}
0\leq-\frac{\alpha_1}{\alpha_2}\leq t
\quad\Longleftrightarrow\quad
0\leq\alpha_1\leq-t\alpha_2,
\end{equation}

which is exactly the negative-\(\alpha_2\) part of \(W_-(t)\). If \(\alpha_2=0\), the equation has a root if
and only if \(\alpha_1=0\); this origin belongs to both closed wedges. These three exhaustive cases prove the
event equality.

Because \(t\leq1\), the bounds \(\lvert\alpha_1\rvert\leq t\lvert\alpha_2\rvert\leq1\) keep both wedges
inside the coefficient square, including at \(t=1\). Away from the axes, the inequalities force strict opposite
signs. The candidate root is \(u=0\) when \(\alpha_1=0\) and \(u=t\) on the sloping equality boundary, proving
the two endpoint assertions rather than discarding their coefficient sets. 
\end{proof}

\begin{proposition}[Exact wedge probability for the uniform source law]
\label{prop:step-011-wedge-probability}

Under Assumption~\ref{assump:cube-density-laws}, let \(\alpha_1,\alpha_2\) be independent uniform random
variables on \([-1,1]\), solely for the required source example. Under the event identity of
Lemma~\ref{lem:step-011-root-wedges}, for every \(0<\epsilon\leq\delta\leq1\),

\begin{equation}
\label{eq:appendix-display-360}
\lambda_2(W_+(\epsilon/\delta)\cup W_-(\epsilon/\delta))
=\frac\epsilon\delta,
\end{equation}

and

\begin{equation}
\label{eq:appendix-display-361}
\Pr[\exists\theta\in[0,\epsilon]:
\alpha_1+\alpha_2\theta/\delta=0]
=\frac{\epsilon}{4\delta}.
\end{equation}


\end{proposition}

\begin{proof}
Retain \(t=\epsilon/\delta\in(0,1]\). At a fixed
\(\alpha_2=s\in[0,1]\), the \(W_+(t)\) slice in the \(\alpha_1\)-coordinate is
\([-ts,0]\), of length \(ts\). Hence

\begin{equation}
\label{eq:appendix-display-362}
\lambda_2(W_+(t))
=\int_0^1 ts\,ds
=\frac t2.
\end{equation}

For \(W_-(t)\), substitute \(s=-\alpha_2\in[0,1]\); its slice length is again \(ts\), so

\begin{equation}
\label{eq:appendix-display-363}
\lambda_2(W_-(t))
=\int_0^1 ts\,ds
=\frac t2.
\end{equation}

The two wedges intersect only at the origin. Therefore

\begin{equation}
\label{eq:appendix-display-364}
\lambda_2(W_+(t)\cup W_-(t))
=2\int_0^1ts\,ds
=t
=\frac\epsilon\delta.
\end{equation}

This is the combined area of both sign quadrants, not the area of one wedge. The coefficient axes, the two
sloping equality edges, and the two outer square-edge segments are finite unions of line segments and hence
have planar Lebesgue measure zero. All closed boundaries were included in the exact event in
Lemma~\ref{lem:step-011-root-wedges}; their nullity is used only when evaluating area, not to remove either
endpoint root or any square-boundary coefficient from the event.

The selected product law has joint density exactly \(1/4\) on the square. Integrating that density over the
exact event gives

\begin{equation}
\label{eq:appendix-display-365}
\Pr[\exists\theta\in[0,\epsilon]:
\alpha_1+\alpha_2\theta/\delta=0]
=\frac14\lambda_2(W_+(t)\cup W_-(t))
=\frac{\epsilon}{4\delta}.
\end{equation}

At \(\epsilon=\delta\), \(t=1\), the two triangles have total area one and the probability is \(1/4\). At
\(\delta=1\), the same calculation holds for every \(0<\epsilon\leq1\). The \(\theta=0\) endpoint roots and
the \(\theta=\epsilon\) endpoint roots are already included exactly. 
\end{proof}

\begin{proposition}[Exact normalized upper/lower scale comparison]
\label{prop:step-011-scale-comparison}

Under Lemma~\ref{lem:step-011-shear-certificate} and
Propositions~\ref{prop:step-011-upper-coefficient} and
\ref{prop:step-011-wedge-probability}, for the selected uniform source law and every
\(0<\epsilon\leq\delta\leq1\),

\begin{equation}
\label{eq:appendix-display-366}
\frac{
\Pr[\exists\theta\in[0,\epsilon]:
\alpha_1+\alpha_2\theta/\delta=0]
}{\lvert[0,\epsilon]\rvert}
=\frac1{4\delta}.
\end{equation}

The three quantities are therefore:

\begin{equation}
\label{eq:appendix-display-367}
\begin{aligned}
\text{selected-law lower ratio}&=\frac1{4\delta},\\
\text{general upper coefficient}&=\frac1\delta,\\
\text{raw one-entry certificate}&=\frac1\delta.
\end{aligned}
\end{equation}

The first quantity is smaller by the literal factor four; it is not asserted to equal either of the latter two,
and no optimal upper constant is claimed.
At \(\epsilon=\delta\) and at \(\delta=1\) these values remain literal. For fixed \(\delta>0\), the notation
\(\epsilon\downarrow0\) refers only to the limit through positive interval lengths, for which the ratio remains
\(1/(4\delta)\); no ratio is asserted at \(\epsilon=0\).


\end{proposition}

\begin{proof}
Since \(\epsilon>0\), the interval has positive length
\(\lvert[0,\epsilon]\rvert=\epsilon\), so Proposition~\ref{prop:step-011-wedge-probability} may be divided by
that length:

\begin{equation}
\label{eq:appendix-display-368}
\frac{\epsilon/(4\delta)}{\epsilon}=\frac1{4\delta}.
\end{equation}

Proposition~\ref{prop:step-011-upper-coefficient} supplies \(1/\delta\) from the general theorem, not
from the direct wedge calculation. Separately, Lemma~\ref{lem:step-011-shear-certificate} supplies the raw
Euclidean coefficient height \(1/\delta\) directly from the single matrix entry. Thus the lower test ratio is
\(1/4\) times each comparison scale, while the upper theorem and certificate retain logically distinct
mathematical roles. If \(\epsilon\downarrow0\) with \(\epsilon>0\), the probability
\(\epsilon/(4\delta)\) tends to zero while the normalized ratio remains constant for fixed \(\delta>0\); at
\(\epsilon=0\), division is not defined and is not used. 
\end{proof}


\begin{proof}[Synthesis of the Counter-example 1 scales]
For $0<\delta\leq1$, Lemma~\ref{lem:step-011-shear-certificate} gives
\begin{equation}
\label{eq:appendix-display-369}
(q,M,\Delta,N,m)=(0,0,1,2,0),
\qquad B_{2,1}=\frac1\delta,
\qquad \widehat\Lambda_{B,T}=\frac1\delta,
\end{equation}
with no other nonzero closure entry.  Lemma~\ref{lem:step-011-projective-speed}
independently gives
\begin{equation}
\label{eq:appendix-display-370}
\|\gamma_F'(\theta)\|_2
=\frac1{\delta(1+(\theta/\delta)^2)},
\qquad
\Gamma_{\mathrm{proj}}(F)=\frac1\delta.
\end{equation}
With $A=1$ and $N=2$, Proposition~\ref{prop:step-011-upper-coefficient}
therefore gives the coefficient $1/\delta$ for every admissible, possibly
correlated law.  For the selected uniform law only,
Lemma~\ref{lem:step-011-root-wedges} identifies both closed opposite-sign
wedges, and Proposition~\ref{prop:step-011-wedge-probability} computes their
combined area and probability as
\begin{equation}
\label{eq:appendix-display-371}
\lambda_2(W_+(\epsilon/\delta)\cup W_-(\epsilon/\delta))
=\frac\epsilon\delta,
\qquad
\Pr[\exists\theta\in[0,\epsilon]:
\alpha_1+\alpha_2\theta/\delta=0]
=\frac{\epsilon}{4\delta}.
\end{equation}
Proposition~\ref{prop:step-011-scale-comparison} divides only by positive
$\epsilon$ and gives the selected-law ratio $1/(4\delta)$, distinct from
the all-law upper coefficient and raw certificate, both $1/\delta$.
Multiplying the raw certificate by the literal Ball-section factor gives the
separate geometric normalization
\begin{equation}
\label{eq:appendix-display-372}
\sqrt2\,\widehat\Lambda_{B,T}=\frac{\sqrt2}{\delta};
\end{equation}
this identity is not a probability coefficient or an optimality statement.
The deterministic-leading-coefficient monic baseline is unaffected.
\end{proof}

\subsection{Anchored Coefficient-Sweep Theorem}\label{app:step-012}

\begin{proposition}[Anchored derivative-closure coefficient sweep]
\label{prop:step-012-anchored-coefficient-sweep}

Under Assumptions~\ref{assump:parameter-regime}, \ref{assump:balcan-common-chain},
\ref{assump:anchored-derivative-closure}, and \ref{assump:cube-density-laws}, with each conclusion using only
the assumption subsets stated in the listed results, use the following named
theorem-style conclusions with exactly their stated scopes:

\begin{equation}
\label{eq:appendix-display-373}
\begin{gathered}
\text{Lemma~\ref{lem:step-001-anchor}, Lemma~\ref{lem:step-001-height},
Lemma~\ref{lem:step-001-homogeneous-block},}\\
\text{Proposition~\ref{prop:step-001-projective}, Proposition~\ref{prop:step-001-boundary};}\\
\text{Lemma~\ref{lem:step-003-measurable-domains}, Lemma~\ref{lem:step-003-finite-chart},
Proposition~\ref{prop:step-003-finite-area},}\\
\text{Lemma~\ref{lem:step-003-root-coverage}, Proposition~\ref{prop:step-003-pivot-sweep};}\\
\text{Proposition~\ref{prop:step-007-s2-general-affine-rate};}\\
\text{Lemma~\ref{lem:step-008-s2-radial-cancellation},
Proposition~\ref{prop:step-008-s2-stationary-projective},}\\
\text{Lemma~\ref{lem:step-008-s2-literal-algebra},
Proposition~\ref{prop:step-008-s2-pairwise-homogeneous-rate},
Proposition~\ref{prop:step-008-s2-pf-closure};}\\
\text{Proposition~\ref{prop:step-010-complete-affine-monic-wrapper};}\\
\text{Lemma~\ref{lem:step-011-shear-certificate}, Lemma~\ref{lem:step-011-projective-speed},
Proposition~\ref{prop:step-011-upper-coefficient},}\\
\text{Lemma~\ref{lem:step-011-root-wedges},
Proposition~\ref{prop:step-011-wedge-probability},
Proposition~\ref{prop:step-011-scale-comparison}.}
\end{gathered}
\end{equation}

Then all of the following conclusions hold simultaneously.

1. \textbf{Static certificate and homogeneous closure.} For every \(\theta\in\Theta\),
   
\begin{equation}
\label{eq:appendix-indented-034}
   F_{j_*}(\theta)=1,
   \qquad \lVert F(\theta)\rVert_2\geq1,
   \qquad F(\theta)\neq0,
\end{equation}
   
   and
   
\begin{equation}
\label{eq:appendix-indented-035}
   \lVert B(\theta)\rVert_{\mathrm{op}}
   \leq\lVert B(\theta)\rVert_{\mathrm F}
   \leq\widehat\Lambda_{B,T},
   \qquad
   \sup_{\vartheta\in\Theta}\lVert B(\vartheta)\rVert_{\mathrm{op}}
   \leq\widehat\Lambda_{B,T}.
\end{equation}
   
   If \(F_0\equiv0\), then
   
\begin{equation}
\label{eq:appendix-indented-036}
   F'=B_FF,
   \qquad
   \gamma_F'
   =(I_N-\gamma_F\gamma_F^{\mathsf T})B_F\gamma_F,
   \qquad
   \Gamma_{\mathrm{proj}}(F)\leq\widehat\Lambda_{B,T}.
\end{equation}

   These conclusions remain literal for \(q=0\), \(m=0\), constant \(B\), constant \(\widetilde F\),
   \(N=1\), both endpoints of \(\Theta\), zero certificate height, and stationary \(\gamma_F\). The two
   fixed-presentation certificate reductions are

\begin{equation}
\label{eq:appendix-indented-037}
   \widehat\Lambda_{B,T}=\left(\sum_{k=1}^d k^2\right)^{1/2}
   \quad\text{for the augmented monic shift},
   \qquad
   \widehat\Lambda_{B,T}=\frac1\delta
   \quad\text{for the one-entry shear}.
\end{equation}

2. \textbf{Equivalent chart form and general affine theorem.} For every
   \(\mu\in\mathcal D_{N,R,\kappa}\), every interval \(I\subseteq\Theta\) with
   \(\lvert I\rvert>0\), and every Lebesgue-measurable legal partition
   \(I=\bigsqcup_{j=1}^N E_j\) with \(F_j\neq0\) on \(E_j\),
   
\begin{equation}
\label{eq:appendix-indented-038}
   \begin{aligned}
   \Pr_{\alpha\sim\mu}[\exists\theta\in I:\phi_\alpha(\theta)=0]
   &\leq
   \kappa\sum_{j=1}^{N}\int_{E_j}\int_{[-R,R]^{N-1}}
   \mathbf 1\{\lvert T_j(\theta,\beta)\rvert\leq R\}
   \lvert\partial_\theta T_j(\theta,\beta)\rvert\,d\beta\,d\theta\\
   &\leq
   \kappa\sum_{j=1}^{N}\int_{E_j}\int_{[-R,R]^{N-1}}
   \lvert\partial_\theta T_j(\theta,\beta)\rvert\,d\beta\,d\theta.
   \end{aligned}
\end{equation}
   
   Independently of the choice of legal partition, every such law and positive-length interval satisfies the
   coordinate-free chain
   
\begin{equation}
\label{eq:appendix-indented-039}
   \begin{aligned}
   \Pr_{\alpha\sim\mu}[\exists\theta\in I:\phi_\alpha(\theta)=0]
   &\leq
   \kappa\int_I\int_{H_\theta\cap[-R,R]^N}
   \frac{\lvert F_0'(\theta)+\langle a,F'(\theta)\rangle\rvert}
   {\lVert F(\theta)\rVert_2}
   \,d\mathcal H^{N-1}(a)\,d\theta\\
   &\leq
   \kappa\sqrt2(2R)^{N-1}(1+NR^2)\widehat\Lambda_{B,T}\lvert I\rvert\\
   &=\frac{A(1+NR^2)\widehat\Lambda_{B,T}}{\sqrt2R}\lvert I\rvert,
   \end{aligned}
\end{equation}
   
   and
   
\begin{equation}
\label{eq:appendix-indented-040}
   C^{\mathrm{aff}}_{\mathcal D}(F_0,F;\Theta)
   \leq\frac{A(1+NR^2)\widehat\Lambda_{B,T}}{\sqrt2R}.
\end{equation}
   
   If \(\widehat\Lambda_{B,T}=0\), the event probability is exactly zero for every such law and interval.
   The chart conclusion includes completed-measurable validation, finite localized charts, exact area/multiplicity
   multiplicity, persistent-root probability zero, arbitrary correlation, empty cells, zero Jacobians,
   tangent and multiple roots, every endpoint convention, pivots with no uniform margin, and \(N=1\), without
   adding the finite exhaustion variables to the conclusion.

3. \textbf{Sharper homogeneous theorem.} If \(F_0\equiv0\), then on the actual central section
   
\begin{equation}
\label{eq:appendix-indented-041}
   H_\theta=F(\theta)^\perp=\gamma_F(\theta)^\perp,
   \qquad
   \frac{\lvert\langle a,F'(\theta)\rangle\rvert}{\lVert F(\theta)\rVert_2}
   =\lvert\langle a,\gamma_F'(\theta)\rangle\rvert
   \quad(a\in H_\theta).
\end{equation}
   
   For every admissible law and every positive-length interval,
   
\begin{equation}
\label{eq:appendix-indented-042}
   \Pr_{\alpha\sim\mu}[\exists\theta\in I:\langle\alpha,F(\theta)\rangle=0]
   \leq A\sqrt{\frac N2}\,\Gamma_{\mathrm{proj}}(F)\lvert I\rvert
   \leq A\sqrt{\frac N2}\,\widehat\Lambda_{B,T}\lvert I\rvert,
\end{equation}
   
   and
   
\begin{equation}
\label{eq:appendix-indented-043}
   C^{\mathrm{Pf}}_{\mathcal D}(F;\Theta)
   \leq A\sqrt{\frac N2}\,\widehat\Lambda_{B,T}.
\end{equation}
   
   If \(\Gamma_{\mathrm{proj}}(F)=0\), the normalized curve is constant, every nonempty-interval root event is
   one fixed proper central hyperplane, and every admissible full-density law assigns it probability zero.
   The coefficient identity
   
\begin{equation}
\label{eq:appendix-indented-044}
   \kappa R\sqrt N\,\sqrt2(2R)^{N-1}=A\sqrt{\frac N2}
\end{equation}
   
   is literal. The statements retain \(N=1\), stationary and zero-certificate branches, radial rescaling,
   \(a=0\), literal endpoints, and arbitrarily short positive intervals.

4. \textbf{Complete exact affine-monic baseline.} For every \(d\geq1\), \(R>0\),
   \(0<\kappa<\infty\), every bounded interval \(J\subset\mathbb R\), and every arbitrary possibly correlated
   \(d\)-dimensional law described in Proposition~\ref{prop:step-010-complete-affine-monic-wrapper}, all six clauses of
   Proposition~\ref{prop:step-010-complete-affine-monic-wrapper} hold verbatim. In particular, the exact tuple,
   object identity, descriptor values, and certificate are
   
\begin{equation}
\label{eq:appendix-indented-045}
   (F_0,F_1,\ldots,F_d)=(\theta^d,1,\theta,\ldots,\theta^{d-1}),
   \quad
   p_\alpha(\theta)=F_0(\theta)+\langle\alpha,F(\theta)\rangle
   =\theta^d+\sum_{k=0}^{d-1}\alpha_k\theta^k,
\end{equation}
   
\begin{equation}
\label{eq:appendix-indented-046}
   q=M=m=0,
   \quad \Delta=N=d,
   \quad A=(2R)^d\kappa,
   \quad
   \widehat\Lambda_{B,T}=\left(\sum_{k=1}^{d}k^2\right)^{1/2},
\end{equation}
   
   with exactly the derivative-shift entries
   
\begin{equation}
\label{eq:appendix-indented-047}
   B_{0,d}=d,
   \qquad B_{k+1,k}=k\quad(1\leq k\leq d-1).
\end{equation}
   
   The random vector contains exactly the \(d\) lower coefficients; the leading coefficient one is
   deterministic. For \(d\geq2\), the exact legal cells, original-coordinate pivots, charts, insertion maps,
   beta volume, same-polynomial identities, and velocity formulas and bounds are
   
\begin{equation}
\label{eq:appendix-indented-048}
   E_1=J\cap\{\lvert\theta\rvert\leq1\},
   \qquad E_d=J\cap\{\lvert\theta\rvert>1\},
   \qquad E_j=\varnothing\quad(2\leq j\leq d-1),
   \qquad \lvert E_1\rvert+\lvert E_d\rvert=\lvert J\rvert,
\end{equation}
   
\begin{equation}
\label{eq:appendix-indented-049}
   T_1=-\theta^d-\sum_{k=1}^{d-1}\beta_k\theta^k,
   \qquad \Psi_1=(T_1,\beta_1,\ldots,\beta_{d-1}),
\end{equation}
   
\begin{equation}
\label{eq:appendix-indented-050}
   T_d=-\theta-\sum_{k=0}^{d-2}\beta_k\theta^{k-d+1},
   \qquad \Psi_d=(\beta_0,\ldots,\beta_{d-2},T_d),
\end{equation}
   
\begin{equation}
\label{eq:appendix-indented-051}
   \lambda_{d-1}([-R,R]^{d-1})=(2R)^{d-1},
   \qquad p_{\Psi_1(\theta,\beta)}(\theta)
   =p_{\Psi_d(\theta,\beta)}(\theta)=0,
\end{equation}
   
\begin{equation}
\label{eq:appendix-indented-052}
   \partial_\theta T_1
   =-d\theta^{d-1}-\sum_{k=1}^{d-1}k\beta_k\theta^{k-1},
   \qquad
   \lvert\partial_\theta T_1\rvert\leq d+\frac{Rd(d-1)}2,
\end{equation}
   
\begin{equation}
\label{eq:appendix-indented-053}
   \partial_\theta T_d
   =-1+\sum_{k=0}^{d-2}(d-1-k)\beta_k\theta^{k-d},
\end{equation}
   
\begin{equation}
\label{eq:appendix-indented-054}
   \lvert\partial_\theta T_d\rvert
   \leq1+R\sum_{m=1}^{d-1}\frac{m}{\lvert\theta\rvert^{m+1}}
   \leq1+\frac{Rd(d-1)}2
   \leq d+\frac{Rd(d-1)}2.
\end{equation}
   
   At \(d=1\), \(B=\left(\begin{smallmatrix}0&1\\0&0\end{smallmatrix}\right)\),
   \(\widehat\Lambda_{B,T}=1\), \(E_1=J\), the sole pivot is \(\alpha_0\), the beta mass is one,
   \(T_1=-\theta\), and \(\lvert\partial_\theta T_1\rvert=1\).
   Every boundary, sign, empty-cell, dimension, and interval-location branch stated by the wrapper is retained.
   Finally, for every such bounded \(J\), including empty and singleton intervals,
   
\begin{equation}
\label{eq:appendix-indented-055}
   \Pr_{\alpha\sim\mu}[\exists\theta\in J:p_\alpha(\theta)=0]
   \leq
   \kappa(2R)^{d-1}
   \left(d+\frac{Rd(d-1)}2\right)\lvert J\rvert.
\end{equation}

5. \textbf{Counter-example 1 scale calculation.} For \(0<\delta\leq1\),
   \(\Theta=[-1,1]\), \(\widetilde F=(0,1,\theta/\delta)\), \(R=1\), \(\kappa=1/4\), and \(A=1\),
   
\begin{equation}
\label{eq:appendix-indented-056}
   q=M=0,
   \qquad \Delta=1,
   \qquad N=2,
   \qquad m=0,
\end{equation}
   
   and the sole nonzero closure entry is \(B_{2,1}=1/\delta\), with
   
\begin{equation}
\label{eq:appendix-indented-057}
   \widehat\Lambda_{B,T}
   =\Gamma_{\mathrm{proj}}(F)
   =\frac1\delta.
\end{equation}
   
   For every possibly correlated law in \(\mathcal D_{2,1,1/4}\) and every positive-length interval
   \(I\subseteq[-1,1]\),
   
\begin{equation}
\label{eq:appendix-indented-058}
   \Pr_{\alpha\sim\mu}[\exists\theta\in I:
   \alpha_1+\alpha_2\theta/\delta=0]
   \leq\frac1\delta\lvert I\rvert,
   \qquad
   C^{\mathrm{Pf}}_{\mathcal D}(F;[-1,1])\leq\frac1\delta.
\end{equation}
   
   For the selected independent uniform law on \([-1,1]^2\) and every
   \(0<\epsilon\leq\delta\), the complete closed two-wedge event includes both sign branches, coefficient
   axes, square boundaries, and the roots at \(\theta=0\) and \(\theta=\epsilon\), and
   
\begin{equation}
\label{eq:appendix-indented-059}
   \Pr[\exists\theta\in[0,\epsilon]:
   \alpha_1+\alpha_2\theta/\delta=0]
   =\frac{\epsilon}{4\delta}.
\end{equation}
   
   Consequently the selected-law positive-length lower ratio is exactly \(1/(4\delta)\), separately from the
   all-law upper coefficient \(1/\delta\), the raw certificate \(1/\delta\), and the geometric normalization
\begin{equation}
\label{eq:appendix-indented-060}
   \sqrt2\,\widehat\Lambda_{B,T}=\frac{\sqrt2}{\delta}.
\end{equation}
   The last quantity is not a probability coefficient. No optimality claim is made.

6. \textbf{Dependence, mode, baseline, and nonoutput closure.} The deterministic instance
   \((\Theta,T,q,M,\Delta,N,R,\kappa,A,m,B)\) is fixed before a law and interval are selected. The law may have
   arbitrary coefficient correlation; probability is ordinary probability for each fixed law; every general
   interval statement ranges over all positive-length subintervals with its literal endpoint convention; and
   each capacity takes the interval supremum for a fixed law before the outer law supremum. All vector,
   operator, and Frobenius norms are Euclidean, and section measure is Euclidean Hausdorff measure, with the
   supplied (N=1) convention. Every stated constant is literal, there is no hidden constant or confidence
   parameter, and once \(\widehat\Lambda_{B,T}\) is fixed its additional dependence on \(q,M,\Delta\) is exactly
   degree zero.

   The exact monic deterministic-leading-coefficient theorem and the distinct Counter-example scales
   \(1/(4\delta)\), \(1/\delta\), and \(\sqrt2/\delta\) are preserved without a remainder, changed mode,
   or surrogate object. No result here asserts that every unrestricted raw Pfaffian presentation admits the
   supplied derivative-closure certificate or a polynomial presentation-format bound. No independence is
   claimed outside the selected uniform lower-law computation; no random leading coordinate, chart-count
   factor, interval-location factor, auxiliary tolerance, new source result, term absorption, new probability
   conversion, or proof-local general exhaustion object is introduced.


\end{proposition}

\begin{proof}
We verify each clause of the conjunction.

For conclusion 1, Lemma~\ref{lem:step-001-anchor}, Lemma~\ref{lem:step-001-height},
Lemma~\ref{lem:step-001-homogeneous-block}, Proposition~\ref{prop:step-001-projective}, and
Proposition~\ref{prop:step-001-boundary} give exactly the static, homogeneous, projective, endpoint,
stationary, zero-height, and certificate-baseline statements. Nothing is inferred beyond those conclusions.

For the chart portion of conclusion 2, Lemma~\ref{lem:step-003-measurable-domains} supplies measurable
completed domains and exact exhaustion, Lemma~\ref{lem:step-003-finite-chart} supplies finite localized charts
and the determinant, Proposition~\ref{prop:step-003-finite-area} supplies exact area/multiplicity and joint
density control, Lemma~\ref{lem:step-003-root-coverage} supplies persistent-root removal and full finite-chart
coverage, and Proposition~\ref{prop:step-003-pivot-sweep} supplies both public chart inequalities.
The finite domains, images, and exhaustion levels establish the chart bounds
but do not appear in the present conclusion. For the coordinate-free rate, capacity, and zero-height portion,
Proposition~\ref{prop:step-007-s2-general-affine-rate} supplies every conclusion with the same
objects and modes.

For conclusion 3, Lemma~\ref{lem:step-008-s2-radial-cancellation} supplies the exact actual-section identity;
Proposition~\ref{prop:step-008-s2-stationary-projective} supplies the fixed law-null stationary branch;
Lemma~\ref{lem:step-008-s2-literal-algebra} supplies the literal coefficient equality;
Proposition~\ref{prop:step-008-s2-pairwise-homogeneous-rate} supplies the two pairwise rates; and
Proposition~\ref{prop:step-008-s2-pf-closure} supplies the capacity bound and its exact supremum order.

For conclusion 4, Proposition~\ref{prop:step-010-complete-affine-monic-wrapper} gives the result directly.
Its six-clause statement already contains the tuple, coefficient dimension, deterministic leader, descriptor
values, full shift matrix and (d=1) branch, certificate, cells, pivots, charts, insertion maps, beta volume,
velocity bounds, all boundary/location branches, zero-length nullity, and exact probability inequality. No
other monic result is needed.

For conclusion 5, Lemma~\ref{lem:step-011-shear-certificate} supplies the instance and exact certificate,
Lemma~\ref{lem:step-011-projective-speed} supplies the exact Euclidean projective speed,
Proposition~\ref{prop:step-011-upper-coefficient} supplies the all-law upper statements,
Lemma~\ref{lem:step-011-root-wedges} supplies the complete closed event geometry,
Proposition~\ref{prop:step-011-wedge-probability} supplies the selected-law exact probability, and
Proposition~\ref{prop:step-011-scale-comparison} supplies the distinct positive-length lower ratio and scale
comparison.

Conclusion 6 repeats the exact dependence, probability, interval, norm, boundary, baseline, and nonoutput
statements already contained in those named results. Conjunction changes no quantifier, object, dimension,
constant, law class, interval class, probability mode, norm, or measure. Each result is applied to the
identical setting object. There is no term absorption, new cited result, new source invocation,
new probability conversion, hidden constant, or additional
mathematical claim. This proves the proposition. 
\end{proof}

\begin{proof}[Synthesis of the theorem clauses]
The static and projective clauses are exactly the conclusions of
Lemmas~\ref{lem:step-001-anchor}, \ref{lem:step-001-height}, and
\ref{lem:step-001-homogeneous-block} and
Propositions~\ref{prop:step-001-projective} and
\ref{prop:step-001-boundary}.  The equivalent chart clause is exactly
Proposition~\ref{prop:step-003-pivot-sweep}, with its measurability,
multiplicity, and persistent-root conclusions supplied by the other named
results in that subsection.  Proposition~\ref{prop:step-007-s2-general-affine-rate}
supplies the coordinate-free affine rate, capacity, and zero-certificate
case.  The five named homogeneous results in the projective-rate subsection
supply the sharper rate and preserve the interval-then-law supremum order.
Proposition~\ref{prop:step-010-complete-affine-monic-wrapper} gives the complete
deterministic-leading-coefficient monic clause.
Finally, the six named Counter-example results supply the exact certificate,
projective speed, all-law upper coefficient, closed wedge geometry, selected-law
probability, and normalized lower ratio.  The identity
$\sqrt2\widehat\Lambda_{B,T}=\sqrt2/\delta$ is direct scalar substitution
and is kept separate from both probability quantities.

All six groups use their original deterministic presentation, coefficient
dimension, law class, interval class, probability mode, and Euclidean norm
and measure conventions. Their conjunction changes no assumption, quantifier,
object, constant, probability conversion,
or supremum order.  This gives Proposition~\ref{prop:step-012-anchored-coefficient-sweep}
with the exact deterministic-presentation-first order and both baseline
invariance conclusions.
\end{proof}

\subsection{Proof of the Main Theorem}\label{app:main-proof}

\begin{proof}[Proof of Theorem~\ref{thm:main}]
Under Assumptions~\ref{assump:parameter-regime}--\ref{assump:cube-density-laws},
Proposition~\ref{prop:step-012-anchored-coefficient-sweep} applies to the
identical fixed presentation.  Its first conclusion is the theorem's static
derivative-closure and normalized projective certificate.  Its second
conclusion gives both original-coordinate chart forms, the coordinate-free
affine rate and capacity, and the zero-certificate branch.  Its third
conclusion is the sharper homogeneous rate and capacity, including the
stationary branch and the interval-before-law supremum order.  Its fourth
conclusion is the exact affine-monic presentation, certificate, chart, and
ordinary-probability recovery with the deterministic leading coefficient
outside the random vector.  Its fifth conclusion is Counter-example 1 with
the distinct scales $1/(4\delta)$, $1/\delta$, and $\sqrt2/\delta$, the last
being geometric rather than probabilistic.  Its sixth conclusion supplies
all dependence, mode, norm, measure, baseline-invariance, and nonoutput
declarations, including the boundary excluding any claim for unrestricted
raw Pfaffian presentations.  These are exactly the clauses of
Theorem~\ref{thm:main}, so the theorem follows.
\end{proof}
